%=====================================================================
%  E-COMMERCE APPLICATION DEVELOPMENT --- FALL 2026
%  Unified Master Lecture Handout
%
%  Taught as INFT-4041, BSIT 5th semester, 3 credit hours.
%
%  Written against the HEC 2017 outline (BS & MS (CS, SE & IT)
%  Curriculum 2017, p. 94), which is the ONLY published specification of
%  this course in any HEC edition.  There it is called simply
%  "E-Commerce" and is an ELECTIVE in BS Software Engineering (2017
%  p. 38, "select any FIVE"), not a BSIT course at all.
%
%  The 2023 edition drops the course entirely -- it appears nowhere in
%  BS Curriculum Computing Disciplines 2023 except as an incidental
%  application domain on p. 23.  The 2025 edition lists "Digital
%  Marketing & E-Commerce" as an Interdisciplinary/Allied course (p. 16,
%  item 11) with credit hours but no contents.
%
%  So there is one outline to trace against, it is nine years old, and
%  its "Course Content" paragraph is the table of contents of its own
%  reference book number 2 (Peacock, PHP 5 E-commerce Development,
%  Packt, 2010), duplicated clauses and all.  The Course Specification
%  in the front matter records all of that, traces every distinct clause
%  to a chapter, and justifies the three topics taught that the outline
%  does not name.
%
%  Build:  latexmk -pdf ECAD_Fall2026_Handout.tex
%     or:  .\build.ps1
%
%  Structure: 6 parts, 23 chapters, 7 appendices.
%  All style and pedagogy macros live in ecadhandout.sty.
%=====================================================================
\documentclass[11pt,a4paper,oneside]{report}
\usepackage{ecadhandout}

\begin{document}

%---------------------------------------------------------------------
%=====================================================================
% FRONT MATTER
%=====================================================================
\pagenumbering{roman}
\structuralfigures{F.}
\thispagestyle{empty}

\begin{center}
\vspace*{1.0cm}

{\color{secondary}\rule{0.85\textwidth}{2pt}}\\[6pt]
{\large\textsc{\color{secondary}Bachelor of Science in Information Technology}}\\[18pt]

{\Huge\bfseries\color{primary}E-COMMERCE APPLICATION\\[4pt]DEVELOPMENT}\\[10pt]
{\LARGE\bfseries\color{primary}Building, Running and Securing an Online Shop}\\[16pt]

{\color{secondary}\rule{0.85\textwidth}{2pt}}\\[16pt]

{\Large\textsc{Unified Lecture Handout}}\\[6pt]
{\large\color{gray!70!black}Complete Course in One Volume $\bullet$ Fall 2026}\\[24pt]

% --- the claim the whole book rests on ---
\begin{tikzpicture}
  \node[rounded corners=4pt, draw=primary, line width=1.1pt, fill=primary!6,
        text=primary, align=center, font=\small\bfseries,
        minimum width=11.6cm, minimum height=1.5cm]
       {\faIcon{store}\\[2pt]Anyone can install a storefront in an afternoon.
        This course is everything after that.};
\end{tikzpicture}\\[10pt]
{\small\itshape\color{gray!70!black}A shop is not a website. It takes money,
holds other people's data, and is answerable in law --- and each of those
changes what ``finished'' means.}

\vspace{1.0cm}

\begin{tcolorbox}[enhanced, width=0.86\textwidth, colback=lightbg,
  colframe=primary, boxrule=0pt, leftrule=4pt, arc=2pt, top=6pt, bottom=6pt]
\small
\textbf{\color{primary}What this volume covers.} Every distinct clause of the
HEC outline for \emph{E-Commerce}, in the order the outline gives them, traced
chapter by chapter in the Course Specification overleaf.

\smallskip
That outline is published once only, in the \textbf{2017} curriculum (p.~94),
where the course is an \emph{elective in BS Software Engineering} rather than a
BSIT course at all. The 2023 curriculum drops it completely; the 2025
curriculum renames it \emph{Digital Marketing \& E-Commerce} and publishes no
contents. The department teaches it as \textbf{INFT-4041}, BSIT, 5th semester,
3 credit hours. All of that is recorded in the Course Specification rather
than smoothed over.

\smallskip
The course is taught on a \textbf{running shop}. One storefront is installed in
Chapter~5 and every chapter afterwards adds a layer to it; by Chapter~23 it is
the semester project. Appendix~A sets up the stack in one sitting.
\end{tcolorbox}

\vspace{0.9cm}
{\small\color{gray!60!black}Six Parts $\bullet$ Twenty-Three Chapters $\bullet$ Seven Appendices}
\vspace{1cm}
\end{center}

\newpage

%=====================================================================
\chapter*{How to Use This Handout}
\addcontentsline{toc}{chapter}{How to Use This Handout}
\markboth{How to Use This Handout}{}

There are two ways to teach this subject badly.

The first is to teach it as business studies: models, channels, trends, the
history of the dot-com boom, and a graduate who can describe an online shop but
has never configured a tax rule. The second is to teach it as programming
exercises: write a cart class, write a checkout function, and a graduate who has
built a toy that no one would take money through.

This handout takes the third path, which is also the one the published outline
describes if you read it closely. You install a real, open-source storefront in
Chapter~5, and you spend the rest of the course making it into a shop that could
actually trade: a catalogue that models real products, a checkout that handles
shipping and tax, a payment path that survives a failed transaction, an admin
who can process a refund, a server that does not leak its database, and a listing
a search engine can find. Where the platform cannot do what you need --- and in
Chapter~19 it cannot --- you write the code that extends it.

\section*{The six parts}

\textbf{Part~I --- E-Commerce and Its Models} (Chapters 1--4) is the business
grounding, and it is short on purpose. What electronic commerce actually changed;
the models and where the revenue comes from; auctions, portals and communities;
and the build-buy-or-extend decision that justifies everything the rest of the
book does.

\textbf{Part~II --- Standing the Shop Up} (Chapters 5--8) gets a storefront
running and fills it. The stack and what each layer does; products and categories,
with the tables underneath them; variants, options and uploads; then the theme.

\textbf{Part~III --- From Basket to Order} (Chapters 9--12) is the pipeline a
customer walks through. Where basket state lives; the checkout and the order
state machine; shipping, tax and zones; discounts, vouchers and referrals.

\textbf{Part~IV --- Money and Identity} (Chapters 13--16) is the part where
mistakes cost money. Payment systems and who carries the risk; integrating a
gateway and testing it in a sandbox; accounts, sessions and password storage;
then securing the whole thing.

\textbf{Part~V --- Running and Extending the Shop} (Chapters 17--20) is the half
of the job nobody demonstrates: the dashboard and what to watch, the back office,
writing your first module, and deployment with everything that follows it.

\textbf{Part~VI --- Reach, Law and Practice} (Chapters 21--23) closes with being
found, being lawful, and your project.

\begin{conceptbox}
Why Chapter~19 is the hinge of the course. Everything before it is
\emph{configuration}: you are choosing among the options a platform already
offers, and a careful non-programmer could do most of it. Chapter~19 is the
first time the answer to ``can the shop do this?'' is no, and you have to write
the code that makes it yes. That is the difference between the outline's bare
title, \emph{E-Commerce}, and the one this department teaches,
\emph{E-Commerce Application Development}.
\end{conceptbox}

\section*{The furniture of every chapter}

\rowcolors{2}{white}{lightbg}
\begin{longtable}{@{}L{4.6cm}L{10.6cm}@{}}
\toprule
\rowcolor{primary!15}
\textbf{Box} & \textbf{What it is for} \\
\midrule
\endhead
\faIcon{bullseye}~\textbf{Outcomes} & What you will be able to \emph{do}. Bloom
verbs, not topics. \\
\faIcon{link}~\textbf{Before you start} & Which earlier chapters this one stands
on, and what state your shop should be in. \\
\faIcon{tags}~\textbf{Key terms} & The vocabulary introduced; all of it is in the
glossary of Appendix~D. \\
\faIcon{book}~\textbf{Definition} & The canonical statement. In this subject
several ordinary English words --- order, basket, customer, product --- are
technical terms with narrower meanings, and confusing the two causes real bugs. \\
\faIcon{lightbulb}~\textbf{Concept} & The judgement call: what an experienced
practitioner would decide, and why. \\
\faIcon{exclamation-triangle}~\textbf{Warning} & Something that will bite you,
usually in production rather than in the examination. \\
\faIcon{calculator}~\textbf{Worked example} & Real arithmetic --- a tax
calculation, a gateway fee, a discount applied in the wrong order --- with every
step shown. \\
\faIcon{tachometer-alt}~\textbf{Measured} & A number produced by running
something, not by asserting it. Twelve chapters carry one; see below. \\
\faIcon{exclamation-circle}~\textbf{Common pitfalls} & The mistakes that are
actually made. \\
\faIcon{flask}~\textbf{Try it yourself} & The laboratory. Every one of them
changes the same shop, and they are cumulative: skip Chapter~9's and
Chapter~10's will have nothing to check out. \\
\faIcon{question-circle}~\textbf{Checkpoint} & Three questions. Attempt them
before reading Appendix~C. \\
\faIcon{clipboard-check}~\textbf{Chapter summary} & What to retain. \\
\faIcon{pen-fancy}~\textbf{Review questions} & Examination practice: multiple
choice, short answer, applied. \\
\bottomrule
\end{longtable}

\newpage

Two typographic conventions carry a lot of weight in a course taught through a
graphical interface, and they are worth learning on this page:

\begin{itemize}[leftmargin=1.6em, itemsep=3pt, topsep=4pt]
  \item \adminpath{Catalog, Products, Add New} is a \textbf{click path} through
        the administration panel. Read the chevrons as ``then''.
  \item \uiel{Save} on its own is a single control --- a button, a field, a tab
        or a menu item. Anything in \texttt{typewriter} is text you type or a
        file you edit, which is a different thing entirely.
\end{itemize}

\begin{alertbox}[The measured strand, and why it is only in twelve chapters]
Twelve of the twenty-three chapters end with a number that was produced rather
than asserted: a page time, a query count, a request rate, a hash cost, a list
of security headers present and absent.

\smallskip
They are the twelve where measuring genuinely teaches something. ``Add an index
and the category page gets faster'' is a claim a student can be told and will
forget; watching the same page go from a table scan at fifty thousand products
to an index seek, and reading both numbers off your own machine, is not. The
other eleven chapters argue in prose, because dressing up a legal obligation or
a business model with a stopwatch would be theatre.

\smallskip
Every figure in a \faIcon{tachometer-alt}~\textbf{Measured} box was produced by
the file named beside it, in \texttt{code/}, against the stack in Appendix~A.
Rerun them. Your numbers will differ from the book's --- different machine,
different disk, different day --- and the \emph{shape} should not. Where a
result depends on the particular machine, the box says so.
\end{alertbox}

\newpage

%=====================================================================
\chapter*{Course Specification}
\addcontentsline{toc}{chapter}{Course Specification}
\markboth{Course Specification}{}

This page is the course file. It states what HEC requires of a course called
\emph{E-Commerce} and shows where each requirement is met.

That is harder here than in most courses, because the requirement has been
published exactly once and then withdrawn. The three editions are traced below
so that an accreditation review against any of them has something to check.

\section*{Course particulars}

\rowcolors{2}{white}{lightbg}
\begin{longtable}{@{}L{4.2cm}L{11.0cm}@{}}
\toprule
\rowcolor{primary!15}
\textbf{Item} & \textbf{Value} \\
\midrule
\endhead
Course title & \emph{E-Commerce Application Development}, as taught. HEC 2017
calls it simply \emph{E-Commerce}; HEC 2025 calls it \emph{Digital Marketing \&
E-Commerce}. The department's title is the accurate one, and the Course
Specification justifies the two added words \\
Course code & INFT-4041 \\
Programme & BS Information Technology, 5th semester. Under HEC 2017 the course
is an \textbf{elective in BS Software Engineering} (2017, p.~38, ``select any
FIVE'') and does not appear in the BSIT scheme at all \\
Credit hours & 3: three credit hours of lectures, no separate laboratory credit.
The laboratory work in this handout is therefore done in lecture time and as
homework, which is why every lab is small and cumulative \\
Prerequisite & \emph{Web Technologies} (CS-207 / INFT-4005, 4th semester). HEC
2017 names the prerequisite \emph{Web Engineering}; the department's course of
that content is called Web Technologies. The name differs, the dependency is
satisfied --- see Appendix~F \\
Duration & Sixteen teaching weeks, midterm in Week~8, project due Week~16
(Appendix~G) \\
Teaching methodology & As published (HEC 2017): lecturing, written assignments,
\textbf{project}, report writing \\
Assessment & As published: sessional examination, home assignments, quizzes,
project, presentations, final examination \\
Set texts & As published: Laudon and Traver, \emph{E-Commerce}; Peacock,
\emph{PHP 5 E-commerce Development}; Rayport, \emph{Introduction to E-Commerce};
Schneider, \emph{Electronic Commerce}. See defect~5 below, and Appendix~E \\
\bottomrule
\end{longtable}

\section*{The three editions}

\rowcolors{2}{white}{lightbg}
\begin{longtable}{@{}L{1.5cm}L{3.6cm}L{6.4cm}L{3.3cm}@{}}
\toprule
\rowcolor{primary!15}
\textbf{Ed.} & \textbf{Title} & \textbf{Status} & \textbf{Contents} \\
\midrule
\endhead
2017 & E-Commerce & Elective in BS Software Engineering (p.~38); absent from
BSIT. 3(3,0); prerequisite \emph{Web Engineering} & \textbf{Yes, p.~94 --- the
only published outline in any edition} \\
2023 & --- & \textbf{Absent entirely.} The course appears nowhere in the
computing-disciplines curriculum except as an incidental application domain
(p.~23) & No \\
2025 & Digital Marketing \& E-Commerce & Interdisciplinary / Allied course,
item~11 (p.~16). Not a computing core; universities may add others & No --- title
and credit hours only \\
\bottomrule
\end{longtable}

\begin{alertbox}[Five defects in the published document]
Recorded here so that a reader comparing this handout with the HEC booklet is
not confused, and so that the department's own course file does not inherit
them.

\smallskip
\textbf{1. Two learning outcomes, both too broad to assess.} HEC publishes
exactly two: ``Understand the concepts and standards related to the discipline
of E-Commerce'' and ``Analyze complex real world problems found in E-Commerce''.
Sibling courses publish eight. Two cannot carry sixteen weeks or a question
paper, so this department adds six, each traced back to one of HEC's two and
each marked as departmental in the table below.

\smallskip
\textbf{2. The Bloom's Taxonomy column is blank} for both outcomes --- only the
Domain column (``C'') is filled. The levels below are therefore
\emph{assigned by this department}, and are marked as such.

\smallskip
\textbf{3. The ``Course Content'' paragraph is the table of contents of its own
reference book.} Clause for clause, in order, it is the chapter list of
Michael Peacock's \emph{PHP 5 E-commerce Development} (Packt, 2010), which is
reference~2 on the same page. HEC did not write an outline for this course; it
pasted a contents page.

\smallskip
\textbf{4. Two clauses are duplicated}, which is what pasting a contents page
does. ``Checkout'' appears twice --- once as ``The Checkout and Order Process''
--- and ``Managing Products and Categories'' appears twice, once on its own and
once under Administration. Twenty-four clauses; twenty-two distinct.

\smallskip
\textbf{5. The prescribed technology is end-of-life and the editions are stale.}
PHP~5 lost security support on 31~December 2018: the prescribed text targets a
runtime that has been unsafe to deploy for longer than most of this class has
been studying. Laudon and Traver is cited at the 13th edition (2017) and Rayport
at 2007. Appendix~E cites current editions and explains each gap.
\end{alertbox}

\section*{Coverage of the HEC 2017 course outline}

The published outline is a single run-on paragraph. Each of its clauses is
traced below, in the order HEC writes them. Nothing is omitted, and the two
duplicates are shown where they fall rather than quietly merged.

\rowcolors{2}{white}{lightbg}
\begin{longtable}{@{}L{6.4cm}L{2.0cm}L{6.6cm}@{}}
\toprule
\rowcolor{primary!15}
\textbf{HEC 2017 clause} & \textbf{Chapter} & \textbf{Where exactly} \\
\midrule
\endhead
``An overview of E-Commerce \& its business models and concepts'' & 1, 2 & What
changed and why; then the models, with the revenue model treated as the
question that decides the rest \\
``Planning an E-Commerce Framework'' & 4 & Requirements, and the
build-buy-or-extend decision this course's answer rests on \\
``Managing Products and Categories'' & 6 & The admin screens, and the tables
beneath them \\
``Product Variations and User Uploads'' & 7 & Options and variants, the
combinatorics they generate, and why an upload field is a security question \\
``Enhancing the User Experience'' & 8 & Themes and template overrides, measured
by page weight rather than taste \\
``The Shopping Basket'' & 9 & Session, cookie or database, and what survives a
restart \\
``The Checkout and Order Process'' & 10 & The order state machine, step by
step \\
``Shipping and Tax'' & 11 & Geo zones, tax classes, and sales tax on a
Pakistani online sale \\
``Discounts, Vouchers, and Referrals'' & 12 & All three, and the order in which
they must be applied \\
``Checkout'' \emph{(duplicate of clause~7)} & 10 & Same chapter; see defect~4 \\
``Taking Payment for Orders'' & 14 & The gateway handshake, the sandbox, and
what happens when it fails halfway \\
``User Account Management'' & 15 & Registration, sessions, password storage,
reset flows \\
``Administration: Dashboard'' & 17 & The dashboard, and which of its numbers
are worth acting on \\
``Managing Products and Categories'' \emph{(duplicate of clause~3)} & 6, 17 &
Same material from the back office; see defect~4 \\
``Managing Orders'' & 18 & Order states, fulfilment, partial shipment \\
``Customers'' & 18 & Customer records, groups and approval \\
``Refunds'' & 18 & Returns, credit and what the order state machine does next \\
``Voucher Codes'' & 12, 18 & Created in 12, administered in 18 \\
``Shipping'' \emph{(administration)} & 11, 18 & Configured in 11, operated
in 18 \\
``Deploying, Security, and Maintenance'' & 16, 20 & Security is given its own
chapter --- see the note below --- and deployment, backup and maintenance
share 20 \\
``Web Payment Systems'' & 13 & Cards, wallets, bank transfer and cash on
delivery; cost and risk for each \\
``Social, Legal, and Ethical Issues of E-Commerce'' & 22 & Consumer rights,
privacy, dark patterns, accessibility, and the Pakistani statute \\
``Auctions, Portals, and Communities'' & 3 & Moved forward to Part~I, where it
belongs: these are business models \\
``SEO'' & 21 & Crawlability, URL structure, metadata, sitemaps \\
\bottomrule
\end{longtable}

\section*{Course learning outcomes}

The two outcomes HEC publishes, verbatim, followed by the six this department
adds. Bloom levels are assigned by this department throughout, because the
published table leaves the column blank.

\rowcolors{2}{white}{lightbg}
\begin{longtable}{@{}L{1.3cm}L{7.6cm}L{2.1cm}L{3.8cm}@{}}
\toprule
\rowcolor{primary!15}
\textbf{CLO} & \textbf{By the end of the course the student will be able to} &
\textbf{Bloom} & \textbf{Chapters} \\
\midrule
\endhead
CLO-1 & \emph{(HEC, verbatim)} Understand the concepts and standards related to
the discipline of E-Commerce & C2 Understand & 1--4, 13, 22 \\
CLO-2 & \emph{(HEC, verbatim)} Analyze complex real world problems found in
E-Commerce & C4 Analyse & 6, 10, 16, 20, 23 \\
CLO-3 & \emph{(dept.)} Explain e-commerce business and revenue models, and
classify a given venture against them & C2 Understand & 2, 3 \\
CLO-4 & \emph{(dept.)} Install, configure and secure an open-source storefront
on a reproducible stack & C3 Apply & 5, 16, 20 \\
CLO-5 & \emph{(dept.)} Model a catalogue --- products, categories, options and
variants --- and justify the schema that results & C3 Apply & 6, 7 \\
CLO-6 & \emph{(dept.)} Configure the order pipeline end to end: basket,
checkout, shipping, tax, discounts and payment & C3 Apply & 9--12, 14 \\
CLO-7 & \emph{(dept.)} Analyse a storefront's performance and security posture
from measurement rather than assertion & C4 Analyse & 6, 8, 15, 16, 20, 21 \\
CLO-8 & \emph{(dept.)} Extend the platform with a module meeting a stated
requirement, then deploy, present and defend the result & C5 Evaluate &
19, 20, 23 \\
\bottomrule
\end{longtable}

\begin{conceptbox}[Three topics taught that the outline does not name]
Each is justified here rather than smuggled in, because a reviewer will ask.

\smallskip
\textbf{Chapter~6's data model.} The outline names the screens --- ``Managing
Products and Categories'' --- but never the tables. A graduate who can fill in
an admin form but cannot say why a product's name lives in a different table
from its price has not understood the catalogue, and \emph{Database Systems} is
a prerequisite that ought to be paid off rather than left on the shelf.

\smallskip
\textbf{Chapter~16, Security, as a chapter of its own.} HEC buries it inside
``Deploying, Security, and Maintenance'', one clause among three. A shop holds
addresses, order histories and password hashes, and CLO-7 asks for a security
posture to be \emph{analysed}. One third of one clause will not do it.

\smallskip
\textbf{Chapter~19, module development.} The outline is configuration from end
to end; nothing in it requires a student to write a line of code. The
department's title says \emph{Application Development}, the HEC teaching
methodology requires a project, and a project that is entirely configuration is
not a computing project.
\end{conceptbox}

\begin{alertbox}[Where this course sits, and the one awkwardness]
The outline this handout traces was written for \emph{software engineering}
students as an elective, and it shows: it assumes the reader will build a cart
from scratch in PHP. This department teaches it to \emph{information technology}
students as a 5th-semester course, where the useful skill is standing up,
configuring, securing and extending a real platform rather than reimplementing
one badly.

\smallskip
This handout therefore follows the outline's \emph{clauses} faithfully and its
\emph{method} not at all. Every topic HEC lists is taught; none of it is taught
by writing a shopping cart in PHP~5. Chapter~4 makes that argument explicitly,
so that the student can defend the choice rather than merely inherit it.
\end{alertbox}

\newpage

%=====================================================================
\chapter*{The Course at a Glance}
\addcontentsline{toc}{chapter}{The Course at a Glance}
\markboth{The Course at a Glance}{}

\dsfigh{%
\begin{tikzpicture}[
  part/.style={rounded corners=3pt, fill=#1, text=white, font=\bfseries\footnotesize,
               minimum width=2.5cm, text width=2.2cm, minimum height=0.95cm,
               align=center},
  ch/.style={rounded corners=2pt, draw=#1, line width=0.7pt, fill=#1!7, text=#1!60!black,
             font=\scriptsize, minimum width=2.5cm, minimum height=0.52cm, align=center},
  arr/.style={-{Stealth[length=2mm]}, primary!45, line width=0.8pt}
]
\node[part=primary]   (p1) at (0,0)    {Part I\\Models};
\node[part=secondary] (p2) at (3.0,0)  {Part II\\Standing It Up};
\node[part=greenhl]   (p3) at (6.0,0)  {Part III\\Basket to Order};
\node[part=purplehl]  (p4) at (9.0,0)  {Part IV\\Money and Identity};
\node[part=tealhl]    (p5) at (12.0,0) {Part V\\Running It};
\node[part=goldhl]    (p6) at (15.0,0) {Part VI\\Reach and Law};

\foreach \a/\b in {p1/p2,p2/p3,p3/p4,p4/p5,p5/p6}{\draw[arr] (\a) -- (\b);}

\foreach \i/\t in {1/{1. What changed}, 2/{2. Business models},
                   3/{3. Auctions, portals}, 4/{4. Build or extend}}{
  \pgfmathsetmacro{\yy}{-0.90-(\i-1)*0.62}
  \node[ch=primary] at (0,\yy) {\t};}

\foreach \i/\t in {1/{5. The stack}, 2/{6. Catalogue},
                   3/{7. Variations}, 4/{8. Theme and UX}}{
  \pgfmathsetmacro{\yy}{-0.90-(\i-1)*0.62}
  \node[ch=secondary] at (3.0,\yy) {\t};}

\foreach \i/\t in {1/{9. The basket}, 2/{10. Checkout},
                   3/{11. Shipping, tax}, 4/{12. Discounts}}{
  \pgfmathsetmacro{\yy}{-0.90-(\i-1)*0.62}
  \node[ch=greenhl] at (6.0,\yy) {\t};}

\foreach \i/\t in {1/{13. Payment systems}, 2/{14. The gateway},
                   3/{15. Accounts}, 4/{16. Security}}{
  \pgfmathsetmacro{\yy}{-0.90-(\i-1)*0.62}
  \node[ch=purplehl] at (9.0,\yy) {\t};}

\foreach \i/\t in {1/{17. Dashboard}, 2/{18. Orders, refunds},
                   3/{19. Your first module}, 4/{20. Deployment}}{
  \pgfmathsetmacro{\yy}{-0.90-(\i-1)*0.62}
  \node[ch=tealhl] at (12.0,\yy) {\t};}

\foreach \i/\t in {1/{21. SEO}, 2/{22. Law and ethics},
                   3/{23. Project}}{
  \pgfmathsetmacro{\yy}{-0.90-(\i-1)*0.62}
  \node[ch=goldhl] at (15.0,\yy) {\t};}

\node[rounded corners=3pt, fill=primary!7, draw=primary!30,
      minimum width=17.3cm, minimum height=0.66cm,
      font=\scriptsize\bfseries, text=primary] at (7.5,-4.35)
      {\faIcon{store}~~One shop, installed in Chapter 5 and added to every week
       until it is the project in Chapter 23.};
\end{tikzpicture}}%
{The six parts and the chapters in each. Chapter 19 is the hinge: it is the
first chapter in which configuring the platform is not enough and you have to
write code.}{coursemap}

\vspace{4pt}
{\small\itshape\color{gray!70!black}
This course stands on \emph{Web Technologies} and \emph{Database Systems}, and
shares its deployment ground with \emph{Virtual Systems and Services} in the
same semester. Appendix~F states what belongs to each.}

\newpage

%=====================================================================
\chapter*{The Shop We Build}
\addcontentsline{toc}{chapter}{The Shop We Build}
\markboth{The Shop We Build}{}

Every lab in this book changes the same shop. Nothing is thrown away, and
nothing is built twice. This page is the map of what gets added when, so that
you can tell at any point in the semester whether you are behind.

\dsfigh{%
\begin{tikzpicture}[
  lay/.style={rounded corners=2pt, draw=#1, line width=0.8pt, fill=#1!8,
              text=#1!55!black, font=\scriptsize, minimum height=0.62cm,
              text width=4.6cm, align=center},
  sidenote/.style={font=\scriptsize\bfseries, text=primary, anchor=west}
]
\node[rounded corners=3pt, fill=primary, text=white, font=\bfseries\footnotesize,
      minimum width=5.0cm, minimum height=0.8cm, align=center] (base) at (0,0)
     {Ch 5\\The stack, running};

\foreach \i/\c/\t in {
   1/secondary/{Ch 6--7: a catalogue},
   2/secondary/{Ch 8: a theme},
   3/greenhl/{Ch 9--10: basket and checkout},
   4/greenhl/{Ch 11--12: shipping, tax, vouchers},
   5/purplehl/{Ch 14: a payment gateway},
   6/purplehl/{Ch 15: customer accounts},
   7/purplehl/{Ch 16: hardened},
   8/tealhl/{Ch 19: your own module},
   9/goldhl/{Ch 20--21: deployed, indexed}}{
  \pgfmathsetmacro{\yy}{0.95+(\i-1)*0.78}
  \node[lay=\c] at (0,\yy) {\t};}

\draw[-{Stealth[length=2mm]}, primary!40, line width=1pt]
      (-3.3,-0.45) -- (-3.3,8.15);
\node[rotate=90, font=\scriptsize\bfseries, text=primary!70]
      at (-3.65,3.85) {sixteen weeks};

\node[rounded corners=3pt, fill=goldhl, text=white, font=\bfseries\footnotesize,
      minimum width=5.0cm, minimum height=0.8cm, align=center] at (0,8.55)
     {Ch 23\\Defended as the project};

% Each note sits at the mid-height of the layers it describes, which is why
% these y-values are not evenly spaced: the layers are, the groups are not.
\node[sidenote, text width=6.0cm] at (2.9,0)    {Docker Compose brings up the
   storefront, the database and a database console. Nothing is installed on
   your own machine.};
\node[sidenote, text width=6.0cm] at (2.9,1.34) {Real products, with real option
   combinations, in a schema you can explain --- and a shop that looks like
   somebody's.};
\node[sidenote, text width=6.0cm] at (2.9,2.90) {A customer can now get from the
   front page to a confirmed order, at the right price.};
\node[sidenote, text width=6.0cm] at (2.9,4.85) {That order is now paid for, by
   an identified customer, over a shop that has been attacked and hardened.};
\node[sidenote, text width=6.0cm] at (2.9,6.41) {The platform does something it
   was not shipped able to do.};
\node[sidenote, text width=6.0cm] at (2.9,7.19) {Somebody who is not you can
   reach it, and find it.};
\end{tikzpicture}}%
{What each part adds to the one shop. A lab skipped is a lab the next chapter
cannot build on --- Chapter 10 has nothing to check out if Chapter 9's basket
was never configured.}{shoplayers}

\begin{alertbox}[Before Week 2]
The stack runs in containers, so nothing in this course asks you to install a
web server, PHP or a database on your own machine. It does ask you to install
\textbf{one} thing: a container runtime. Appendix~A does it in one sitting and
lists what to do when it will not start, which on a university laptop it
sometimes will not.

\smallskip
Do it before Week~2. From Chapter~5 onward every chapter assumes a shop you can
reach in a browser, and a student without one spends the semester reading about
a course the rest of the class is taking.
\end{alertbox}

\newpage

%=====================================================================
\chapter*{Prerequisite Self-Check}
\addcontentsline{toc}{chapter}{Prerequisite Self-Check}
\markboth{Prerequisite Self-Check}{}

Twelve questions. If you can answer ten, begin at Chapter~1. If you cannot,
Appendix~B covers the PHP, SQL and template syntax this book assumes, your
\emph{Web Technologies} and \emph{Database Systems} notes cover the rest --- and
the gap is worth closing in Week~1 rather than Week~9.

\begin{enumerate}[leftmargin=2.4em, itemsep=4pt, topsep=4pt]
  \item What is the difference between an HTTP \texttt{GET} and an HTTP
        \texttt{POST}, and which one must a form that charges a card use? Why?
        \hfill{\scriptsize\color{neutral}(Web Technologies)}
  \item What do the status codes 200, 302, 404 and 500 mean, and which one does
        a browser follow without telling the user?
        \hfill{\scriptsize\color{neutral}(Web Technologies)}
  \item HTTP is stateless. How, then, does a server know that two requests came
        from the same visitor?
        \hfill{\scriptsize\color{neutral}(Web Technologies, \chref{basket})}
  \item Write a \texttt{SELECT} that joins two tables and returns one row per
        match. \hfill{\scriptsize\color{neutral}(Database Systems, Appendix~B)}
  \item What is a database index for, and what does adding one cost?
        \hfill{\scriptsize\color{neutral}(Database Systems, \chref{catalogue})}
  \item Distinguish a primary key from a foreign key, and say what the database
        guarantees about each.
        \hfill{\scriptsize\color{neutral}(Database Systems)}
  \item One product has many photographs. Which table holds the photographs, and
        why not the product table?
        \hfill{\scriptsize\color{neutral}(Database Systems, \chref{catalogue})}
  \item In an HTML form, what does the \texttt{name} attribute do, and where
        does its value arrive on the server?
        \hfill{\scriptsize\color{neutral}(Web Technologies)}
  \item A theme sets a heading to 18\,px and you want 24\,px. Give two ways to
        override it in CSS, and say which one you would not use.
        \hfill{\scriptsize\color{neutral}(Web Technologies, \chref{ux})}
  \item In PHP, what does \texttt{\$\_POST['qty']} contain, and why may you
        never use it in a query without doing something to it first?
        \hfill{\scriptsize\color{neutral}(Appendix~B, \chref{security})}
  \item What does HTTPS protect, and name one thing it does \emph{not} protect.
        \hfill{\scriptsize\color{neutral}(Information Security, \chref{security})}
  \item Open a terminal in a named directory, run a command, and read its
        output. If that sentence is not routine, start with Appendix~A.
        \hfill{\scriptsize\color{neutral}(Appendix~A)}
\end{enumerate}

\begin{conceptbox}
Notice what is \emph{not} on that list. You do not need to be a strong PHP
programmer: until Chapter~19 you will read far more code than you write, and
Appendix~B gives you enough syntax to read it. You do not need any business
background. You do not need a powerful machine.

\smallskip
What you do need is the habit of \emph{checking} --- opening the database to
see what the form actually wrote, reading the error log rather than guessing,
rerunning a measurement after changing one thing. Most of what goes wrong in
this course is invisible from the front page of the shop, and every chapter
tries to teach you where to look.
\end{conceptbox}

\newpage
\tableofcontents
\newpage
\listoffigures
\newpage

\pagenumbering{arabic}

%---------------------------------------------------------------------

\part{E-Commerce and Its Models}
%=====================================================================
% PART I OPENER
%=====================================================================
\thispagestyle{plain}
\structuralfigures{P1.}

\begin{center}
\begin{tcolorbox}[enhanced, width=0.92\textwidth, colback=primary!5,
  colframe=primary, boxrule=0pt, leftrule=5pt, arc=3pt, top=8pt, bottom=8pt]
\large\itshape\color{primary!80!black}
Before you install anything, four questions have to have answers: what electronic commerce actually changed, how a shop makes money, what shapes other than a shop exist, and --- given all that --- whether you should build a storefront, buy one, or extend one.
\end{tcolorbox}
\end{center}

\vspace{6pt}

\dsfigh{%
\begin{tikzpicture}[
  cbox/.style={rounded corners=3pt, draw=primary, line width=1pt, fill=primary!7,
               text=primary!80!black, font=\small\bfseries, align=center,
               minimum width=3.5cm, minimum height=1.0cm},
  hbox/.style={rounded corners=3pt, draw=primary, line width=1.8pt, fill=primary!16,
               text=primary!80!black, font=\small\bfseries, align=center,
               minimum width=3.5cm, minimum height=1.0cm},
  arr/.style={-{Stealth[length=2.4mm]}, primary!60, line width=1pt},
  note/.style={font=\scriptsize\itshape, text=primary!80!black, align=center,
               fill=white, inner sep=2pt}
]
  \node[cbox] (c0) at (0.00,0) {Ch.\ 1\\What Changed};
  \node[cbox] (c1) at (4.30,0) {Ch.\ 2\\Business Models};
  \node[cbox] (c2) at (8.60,0) {Ch.\ 3\\Auctions, Portals};
  \node[hbox] (c3) at (12.90,0) {Ch.\ 4\\Build or Extend};
  \draw[arr] (c0) -- (c1);
  \draw[arr] (c1) -- (c2);
  \draw[arr] (c2) -- (c3);
  \node[note] at (2.15,-1.0) {why it is different};
  \node[note] at (6.45,-1.0) {where the money is};
  \node[note] at (10.75,-1.0) {what else it can be};
  \node[font=\scriptsize\itshape, text=accent, align=center] at (12.90,1.15)
       {the decision the rest of the book carries out};
\end{tikzpicture}}%
{Reading path through Part~I. Chapters 1--3 are the business grounding; Chapter 4 turns it into an engineering decision, and is the one to argue with.}{part1map}

\newpage

%=====================================================================
\chapter{What E-Commerce Is, and What It Changed}\label{ch:what}
%=====================================================================
\locator{1}{Part I --- E-Commerce and Its Models}

\begin{outcomes}
By the end of this chapter you will be able to:
\begin{tight}
  \item \textbf{Define} electronic commerce, and \textbf{distinguish} it from
        the wider category of electronic business.
  \item \textbf{Name} the four functions every shop must perform, and
        \textbf{locate} each of them in a storefront you have not seen before.
  \item \textbf{Explain} why selling online is a software problem rather than a
        marketing one, giving two consequences that follow.
  \item \textbf{Distinguish} the parts of a transaction that electronic commerce
        actually changed from the parts it only moved.
  \item \textbf{Classify} a given online business by the four functions, and
        \textbf{identify} which of them it has chosen not to perform itself.
\end{tight}
\end{outcomes}

\begin{prereq}
Nothing in this book. You need to have used an online shop as a customer, which
is the only prerequisite this chapter has. Appendix~A sets up the software you
will need from \chref{stack} onward; start it this week, because it sometimes
needs a second attempt.
\end{prereq}

\begin{keyterms}
electronic commerce $\bullet$ electronic business $\bullet$ storefront
$\bullet$ catalogue $\bullet$ basket $\bullet$ checkout $\bullet$ fulfilment
$\bullet$ transaction $\bullet$ disintermediation $\bullet$ long tail
$\bullet$ friction $\bullet$ conversion $\bullet$ abandonment
$\bullet$ trust
\end{keyterms}

\section{A Shop With No Closing Time}

Consider a bookshop on a street in Bahawalpur. It opens at nine and closes at
nine. It holds perhaps four thousand titles, because that is what the shelves
take. Its customers are the people who can reach the street. If you want a book
it does not stock, the shopkeeper can order it, and you will come back next
week.

Now consider the same bookshop online. It is open at three in the morning. It
can list four hundred thousand titles, because a listing is a database row and
rows are cheap. Its customers are anyone with a connection. If you want a book
it does not stock, the website can know that instantly, and can offer to tell
you when it arrives.

Every one of those differences is real, and none of them is the interesting one.

The interesting difference is this: in the street shop, a human being stands
between the customer and every mistake. If the price label is wrong, the
shopkeeper notices. If you try to buy two of something when one is left, the
shopkeeper tells you. If you hand over a torn note, the shopkeeper decides. The
shop runs on a continuous supply of human judgement, and that judgement is
invisible precisely because it never fails conspicuously.

Online, there is no shopkeeper. Every one of those judgements has to have been
written down, in advance, as a rule --- and anything that was not written down
will happen.

\begin{definitionbox}[Electronic Commerce]
\term{Electronic commerce} is the buying and selling of goods and services over
a computer network, together with the transfer of money and data that completes
those transactions.

\smallskip
It is a subset of \term{electronic business}, which is the use of networks for
\emph{any} business process --- a supplier portal, an internal stock system, a
staff payroll. The distinguishing feature of e-commerce is the
\term{transaction}: value changes hands, and both parties are entitled to
expect that it did so exactly once.

\smallskip
That last clause is not pedantry. It is the reason \chref{checkout} is the
longest chapter in this book.
\end{definitionbox}

\section{The Four Functions}

Strip away the branding and every shop on the internet does the same four
things, in the same order. This is worth fixing now, because it is the skeleton
of the whole course: each function is a part of this book, and each is a cluster
of screens in the software you will install in \chref{stack}.

\dsfig{%
\begin{tikzpicture}[
  fn/.style={rounded corners=3pt, draw=#1, line width=1.1pt, fill=#1!8,
             text=#1!60!black, font=\bfseries\small, align=center,
             minimum width=2.9cm, minimum height=1.25cm},
  arr/.style={-{Stealth[length=2.4mm]}, primary!55, line width=1pt},
  lbl/.style={font=\scriptsize\itshape, text=neutral, align=center}
]
  \node[fn=secondary] (c1) at (0,0)    {Catalogue\\{\scriptsize\mdseries what is
       for sale}};
  \node[fn=greenhl]   (c2) at (3.9,0)  {Basket\\{\scriptsize\mdseries what you
       chose}};
  \node[fn=purplehl]  (c3) at (7.8,0)  {Payment\\{\scriptsize\mdseries money
       changes hands}};
  \node[fn=tealhl]    (c4) at (11.7,0) {Fulfilment\\{\scriptsize\mdseries the
       thing arrives}};
  \draw[arr] (c1) -- (c2);
  \draw[arr] (c2) -- (c3);
  \draw[arr] (c3) -- (c4);

  \node[lbl] at (0,-1.15)    {Part II};
  \node[lbl] at (3.9,-1.15)  {Part III};
  \node[lbl] at (7.8,-1.15)  {Part IV};
  \node[lbl] at (11.7,-1.15) {Part V};

  \node[font=\scriptsize\itshape, text=accent, align=center] at (5.85,1.30)
       {reversible without consequence $\;\longrightarrow\;$ not reversible};
  \draw[dashed, accent!60, line width=0.8pt] (5.85,0.85) -- (5.85,-1.5);
\end{tikzpicture}}%
{The four functions of any shop, and the parts of this book that build them.
The dashed line is the one that matters: everything to its left can be undone
by closing the browser tab, and nothing to its right can.}{fourfunctions}

\begin{conceptbox}[The dashed line is the whole subject]
A customer browsing a catalogue and filling a basket has committed to nothing.
They can close the tab and the shop is exactly as it was. Mistakes on that side
of the line cost a sale at worst.

\smallskip
The moment payment is attempted, that stops being true. Money may have moved
while the network dropped. An order may exist that the customer does not know
about, or the customer may believe in an order that does not exist. Stock may
have been reserved for someone who never paid.

\smallskip
Most of what distinguishes a real shop from a student exercise is the handling
of that line, and most student projects put all their effort on the left of it
because that is the half you can see.
\end{conceptbox}

\section{What Actually Changed}

It is easy to overstate what e-commerce changed. A great deal of it is the same
commerce, conducted through a different window. Three things, though, genuinely
changed, and they explain most of what the rest of this book worries about.

\subsection{The cost of listing fell to nearly nothing}

Shelf space is expensive and finite; a database row is neither. A street shop
must stock what sells often, because an item that sells twice a year occupies
the same shelf as one that sells twice a week. An online shop has no such
pressure, and can profitably carry items that sell twice a year --- the
\term{long tail}.

That is a business observation with an engineering consequence that this course
will measure: a catalogue of four hundred thousand items does not behave like a
catalogue of four thousand. \chref{catalogue} puts fifty thousand products into
your shop and times the result, and the answer is not comfortable.

\subsection{Some middlemen became unnecessary, and others appeared}

\term{Disintermediation} is the removal of intermediaries: a manufacturer can
sell directly to you, and the distributor and the retailer are not needed.

The prediction made in the 1990s was that this would flatten commerce. What
happened instead was that old intermediaries were replaced by new and larger
ones --- marketplaces, payment processors, delivery networks, search engines.
A shop today typically pays a marketplace for access to customers, a gateway
for the ability to take money, and a courier to deliver, and each of those is a
dependency with a price and a failure mode. \chref{payments} is about one of
them in detail.

\subsection{Trust had to be manufactured}

In a street shop, trust is cheap. The shop is physically there; you can see the
goods; if it cheats you, you know where it is.

Online, none of that holds, and every substitute has to be built deliberately:
a visible address and a returns policy, a padlock in the address bar, reviews,
an order confirmation that arrives within seconds, a refund that works. These
are not decoration --- they are the mechanism that replaces a shopfront, and
each of them is a feature somebody has to implement and maintain.

\begin{examplebox}[Three signals, and what each costs to fake]
A customer deciding whether to buy from an unfamiliar shop is reading signals.
It is worth knowing which ones are hard to fake and which are not.

\begin{tight}
  \item \textbf{A professional-looking design.} Costs almost nothing to fake;
        a template is free. Carries little information.
  \item \textbf{A valid certificate and HTTPS.} Costs nothing now that
        certificates are free --- which is why its absence is damning and its
        presence proves very little. \chref{security} returns to exactly what
        it does and does not prove.
  \item \textbf{A confirmation email, with an order number, that arrives in
        four seconds.} Hard to fake, because it requires the shop to actually
        have the machinery. This is one reason \chref{checkout} treats the
        confirmation as part of the transaction rather than as an
        afterthought.
\end{tight}
\end{examplebox}

\section{Why This Is a Software Problem}

The common way to teach this subject is as business studies with screenshots.
That gets the emphasis wrong, and the published outline this course follows ---
see the Course Specification --- gets it right: eighteen of its twenty-two
clauses are about building, configuring and running the software.

Here is the argument in one paragraph. A shop is a set of rules that must hold
without supervision: this price applies to this customer in this quantity in
this region on this date; this voucher may be used once; this stock is reserved
for eleven minutes; this payment either completed or it did not. Every one of
those is a statement about state, every one of them has an edge, and nobody is
watching. That is not a marketing problem. It is the same problem as any other
concurrent, stateful, partially-trusted system, and it is solved with the same
tools.

\begin{alertbox}[The two failures that are specific to this subject]
Almost everything that goes badly wrong in an online shop is one of two things,
and both are invisible from the front page.

\smallskip
\textbf{The state got out of step.} Money taken, no order recorded. Order
recorded, stock never reduced. Refund issued twice. These are not rare, and they
do not announce themselves --- the shop looks fine, and the accounts do not
balance at the end of the month. \chref{checkout} and \chref{orders} are about
not causing them.

\smallskip
\textbf{Something was trusted that should not have been.} A price taken from a
form field rather than from the database. A quantity that was allowed to be
negative. An administrative page that was never actually checked for a login.
\chref{security} is a whole chapter about this, and its laboratory finds several
of them in the shop you will have built by then.
\end{alertbox}

\section{Reading a Shop You Have Not Seen Before}

The skill this chapter is really teaching is a diagnostic one, and it is worth
practising before you build anything. Given any online business, you should be
able to say within a few minutes which of the four functions it performs itself,
which it has handed to somebody else, and what that choice costs it.

\begin{worked}[Classifying three businesses by the four functions]
\textbf{A clothing brand selling from its own site.}
Catalogue: its own. Basket: its own. Payment: a gateway --- it does not hold a
banking licence, so this is handed over. Fulfilment: usually a courier.
\emph{Performs two of four; the two that define it.}

\smallskip
\textbf{The same brand selling only through a large marketplace.}
Catalogue: the marketplace's, in the marketplace's format. Basket: the
marketplace's. Payment: the marketplace's. Fulfilment: possibly the
marketplace's warehouse. \emph{Performs none of the four.} It has bought
access to customers and sold its relationship with them; it does not know who
they are, cannot email them, and can be delisted. The commission is the
visible price and is not the largest one.

\smallskip
\textbf{A food delivery platform.}
Catalogue: aggregated from restaurants. Basket: its own. Payment: its own.
Fulfilment: its own riders --- which is the expensive one, and the one that is
hard to copy. \emph{Performs three and a half of four}, and the half it does
not perform (the cooking) is the one with the thinnest margin.

\smallskip
\textbf{The general rule.} The function a business insists on performing itself
is usually the one it believes is its advantage. Ask which function a business
refuses to outsource, and you have found what it thinks it is.
\end{worked}

\begin{pitfall}
\begin{tight}
  \item \textbf{Confusing e-commerce with having a website.} A site that shows a
        catalogue and a phone number is a brochure. The transaction is the
        thing; without it, none of the hard problems in this book arise.
  \item \textbf{Assuming the hard part is the shopping cart.} It is not, and it
        has not been for twenty years --- you will configure one in an
        afternoon in \chref{basket}. The hard parts are everything after the
        dashed line in \dsref{fourfunctions}.
  \item \textbf{Treating trust signals as design.} A returns policy is not a
        page; it is a promise that somebody has to keep, and a refund path
        somebody has to build (\chref{orders}).
  \item \textbf{Believing disintermediation removed the middlemen.} It replaced
        them. Count your dependencies before claiming you have none.
  \item \textbf{Saying ``we will add payments later''.} Everything to the right
        of the dashed line constrains everything to its left. Deciding how you
        take money late is how projects discover their data model is wrong.
\end{tight}
\end{pitfall}

\begin{lab}[Reading four real shops]
No software is needed for this one. \chref{stack} installs the stack; this week
you are learning to look.

\smallskip
Choose \textbf{four} online businesses that you can actually reach, of
deliberately different shapes --- for example a single-brand store, a large
general marketplace, a services business that takes bookings rather than
shipping goods, and a local business selling through social media only.

\smallskip
For each one, as a customer and without buying anything, record:

\begin{tightnum}
  \item \textbf{The four functions.} Who performs each --- the business, or
        somebody else? How can you tell? For payment, the giveaway is usually
        whether the address bar changes during checkout.
  \item \textbf{Where the catalogue ends.} How many categories deep does it go?
        Can a product appear in two categories? Does it have options --- size,
        colour --- and are they separate products or one product with choices?
        You are looking at a data model through a window; \chref{catalogue} will
        show you the model itself.
  \item \textbf{The furthest point you can reach without an account.} Can you
        fill a basket? Reach the checkout? See the shipping cost before
        committing? Write down the exact step at which it stopped.
  \item \textbf{Three trust signals}, and for each, whether it is cheap or
        expensive to fake.
  \item \textbf{One thing that is broken or confusing.} There is always one.
        Describe it precisely enough that somebody could fix it.
\end{tightnum}

\smallskip
\textbf{Write up one page} comparing the four, ending with a single sentence
for each: what does this business think it is? Keep it --- \chref{planning}
asks you to do the same exercise for the shop you are going to build, and it is
much easier the second time.

\smallskip
\textbf{What to notice.} The four shops will differ far more in what they
\emph{refuse} to do than in what they do. Everyone has a catalogue. Not everyone
will tell you the delivery charge before you have handed over your address, and
the ones that will not have usually decided that on purpose.
\end{lab}

\begin{checkpoint}
\begin{tightnum}
  \item A university lets students view results and download transcripts
        online, but fee payment happens at a bank counter. Is this electronic
        commerce, electronic business, both or neither? Justify your answer from
        the definitions.
  \item Give one judgement a human shopkeeper makes without noticing, and
        describe precisely what has to exist in software to replace it.
  \item A marketplace charges a 15\% commission. A seller calculates that
        running its own site would cost less than 15\% of revenue and concludes
        it should leave. Name two costs the seller has not counted, and one
        thing it would gain that is not money.
\end{tightnum}
\end{checkpoint}

\begin{chaptersummary}
\begin{tight}
  \item Electronic commerce is the buying and selling of goods and services over
        a network, \emph{including} the transfer of money. The transaction is
        what distinguishes it from electronic business generally.
  \item Every shop performs four functions in order: catalogue, basket, payment,
        fulfilment. They are Parts II to V of this book.
  \item The line between basket and payment is the one that matters. Everything
        before it is reversible by closing a tab; nothing after it is.
  \item Three things genuinely changed: the cost of listing collapsed, old
        intermediaries were replaced by new and larger ones, and trust stopped
        being free and had to be manufactured.
  \item The absent shopkeeper is the engineering problem. Every judgement a
        human would have made has to have been written down in advance, and
        whatever was not written down will eventually happen.
  \item Almost all serious failures are one of two kinds: state that got out of
        step, or something trusted that should not have been.
  \item To read an unfamiliar business, ask which of the four functions it
        performs itself. The one it refuses to outsource is what it believes it
        is.
\end{tight}
\end{chaptersummary}

\begin{reviewq}
  \item \textbf{(MCQ)} Which of these is electronic business but \emph{not}
        electronic commerce? (a)~A customer buying a phone from a web shop;
        (b)~a supplier checking a manufacturer's stock levels through a portal;
        (c)~a customer paying a utility bill through an app; (d)~a marketplace
        taking commission on a sale.
  \item \textbf{(MCQ)} In \dsref{fourfunctions}, the dashed line separates:
        (a)~front end from back end; (b)~free actions from paid actions;
        (c)~what a customer can undo by closing the tab from what they cannot;
        (d)~catalogue data from order data.
  \item \textbf{(Short)} Explain the long tail, and give one engineering
        consequence of it that a shop owner would not anticipate.
  \item \textbf{(Short)} ``Disintermediation removed the middleman.'' Argue
        against this with a concrete example.
  \item \textbf{(Short)} A friend says the hard part of building an online shop
        is the shopping cart. Give the strongest version of why they are wrong,
        and name the two failure modes you would be more worried about.
  \item \textbf{(Applied)} A bookshop with four thousand titles in a shop wants
        to list eighty thousand online, most of which it does not hold in stock.
        Identify three decisions this forces that would not arise at four
        thousand titles, and say for each which of the four functions it
        affects.
  \item \textbf{(Applied)} Take any one of the four businesses from this
        chapter's laboratory. Describe, in enough detail that it could be
        checked, an experiment a customer could run from outside that would
        reveal whether the shop reserves stock during checkout. State what
        result would mean ``yes'', what would mean ``no'', and what would be
        inconclusive.
\end{reviewq}

%=====================================================================
\chapter{Business Models and Revenue Models}\label{ch:models}
%=====================================================================
\locator{1}{Part I --- E-Commerce and Its Models}

\begin{outcomes}
By the end of this chapter you will be able to:
\begin{tight}
  \item \textbf{Distinguish} a business model from a revenue model, and
        \textbf{explain} why a venture can have a sound one and a hopeless
        other.
  \item \textbf{Classify} an online business by participant --- B2C, B2B, C2C,
        C2B --- and say what each classification predicts about the software.
  \item \textbf{Name} the principal revenue models and \textbf{identify} which
        one a given shop is using.
  \item \textbf{Explain} the unit economics of a single order, and
        \textbf{compute} the contribution margin from a set of figures.
  \item \textbf{Argue} which revenue model a given venture should adopt, and
        \textbf{state} the software consequences of that choice.
\end{tight}
\end{outcomes}

\begin{prereq}
\chref{what}: the four functions, and the habit of asking which of them a
business performs itself.
\end{prereq}

\begin{keyterms}
business model $\bullet$ revenue model $\bullet$ value proposition $\bullet$
B2C $\bullet$ B2B $\bullet$ C2C $\bullet$ C2B $\bullet$ marketplace $\bullet$
commission $\bullet$ subscription $\bullet$ freemium $\bullet$ advertising
$\bullet$ affiliate $\bullet$ unit economics $\bullet$ contribution margin
$\bullet$ customer acquisition cost $\bullet$ lifetime value $\bullet$
network effect
\end{keyterms}

\section{Two Questions That Get Confused}

Almost every failed online venture answered one of these two questions well and
the other not at all.

\begin{definitionbox}[Business Model and Revenue Model]
A \term{business model} says \emph{what value you create and for whom}: who the
customer is, what problem is solved, and why they would choose you rather than
the alternative. It is the answer to ``why does this deserve to exist?''

\smallskip
A \term{revenue model} says \emph{how the money reaches you}: who pays, for
what, how often, and how much. It is the answer to ``who writes the cheque?''

\smallskip
They are independent. The same business model --- say, connecting people who
want a ride with people who will drive --- can be monetised by commission per
ride, by a driver subscription, by advertising, or by charging nothing and
selling the data. Those four are different companies with different software.
\end{definitionbox}

The confusion matters because the two fail differently and the failures look
alike from outside. A venture with no business model and a clear revenue model
builds something nobody wants and charges correctly for it. A venture with a
strong business model and no revenue model builds something people love and
cannot pay for it --- and it is usually the second that students build.

\section{Classifying by Who Is at Each End}

The oldest classification asks only who is transacting with whom. It is coarse,
and it is still the first question worth asking, because it predicts a
surprising amount about the software.

\rowcolors{2}{white}{lightbg}
\begin{longtable}{@{}L{1.4cm}L{5.2cm}L{8.2cm}@{}}
\toprule
\rowcolor{primary!15}
\textbf{Type} & \textbf{Who to whom} & \textbf{What it implies for the
software} \\
\midrule
\endhead
\textbf{B2C} & A business sells to a consumer & Many small orders, card and
wallet payment, a catalogue tuned for browsing, and heavy traffic at
unpredictable moments. The default assumption of almost every storefront
platform, including the one this course uses \\
\textbf{B2B} & A business sells to another business & Few large orders,
negotiated and customer-specific prices, credit terms rather than immediate
payment, purchase orders, approval chains, and quantity breaks. Almost
everything in \chref{checkout} changes, and most of \chref{discounts} \\
\textbf{C2C} & A consumer sells to another consumer & The operator holds no
stock and never owns the goods. The hard parts move to trust, dispute
resolution and holding money in escrow --- \chref{auctions} is about this \\
\textbf{C2B} & A consumer sells to a business & Freelance marketplaces, stock
photography, influencer deals. The ``product'' is a person's time or licence,
which breaks the stock model entirely \\
\bottomrule
\end{longtable}

\begin{conceptbox}[Why the B2B/B2C split is the one that bites]
A platform built for B2C assumes one price per product. B2B does not have one
price per product --- it has a price per product \emph{per customer}, often
negotiated, often with quantity breaks, sometimes expiring.

\smallskip
This is the most common way a student project, or a real small business, picks
the wrong tool. Everything looks fine until the first wholesale customer asks
for their agreed rate, and the answer turns out to require a different data
model rather than a setting. \chref{planning} makes you ask this question
before you choose, which is the only time asking it is cheap.
\end{conceptbox}

\section{The Revenue Models}

There are not very many, and most businesses combine two or three.

\dsfig{%
\begin{tikzpicture}[
  rm/.style={rounded corners=3pt, draw=#1, line width=1pt, fill=#1!8,
             text=#1!60!black, font=\bfseries\scriptsize, align=center,
             text width=2.7cm, minimum height=1.5cm},
  lbl/.style={font=\scriptsize\itshape, text=neutral, align=center,
              text width=2.9cm}
]
  \node[rm=secondary] (a) at (0,0)    {Sales\\{\scriptsize\mdseries margin on
       each item sold}};
  \node[rm=greenhl]   (b) at (3.3,0)  {Commission\\{\scriptsize\mdseries a cut
       of somebody else's sale}};
  \node[rm=purplehl]  (c) at (6.6,0)  {Subscription\\{\scriptsize\mdseries paid
       access, repeating}};
  \node[rm=goldhl]    (d) at (9.9,0)  {Advertising\\{\scriptsize\mdseries
       attention sold to a third party}};
  \node[rm=tealhl]    (e) at (13.2,0) {Affiliate\\{\scriptsize\mdseries paid for
       sending a buyer}};

  \node[lbl] at (0,-1.55)    {you hold stock};
  \node[lbl] at (3.3,-1.55)  {you hold none};
  \node[lbl] at (6.6,-1.55)  {predictable; churn is the risk};
  \node[lbl] at (9.9,-1.55)  {needs scale before it pays};
  \node[lbl] at (13.2,-1.55) {no customer of your own};
\end{tikzpicture}}%
{The five revenue models a shop can use, and the single fact about each that
most shapes the software. Most real businesses run two of them at once.}%
{revenuemodels}

\subsection{Sales margin}

You buy at one price and sell at another. The oldest model, and the one the
storefront software in this course is built around by default.

Its distinguishing feature, from a software point of view, is that \emph{you
hold stock}. Stock means a quantity that can be wrong, can be reserved, can go
negative under concurrency, and must be reconciled against physical reality.
\chref{catalogue} and \chref{orders} spend real effort on this.

\subsection{Commission}

Somebody else sells; you take a percentage for having made it possible. This is
the marketplace model, and it is attractive because you hold no stock and your
costs barely rise with volume.

What you take on instead is everything in \chref{auctions}: two parties who do
not trust each other, money that must be held until both are satisfied, and
disputes that someone has to adjudicate. You have also acquired a problem the
sales model does not have --- you must attract \emph{both} sides, and neither
arrives before the other.

\subsection{Subscription}

The customer pays regularly for continued access or delivery. Revenue becomes
predictable, which is why investors like it and why so much of the web has
moved to it.

The software consequence is recurring billing, which is substantially harder
than one-off payment: cards expire, payments fail silently, customers pause and
resume, and the price may change mid-contract. \chref{gateway} explains why
a stored card is a different kind of integration from a single charge.

\subsection{Advertising and affiliate}

Advertising sells the attention of your visitors to somebody else. It needs
large traffic before it pays anything at all, which makes it a poor first
revenue model and a common fantasy.

Affiliate revenue is the mirror image: you send a visitor to somebody else's
shop and are paid if they buy. Cheap to start, and it leaves you without a
customer --- you never learn who bought, cannot serve them again, and the
relationship belongs to the merchant. \chref{seo} is where this becomes a
strategy rather than an accident.

\section{The Arithmetic of One Order}

A revenue model is a claim that can be checked, and the check is simpler than
students expect. Take one order and follow the money out of it.

\begin{worked}[Does this order make money?]
A shop sells a pair of shoes at \textbf{Rs.~4{,}500}, delivered anywhere in
Pakistan, with free returns. Work out what is left.

\smallskip
\textbf{Step 1 --- what came in.}\\
Revenue per order: \textbf{Rs.~4{,}500}.

\smallskip
\textbf{Step 2 --- the cost of the thing itself.}\\
The shoes cost the shop Rs.~2{,}800 to buy.\\
Remaining: $4{,}500 - 2{,}800 = \textbf{Rs.~1{,}700}$ (a gross margin of 38\%).

\smallskip
\textbf{Step 3 --- the cost of taking the money.}\\
The gateway charges 2.5\% plus Rs.~15 per transaction:
$0.025 \times 4{,}500 + 15 = 112.5 + 15 = \text{Rs.~127.50}$.\\
Remaining: $1{,}700 - 127.50 = \textbf{Rs.~1{,}572.50}$.

\smallskip
\textbf{Step 4 --- the cost of delivery.}\\
Courier, Rs.~250, absorbed by the shop because delivery was advertised free.\\
Remaining: $1{,}572.50 - 250 = \textbf{Rs.~1{,}322.50}$.

\smallskip
\textbf{Step 5 --- returns, which are not optional.}\\
20\% of shoe orders come back. A return costs the courier fee again (Rs.~250)
and the item cannot always be resold; assume a further Rs.~200 of loss on
average. Expected cost per \emph{order placed}:
$0.20 \times (250 + 200) = \text{Rs.~90}$.\\
Remaining: $1{,}322.50 - 90 = \textbf{Rs.~1{,}232.50}$.

\smallskip
\textbf{Step 6 --- the cost of finding the customer.}\\
The shop spends Rs.~60{,}000 a month on advertising and gets 50 orders from it:
a \term{customer acquisition cost} of $60{,}000 / 50 = \text{Rs.~1{,}200}$.\\
Remaining: $1{,}232.50 - 1{,}200 = \textbf{Rs.~32.50}$.

\smallskip
\textbf{The answer.} The order that looked like a 38\% margin earns
\textbf{Rs.~32.50}, which is \textbf{0.7\%} of the sale price --- and that is
before rent, salaries, hosting or tax. One extra return, or a 3\% gateway
instead of 2.5\%, and it is a loss.

\smallskip
\textbf{What this is really telling you.} The two numbers that decided it were
\emph{returns} and \emph{acquisition cost}, and neither appears anywhere in the
product's price. A shop that measures only gross margin believes it is making
38\%. This is why \chref{dashboard} is careful about which numbers on a
dashboard are worth acting on.
\end{worked}

\begin{definitionbox}[Contribution Margin, CAC and LTV]
\term{Contribution margin} is what one order contributes after all the costs
that vary with that order --- goods, payment fees, delivery, expected returns.
It is what is left to pay the costs that do not vary.

\smallskip
\term{Customer acquisition cost} (CAC) is the marketing spend divided by the
customers it produced.

\smallskip
\term{Lifetime value} (LTV) is the total contribution margin a customer
produces before they stop buying.

\smallskip
The test a business must pass is $\text{LTV} > \text{CAC}$, by a comfortable
multiple. The worked example above passes only if the customer orders again ---
its CAC is nearly its whole first-order margin, which is survivable for a
business with repeat customers and fatal for one without.
\end{definitionbox}

\begin{alertbox}[Where the network effect comes from, and where it does not]
A \term{network effect} exists when the service gets better for each user as
more users join. A marketplace has one: more sellers make it more useful to
buyers, which attracts more sellers.

\smallskip
A shop selling its own goods has \emph{no} network effect. Your thousandth
customer makes the shop no better for your first. Volume may buy you better
purchase prices, which is a real advantage, but it is not a network effect and
it does not protect you the same way.

\smallskip
This matters because the two justify completely different strategies. Running
at a loss to grow is rational when a network effect will eventually make you
defensible, and is just losing money when it will not. Plenty of ventures have
adopted the first strategy while having the second business.
\end{alertbox}

\begin{pitfall}
\begin{tight}
  \item \textbf{Naming a business model and thinking you have named a revenue
        model.} ``We are a marketplace for tutors'' says nothing about who pays.
  \item \textbf{Quoting gross margin as if it were profit.} The worked example
        went from 38\% to 0.7\% without anything unusual happening.
  \item \textbf{Forgetting returns.} In clothing and footwear they run at 20--40\%
        and they are the difference between a viable shop and a dead one.
  \item \textbf{Advertising as a first revenue model.} It pays at scale and
        nothing before it, so a new venture choosing it has chosen to earn
        nothing for a long time.
  \item \textbf{Assuming the platform is model-agnostic.} It is not. A B2C
        storefront has one price per product, and discovering that during
        implementation is expensive --- see \chref{planning}.
  \item \textbf{Treating ``free delivery'' as marketing.} It is a cost of
        Rs.~250 per order that appears in no price list and is paid out of the
        margin.
\end{tight}
\end{pitfall}

\begin{lab}[The model behind the shop you will build]
Still no software --- \chref{stack} starts that. This laboratory produces the
one-page brief that \chref{planning} turns into a decision and
\chref{project} is eventually marked against, so it is worth real effort.

\smallskip
\textbf{Choose a shop to build.} It must be plausible, small, and local enough
that you can check your assumptions: a bookshop, a bakery, a tailor, a
spare-parts dealer, a plant nursery. Do not choose ``an Amazon competitor''.

\smallskip
Then produce, on one page:

\begin{tightnum}
  \item \textbf{The business model in two sentences.} Who is the customer, what
        problem is solved, and why you rather than the shop down the road or
        the marketplace they already use.
  \item \textbf{The classification}: B2C, B2B, C2C or C2B --- and if you are
        tempted to say ``both'', say which one the \emph{software} must assume,
        because it can only assume one.
  \item \textbf{The revenue model}, from \dsref{revenuemodels}, with the reason
        for choosing it over the next most plausible one.
  \item \textbf{The arithmetic of one order}, following the six steps of the
        worked example with your own figures. Find real numbers where you can:
        gateway fees and courier rates are published, and a shopkeeper will
        usually tell you their return rate if you ask politely. Mark every
        number you guessed.
  \item \textbf{The two numbers that decide it.} From your own arithmetic, name
        the two inputs the result is most sensitive to, and say what would have
        to be true for the shop to lose money.
\end{tightnum}

\smallskip
\textbf{What to notice.} Almost everyone finds that the contribution margin is
far thinner than they expected, and that the deciding numbers are not the ones
they thought they were building a business around. That discovery is the point
of the exercise, and it is much cheaper to make here than in Chapter~23.
\end{lab}

\begin{checkpoint}
\begin{tightnum}
  \item A food delivery app charges restaurants 20\% commission, charges
        customers a delivery fee, and sells promoted placement in search
        results. Name the business model and all the revenue models, and say
        which single one you would expect to be most profitable per unit.
  \item A shop has a contribution margin of Rs.~400 per order and a customer
        acquisition cost of Rs.~1{,}500. Under what condition is this business
        viable? State it as an inequality and say what would need measuring.
  \item Explain why a B2B shop cannot usually be built by configuring a B2C
        platform, naming the specific assumption that breaks.
\end{tightnum}
\end{checkpoint}

\begin{chaptersummary}
\begin{tight}
  \item A business model says what value is created and for whom; a revenue
        model says who pays. They are independent, and a venture can have a
        good one and a fatal other.
  \item B2C, B2B, C2C and C2B classify by who transacts. The split that most
        affects the software is B2B against B2C: B2B has a price per customer,
        not a price per product.
  \item The five revenue models are sales margin, commission, subscription,
        advertising and affiliate. Holding stock or not is the fact about each
        that most shapes the system you have to build.
  \item Follow one order through six deductions --- goods, payment fee,
        delivery, returns, acquisition --- and the margin usually collapses. In
        the worked example, 38\% became 0.7\%.
  \item The viability test is $\text{LTV} > \text{CAC}$ by a comfortable
        multiple, which for a thin first-order margin means repeat custom is
        not optional.
  \item Marketplaces have network effects; shops selling their own goods do
        not, and strategies that assume one are expensive when it is absent.
\end{tight}
\end{chaptersummary}

\begin{reviewq}
  \item \textbf{(MCQ)} ``We connect home cooks with office workers who want
        lunch.'' This is: (a)~a revenue model; (b)~a business model; (c)~both;
        (d)~neither.
  \item \textbf{(MCQ)} Which revenue model most changes the \emph{payment}
        integration a shop needs? (a)~Sales margin; (b)~affiliate;
        (c)~subscription; (d)~advertising.
  \item \textbf{(Short)} Distinguish gross margin from contribution margin, and
        give two costs that appear in one and not the other.
  \item \textbf{(Short)} Explain why a marketplace has a harder launch than a
        shop, despite holding no stock.
  \item \textbf{(Short)} A clothing shop reports 45\% gross margin and is losing
        money. Give the two most likely explanations, and say what you would
        measure first.
  \item \textbf{(Applied)} A shop sells a product at Rs.~2{,}000 that costs it
        Rs.~1{,}100. Gateway: 2.9\% + Rs.~20. Delivery: Rs.~200, charged to the
        customer at Rs.~150. Returns: 12\%, each costing Rs.~200 in courier and
        Rs.~150 in unsellable stock. Compute the contribution margin per order,
        showing each step. Then find the return rate at which the shop breaks
        even.
  \item \textbf{(Applied)} Take the venture from this chapter's laboratory and
        argue for a \emph{different} revenue model than the one you chose. Give
        the strongest case you can, then say which software decisions in Parts
        III and IV would have to change if you adopted it.
\end{reviewq}

%=====================================================================
\chapter{Auctions, Portals and Communities}\label{ch:auctions}
%=====================================================================
\locator{1}{Part I --- E-Commerce and Its Models}

\begin{outcomes}
By the end of this chapter you will be able to:
\begin{tight}
  \item \textbf{Describe} the four common auction formats and \textbf{explain}
        what each one does to a bidder's incentive to bid honestly.
  \item \textbf{Identify} the engineering problems an auction creates that a
        fixed-price shop does not have.
  \item \textbf{Explain} the role of escrow, and \textbf{describe} what a
        marketplace operator must do when two users disagree.
  \item \textbf{Distinguish} a portal, a marketplace and a community, and
        \textbf{state} where each one's revenue comes from.
  \item \textbf{Explain} the two-sided launch problem and \textbf{evaluate}
        strategies for escaping it.
\end{tight}
\end{outcomes}

\begin{prereq}
\chref{models}: the revenue models, especially commission, and the note on
network effects at the end.
\end{prereq}

\begin{keyterms}
auction $\bullet$ English auction $\bullet$ Dutch auction $\bullet$ sealed-bid
auction $\bullet$ Vickrey auction $\bullet$ reserve price $\bullet$ proxy
bidding $\bullet$ sniping $\bullet$ shill bidding $\bullet$ escrow $\bullet$
dispute resolution $\bullet$ reputation system $\bullet$ portal
$\bullet$ marketplace $\bullet$ two-sided market $\bullet$ chicken-and-egg
problem
\end{keyterms}

\section{Why This Chapter Is Here Rather Than at the End}

The published outline lists ``Auctions, Portals, and Communities'' second from
last, after security and deployment. That is where it appears in the book whose
contents page the outline was copied from, and it is the wrong place: these are
not an afterthought to building a shop, they are \emph{alternative shapes} a
business can take, and the decision between them belongs in Part~I alongside
the other models.

So this chapter sits where it belongs. It is also the chapter that explains why
the rest of the book assumes a fixed-price shop: the alternatives are
genuinely harder, and it is worth knowing exactly how.

\section{Letting the Buyer Set the Price}

A fixed-price shop answers ``what is this worth?'' before the customer
arrives. An auction refuses to answer and lets the buyers decide.

That is attractive when the seller genuinely does not know the price --- a
second-hand item, a one-off, something whose value depends on who wants it. It
is unnecessary when the seller does know, which is why most commerce is not
conducted by auction.

\begin{definitionbox}[Four Auction Formats]
\term{English auction} --- ascending, open. The price rises until one bidder
remains. What nearly everyone means by ``auction''.

\smallskip
\term{Dutch auction} --- descending, open. The price starts high and falls
until somebody accepts. Fast, and used for perishable goods where speed matters
more than the last rupee.

\smallskip
\term{Sealed-bid first-price} --- everyone submits one hidden bid; the highest
wins and pays what they bid. Standard for tenders and contracts.

\smallskip
\term{Vickrey auction} (sealed-bid second-price) --- the highest bid wins but
pays the \emph{second}-highest price. Its remarkable property is that bidding
your true valuation is the best strategy whatever anyone else does, so there is
nothing to be gained from strategising.
\end{definitionbox}

\begin{conceptbox}[Why the second-price rule removes the guessing]
Suppose an item is worth Rs.~10{,}000 to you.

\smallskip
In a \emph{first-price} auction, bidding Rs.~10{,}000 guarantees you gain
nothing even if you win --- so you must shade your bid downward, and how far
depends on guessing what others will do. Everyone does this, and everyone
guesses.

\smallskip
In a \emph{second-price} auction, bid Rs.~10{,}000. If the next highest bid is
Rs.~8{,}000, you win and pay Rs.~8{,}000, gaining Rs.~2{,}000. If it is
Rs.~11{,}000, you lose --- and you are glad, because winning would have cost
more than the item was worth to you. Raising your bid above your valuation can
only make you win auctions you did not want; lowering it can only make you lose
auctions you did want. Telling the truth is simply best.

\smallskip
\textbf{Where you have met this already.} The ``proxy bidding'' on consumer
auction sites --- you enter a maximum and the site bids on your behalf up to it
--- is a Vickrey auction wearing an English auction's clothes. The winner pays
one increment above the second-highest maximum, not their own.
\end{conceptbox}

\section{What an Auction Costs You in Engineering}

A fixed-price product is a row that changes rarely. An auction is a small
real-time system, and it brings four problems the rest of this book does not
have to solve.

\rowcolors{2}{white}{lightbg}
\begin{longtable}{@{}L{3.6cm}L{11.4cm}@{}}
\toprule
\rowcolor{primary!15}
\textbf{Problem} & \textbf{Why a fixed-price shop does not have it} \\
\midrule
\endhead
\textbf{Time is part of the state} & A listing must close at an exact moment,
unattended, and settle correctly whether or not anyone is looking. A product
page has no deadline. This needs a scheduler, and a scheduler that did not run
is a silent failure \\
\textbf{Concurrent bids} & Two bids arriving in the same instant must produce
one consistent outcome. This is the classic lost-update problem and it needs a
transaction and a lock, not a quick read-then-write \\
\textbf{Sniping} & A bid placed one second before closing wins without giving
anyone a chance to respond. Fixes exist --- extend the closing time on a late
bid --- and each changes the rules of the auction, which is a design decision
rather than a bug fix \\
\textbf{Shill bidding} & The seller, or a friend, bids to push the price up
with no intention of buying. Detecting it means analysing bidding patterns
across accounts, and it is fraud rather than a feature gap \\
\bottomrule
\end{longtable}

\begin{alertbox}[The lost update, concretely]
Two bidders submit Rs.~5{,}000 at the same instant on an item currently at
Rs.~4{,}500.

\smallskip
A naive implementation reads the current high bid, checks that the new bid
beats it, and writes. Interleaved, both read Rs.~4{,}500, both pass the check,
and both write Rs.~5{,}000 --- the second overwriting the first. One bidder's
money is now owed on an item they did not win, and the audit trail shows a
single bid.

\smallskip
The fix is the one \chref{checkout} uses for stock: do the check and the write
in one atomic statement, inside a transaction, with the row locked. ``Read,
decide, write'' is three steps, and anything can happen between them.
\end{alertbox}

\section{Holding Other People's Money}

An auction site that merely introduces a buyer to a seller has introduced two
strangers and left. Whoever moves first --- sends the money, or sends the goods
--- carries all the risk.

\term{Escrow} is the operator's answer: it holds the buyer's money, tells the
seller to ship, and releases the money when the buyer confirms receipt or a
time limit expires.

\begin{definitionbox}[Escrow and Dispute Resolution]
\term{Escrow} is the holding of a payment by a trusted third party until the
conditions of the transaction are met.

\smallskip
It makes the operator a custodian of money that is not theirs, which in most
jurisdictions --- Pakistan included --- carries regulatory obligations
(\chref{legal}). It also makes the operator the \term{dispute resolution}
authority, because somebody has to decide who is telling the truth when the
buyer says the item never arrived and the seller says it did.

\smallskip
That adjudication cannot be automated away, and it is why marketplace
commissions are higher than payment processing fees: the commission buys
arbitration, not just plumbing.
\end{definitionbox}

\begin{examplebox}[Reputation as the cheap substitute, and what it cannot do]
Escrow is expensive. The cheaper substitute is a \term{reputation system}:
record how every past transaction went and let future counterparties read it.

\smallskip
It works, within limits, and the limits are worth knowing:

\begin{tight}
  \item \textbf{A new seller has no reputation}, so the system gives no
        protection in exactly the case where it is most needed.
  \item \textbf{Reputation can be bought} --- a hundred small honest sales, then
        one large dishonest one. The score is an average; the fraud is a
        maximum.
  \item \textbf{Retaliation distorts it.} If each side rates the other, a buyer
        who reports a problem risks a bad rating in return, so both sides learn
        to stay quiet and the scores drift upward until they mean nothing.
  \item \textbf{It cannot be transferred.} Years of standing on one platform are
        worth nothing on another, which is exactly why the platform likes it.
\end{tight}
\end{examplebox}

\section{Portals, Marketplaces and Communities}

The three are often lumped together and should not be.

A \term{portal} aggregates \emph{information} and sells attention --- a price
comparison site, a news aggregator with shopping links. It carries no stock,
handles no payment, and earns from advertising or affiliate referral. Its
engineering problem is ingesting other people's data and keeping it fresh, and
its commercial weakness is that it owns no transaction and can be cut out the
moment a supplier decides to stop feeding it.

A \term{marketplace} aggregates \emph{sellers} and hosts the transaction. It
earns commission, and it is a shop in which the catalogue is written by
strangers --- which turns catalogue quality, duplicate detection and fraud into
daily operational work rather than a one-time setup.

A \term{community} aggregates \emph{people} who are there for each other rather
than to transact. Commerce, if any, is secondary, and the usual mistake is
assuming an engaged community will convert into a profitable marketplace. Many
do not: the thing people valued was the absence of selling.

\begin{conceptbox}[The two-sided launch problem]
A marketplace is worthless to buyers with no sellers, and worthless to sellers
with no buyers. On day one it has neither, and every new one faces this.

\smallskip
The strategies that actually work are all variations on \emph{refusing to be
two-sided at first}:

\begin{tight}
  \item \textbf{Be one of the sellers yourself} until there are enough others.
        You take on stock you did not want, and you compete with your own
        sellers, and it works.
  \item \textbf{Pick one narrow side and own it completely} --- a single city, a
        single category --- so that the market is small enough to fill.
  \item \textbf{Be useful to one side alone.} Offer sellers a tool worth having
        even with no buyers (stock management, a storefront), then introduce
        buyers once the sellers are already there.
  \item \textbf{Fake the thin side} by listing supply gathered from elsewhere.
        Common, legally risky, and the point at which several well-known
        companies acquired their reputations.
\end{tight}

\smallskip
Notice that none of these is ``build a better marketplace''. The problem is not
a product problem, and a student project that proposes a marketplace without
addressing it has not yet met its hardest part.
\end{conceptbox}

\begin{pitfall}
\begin{tight}
  \item \textbf{Implementing bidding with read-then-write.} Two simultaneous
        bids will eventually produce one lost bid, and it will be
        unreproducible when reported.
  \item \textbf{Closing auctions with a page request.} If closing happens when
        somebody loads the page, an auction nobody looks at never closes. It
        needs a scheduled job, and the job needs monitoring.
  \item \textbf{Treating escrow as a feature.} It is custody of other people's
        money and it is regulated. \chref{legal} is not optional reading if you
        intend to hold funds.
  \item \textbf{Believing a reputation score protects the first transaction.}
        It protects the hundredth.
  \item \textbf{Proposing a marketplace without a plan for the empty side.}
        This is the single most common fatal flaw in student project proposals
        for this course.
  \item \textbf{Assuming a community will monetise.} The engagement you measured
        may exist precisely because nothing was being sold.
\end{tight}
\end{pitfall}

\begin{lab}[Three shapes, and what each one must get right]
No software yet. This is the last of the three analysis laboratories;
\chref{stack} starts building.

\smallskip
Find one live example of each: an \textbf{auction}, a \textbf{marketplace} and
a \textbf{portal}. They may be Pakistani or international, but you must be able
to browse all three.

\smallskip
For each, record:

\begin{tightnum}
  \item \textbf{Who takes the risk} if the goods never arrive --- buyer, seller
        or operator? Find the actual policy page, do not guess. Quote it.
  \item \textbf{Where the money is} between payment and delivery. Is there
        escrow? How long is it held, and what releases it?
  \item \textbf{What the operator earns}, and from which side.
  \item \textbf{One rule that exists because of a specific abuse.} Every mature
        platform has several --- a bidding increment, a minimum feedback score,
        a holding period for new sellers, a limit on listing edits. Name one and
        work out what it is defending against.
\end{tightnum}

\smallskip
Then, for the \textbf{auction} only: find the closing rules. Does a late bid
extend the auction? By how much? Say which auction format that makes it, and
whether it rewards bidding your true valuation.

\smallskip
\textbf{Write up half a page} ending with one sentence: of the three, which
would be hardest to build, and which single problem makes it so?

\smallskip
\textbf{What to notice.} The rules that look arbitrary are nearly always scar
tissue. Each one is a record of something that went wrong, and reading a mature
platform's rules is the cheapest way to learn what will go wrong with yours.
\end{lab}

\begin{checkpoint}
\begin{tightnum}
  \item In a Vickrey auction an item is worth Rs.~15{,}000 to you. Bids of
        Rs.~12{,}000 and Rs.~13{,}500 have been placed by others. What should
        you bid, what do you pay if you win, and why would bidding Rs.~18{,}000
        be a mistake?
  \item An auction site closes listings when a visitor loads the listing page.
        Describe precisely what happens to an auction for an unpopular item,
        and what the seller sees.
  \item A marketplace has plenty of buyers and too few sellers. Give two
        strategies from this chapter that would address it, and name a cost of
        each.
\end{tightnum}
\end{checkpoint}

\begin{chaptersummary}
\begin{tight}
  \item Auctions, portals, marketplaces and communities are alternative shapes
        a business can take, which is why this chapter sits in Part~I rather
        than where the published outline puts it.
  \item The four formats are English, Dutch, sealed-bid first-price and
        Vickrey. Under the Vickrey rule, bidding your true valuation is optimal
        whatever others do --- and consumer proxy bidding is Vickrey in
        disguise.
  \item An auction adds four problems a fixed-price shop does not have: time as
        state, concurrent bids, sniping and shill bidding. The first two are
        engineering; the second two are rule design.
  \item ``Read the high bid, check it, write it'' loses bids under concurrency.
        The check and the write must be one atomic operation.
  \item Escrow makes the operator a custodian of other people's money and the
        arbiter of their disputes. Reputation is the cheap substitute and fails
        exactly on the first transaction and the large one.
  \item Portals aggregate information, marketplaces aggregate sellers,
        communities aggregate people. Only the marketplace owns the
        transaction.
  \item Every two-sided market launches empty, and every strategy that works
        amounts to refusing to be two-sided at the start.
\end{tight}
\end{chaptersummary}

\begin{reviewq}
  \item \textbf{(MCQ)} In a second-price sealed-bid auction, the optimal
        strategy is to bid: (a)~slightly above your valuation; (b)~exactly your
        valuation; (c)~slightly below it; (d)~it depends on the other bidders.
  \item \textbf{(MCQ)} Which problem does a fixed-price shop share with an
        auction site? (a)~Sniping; (b)~shill bidding; (c)~concurrent updates to
        one row; (d)~scheduled closing.
  \item \textbf{(Short)} Explain sniping, and give one rule change that
        discourages it together with the cost of that change.
  \item \textbf{(Short)} Distinguish a portal from a marketplace, and say which
        is more easily cut out of its own business.
  \item \textbf{(Short)} Give three distinct ways a reputation system can be
        gamed or rendered meaningless.
  \item \textbf{(Applied)} Design the database handling for a single bid.
        Give the table columns you would need, the exact check that must happen
        atomically, and state what goes wrong if the check and the write are
        separate statements. You may write it as SQL or as precise prose.
  \item \textbf{(Applied)} A student proposes a marketplace connecting
        Bahawalpur tailors with customers. Identify the two-sided launch
        problem in this specific case, choose one strategy from this chapter,
        and describe concretely what the first eight weeks would consist of.
\end{reviewq}

%=====================================================================
\chapter{Planning an E-Commerce Framework}\label{ch:planning}
%=====================================================================
\locator{1}{Part I --- E-Commerce and Its Models}

\begin{outcomes}
By the end of this chapter you will be able to:
\begin{tight}
  \item \textbf{Distinguish} the four ways a shop can be obtained, and
        \textbf{state} what each costs in money, time and control.
  \item \textbf{Write} requirements for a shop in a form that discriminates
        between those four rather than describing all of them equally.
  \item \textbf{Identify} the small number of requirements that actually decide
        the choice, and \textbf{explain} why most requirements do not.
  \item \textbf{Compute} a total cost of ownership over three years, including
        the costs that do not appear on a price list.
  \item \textbf{Defend} a build-buy-or-extend decision against the obvious
        objection to it.
\end{tight}
\end{outcomes}

\begin{prereq}
\chref{models} for the business and revenue model, and its laboratory --- this
chapter turns that one-page brief into a decision. \chref{auctions} if your
venture is a marketplace rather than a shop.
\end{prereq}

\begin{keyterms}
requirement $\bullet$ functional requirement $\bullet$ non-functional
requirement $\bullet$ SaaS $\bullet$ hosted platform $\bullet$ self-hosted
$\bullet$ open source $\bullet$ bespoke development $\bullet$ total cost of
ownership $\bullet$ vendor lock-in $\bullet$ extensibility $\bullet$ data
portability $\bullet$ decision matrix $\bullet$ minimum viable product
\end{keyterms}

\section{The Question This Chapter Answers}

You have a business model, a revenue model and some idea of what you are
selling. There is now exactly one decision to make before any of the rest of
this book applies, and it is this: \emph{where does the software come from?}

There are four answers, and the whole of Parts II to V is a consequence of
which one you pick.

\dsfig{%
\begin{tikzpicture}[
  opt/.style={rounded corners=3pt, draw=#1, line width=1.1pt, fill=#1!8,
              text=#1!60!black, font=\bfseries\scriptsize, align=center,
              text width=3.0cm, minimum height=1.5cm},
  ax/.style={-{Stealth[length=2.4mm]}, neutral, line width=0.9pt},
  tick/.style={font=\scriptsize, text=neutral, align=center}
]
  \node[opt=greenhl]   (a) at (0,0)    {Marketplace\\only\\{\scriptsize\mdseries
       somebody else's shop}};
  \node[opt=secondary] (b) at (3.7,0)  {Hosted\\platform\\{\scriptsize\mdseries
       rented, monthly}};
  \node[opt=primary]   (c) at (7.4,0)  {Self-hosted\\open source\\{\scriptsize
       \mdseries installed and extended}};
  \node[opt=purplehl]  (d) at (11.1,0) {Bespoke\\build\\{\scriptsize\mdseries
       written from nothing}};

  \draw[ax] (-1.8,-1.5) -- (12.9,-1.5);
  \node[tick] at (0,-1.95)    {no control\\no effort};
  \node[tick] at (11.1,-1.95) {total control\\total effort};
  \node[font=\scriptsize\itshape, text=accent] at (5.6,-2.55)
       {every step right costs more of your time and buys more of your fate};

  \node[font=\scriptsize\bfseries, text=primary] at (7.4,1.25)
       {this course};
  \draw[-{Stealth[length=1.8mm]}, primary, line width=0.9pt]
       (7.4,1.05) -- (7.4,0.80);
\end{tikzpicture}}%
{The four ways to obtain a shop. They are a spectrum, not a menu: each step to
the right buys control and costs effort, and the right answer depends on
which of the two you are shorter of.}{buildbuy}

\section{The Four Options, Honestly}

\rowcolors{2}{white}{lightbg}
\begin{longtable}{@{}L{2.8cm}L{5.8cm}L{6.0cm}@{}}
\toprule
\rowcolor{primary!15}
\textbf{Option} & \textbf{What it is good at} & \textbf{What it costs you} \\
\midrule
\endhead
\textbf{Marketplace only} & Customers on day one. No software, no hosting, no
payment integration. Genuinely the right answer for many small sellers & You
perform none of the four functions (\chref{what}). You do not own the customer
relationship, cannot email them, and can be delisted without appeal. The
commission is the visible cost and the smallest one \\
\textbf{Hosted platform} & Running in an afternoon, someone else patches it,
payment and shipping already integrated & A monthly fee that scales with your
success, a transaction cut on top, and a hard ceiling: when the platform
cannot do something, it cannot do it. Getting your data out ranges from
awkward to impossible \\
\textbf{Self-hosted open source} & You can change anything, including the
things the vendor never anticipated. No per-sale cut. The data is in your
database & You are now responsible for hosting, backups, updates and
\emph{security} --- an unpatched shop is a liability, not an inconvenience.
Needs someone who can do \chref{deploy} \\
\textbf{Bespoke build} & Fits the business exactly, because it was built for
it. The only option when the business is genuinely unusual & Twelve to
eighteen months before you sell anything, and you will reimplement baskets,
tax tables and payment retries that have been solved for twenty years ---
badly, because you will meet each edge case for the first time in production \\
\bottomrule
\end{longtable}

\begin{alertbox}[The outline this course follows assumes the fourth option]
The HEC outline for this course is, clause for clause, the contents page of a
2010 book about writing a shopping cart from scratch in PHP~5 --- see the
Course Specification. Taken literally, it asks a 5th-semester class to choose
\emph{bespoke build}.

\smallskip
That was defensible in 2010 and is not now, for a reason this chapter can
state precisely: the work has moved. In 2010, writing a basket and a checkout
was most of the job. Today a mature, free, auditable implementation of both is
one command away, and the job that remains --- modelling the catalogue,
configuring tax and shipping, integrating payment, securing the result,
extending what the platform cannot do --- is both larger and more
representative of what an IT graduate is actually hired to do.

\smallskip
So this course teaches every \emph{topic} the outline lists, and reaches them
by the third route rather than the fourth. You should be able to defend that
choice, not merely inherit it, which is what this chapter is for --- and
\chref{modules} is where the ``extend'' in ``self-hosted and extend'' stops
being a word and becomes code you write.
\end{alertbox}

\section{Requirements That Actually Discriminate}

Students write requirements like ``the system shall allow customers to add
products to a basket''. Every one of the four options does that. A requirement
that all four satisfy carries no information and cannot help you choose.

\begin{definitionbox}[Functional and Non-Functional Requirements]
A \term{functional requirement} says what the system must \emph{do}: accept a
voucher code, calculate tax by province, email an invoice.

\smallskip
A \term{non-functional requirement} says how well, under what constraints: how
fast, how many at once, how available, how secure, how maintainable, by whom.

\smallskip
For the build-buy-extend decision the useful requirements are nearly all
non-functional or unusual-functional. The ordinary functional ones are table
stakes --- everything does them --- and listing them at length is the most
common way a requirements document ends up being thirty pages that decide
nothing.
\end{definitionbox}

\begin{conceptbox}[The five questions that decide it]
In practice the choice turns on a handful of questions. Ask these first, and
most of the rest is detail.

\begin{tightnum}
  \item \textbf{Is there anything genuinely unusual about how you price, sell
        or deliver?} Per-customer negotiated pricing, made-to-measure items,
        rental with a return date, services booked by time slot. If yes, every
        option except the last two gets harder fast.
  \item \textbf{Who will maintain it in two years?} Not who will build it. A
        self-hosted shop with nobody to patch it becomes a breach
        (\chref{security}). If the honest answer is nobody, choose a hosted
        platform and accept the ceiling.
  \item \textbf{What volume, and how uneven?} Thirty orders a day and three
        thousand are different systems; so are steady traffic and traffic that
        multiplies tenfold for one week a year.
  \item \textbf{How much does leaving cost?} Can you export your products,
        customers and order history in a usable form? Ask it on day one,
        because the answer never improves.
  \item \textbf{What is the budget, over three years, including somebody's
        time?} Not the setup cost. See below.
\end{tightnum}
\end{conceptbox}

\section{The Cost That Is Not on the Price List}

The commonest error here is comparing a monthly fee against ``free'' and
concluding that open source is cheaper. Both numbers are wrong.

\begin{worked}[Three years of a small shop, two ways]
A shop expecting \textbf{Rs.~800{,}000 of sales a month}, growing slowly.
Compare a hosted platform against self-hosted open source over three years.

\smallskip
\textbf{Hosted platform.}

\begin{tight}
  \item Subscription: \$79/month $\approx$ Rs.~22{,}000. Over 36 months:
        \textbf{Rs.~792{,}000}.
  \item Platform transaction fee, 0.5\% of sales --- charged on top of the
        gateway's own fee: $0.005 \times 800{,}000 \times 36 =
        \textbf{Rs.~144{,}000}$.
  \item Two paid add-ons at \$15/month each $\approx$ Rs.~8{,}400/month:
        \textbf{Rs.~302{,}400}.
  \item Theme and setup, one-off: \textbf{Rs.~60{,}000}.
  \item Maintenance labour: near zero; the vendor patches.
  \item \textbf{Total: Rs.~1{,}298{,}400.}
\end{tight}

\smallskip
\textbf{Self-hosted open source.}

\begin{tight}
  \item Licence: \textbf{Rs.~0}. This is the number people stop at.
  \item Hosting, a modest virtual server plus backups, Rs.~9{,}000/month:
        \textbf{Rs.~324{,}000}.
  \item Setup and theming, 70 hours of somebody at Rs.~2{,}500/hour:
        \textbf{Rs.~175{,}000}.
  \item Two custom modules the platform does not ship, 40 hours:
        \textbf{Rs.~100{,}000}.
  \item \textbf{Maintenance: 4 hours a month} --- updates, patches, backups
        checked, the occasional breakage. $4 \times 36 \times 2{,}500 =
        \textbf{Rs.~360{,}000}$.
  \item \textbf{Total: Rs.~959{,}000.}
\end{tight}

\smallskip
\textbf{The result.} Self-hosting wins by Rs.~339{,}400 over three years ---
about 26\%. But look at what decided it: \textbf{maintenance labour is
Rs.~360{,}000}, the single largest line in the self-hosted column and the one
most often entered as zero. Set it to zero and self-hosting appears to win by
87\%. Double it, to eight hours a month because nobody on staff really knows
the platform, and the two options are level.

\smallskip
\textbf{The lesson.} This calculation is not really about software prices. It
is a question about whether you have, and will keep, somebody who can maintain
a server --- and if the answer is no, the hosted option is cheaper at any
subscription price, because the alternative is not a cheaper shop but an
unpatched one.
\end{worked}

\begin{alertbox}[Two costs that appear in neither column]
\textbf{Vendor lock-in} is not a line item until the day you want to leave, at
which point it is the whole bill. Before committing, actually export your
catalogue and customers from a trial account and look at what comes out. ``We
support CSV export'' and ``you can reconstruct your shop from this'' are
different claims.

\smallskip
\textbf{The ceiling} has no price at all. It is the feature you cannot have,
discovered in month fourteen, when the business depends on it. The hosted
option's real cost is the probability that this happens multiplied by what it
does to you --- which is why question~1 above is first.
\end{alertbox}

\section{Making the Decision Defensible}

Write it down as a matrix: criteria in rows, weights, options scored. Not
because the arithmetic is wise --- it is easy to fit the weights to the answer
you wanted --- but because it forces the weights into the open where somebody
can disagree with them.

\begin{examplebox}[A decision matrix for the shop in this course]
Criteria weighted out of 100, each option scored 1--5.

\rowcolors{2}{white}{lightbg}
\begin{longtable}{@{}L{4.8cm}L{1.5cm}L{1.9cm}L{1.9cm}L{1.9cm}@{}}
\toprule
\rowcolor{primary!15}
\textbf{Criterion} & \textbf{Weight} & \textbf{Hosted} & \textbf{Self-host} &
\textbf{Bespoke} \\
\midrule
\endhead
Students must see and change real code & 30 & 1 & 5 & 5 \\
Working shop within three weeks & 25 & 5 & 4 & 1 \\
Zero licence or subscription cost & 15 & 2 & 5 & 5 \\
Runs identically on every laptop & 15 & 3 & 5 & 4 \\
Covers the HEC clause list directly & 15 & 3 & 5 & 4 \\
\midrule
\textbf{Weighted total} & \textbf{100} & \textbf{2.75} & \textbf{4.75} &
\textbf{3.70} \\
\bottomrule
\end{longtable}

\smallskip
\textbf{Read the weights, not the totals.} The first row is worth 30 and it is
where hosted loses the contest --- a course whose students cannot read the code
is not this course. Change that weight to 5, as a small shopkeeper would, and
hosted wins outright. The matrix did not make the decision; it made the
decision \emph{visible}, which is all a matrix is for.
\end{examplebox}

\begin{pitfall}
\begin{tight}
  \item \textbf{Requirements that every option satisfies.} ``Customers can add
        to a basket'' discriminates between nothing.
  \item \textbf{Entering maintenance as zero.} It was the largest single line
        in the worked example, and leaving it out reverses the answer.
  \item \textbf{Comparing setup cost instead of three-year cost.} Subscriptions
        and labour both accumulate; neither appears in a setup quote.
  \item \textbf{Choosing bespoke because the platform is ``bloated''.} You will
        reimplement the parts you called bloat, in production, one incident at
        a time.
  \item \textbf{Choosing self-hosted with nobody to maintain it.} This is not a
        cheaper shop. It is a shop that will be compromised, and
        \chref{security} shows how quickly.
  \item \textbf{Not testing the export before committing.} Lock-in costs
        nothing until it costs everything.
  \item \textbf{Weighting the matrix after seeing the totals.} Everyone is
        tempted; write the weights down first and date them.
\end{tight}
\end{pitfall}

\begin{lab}[Your requirements, and your decision]
The last analysis laboratory. \chref{stack} installs software.

\smallskip
Using the venture from \chref{models}'s laboratory, produce a two-page
document:

\begin{tightnum}
  \item \textbf{Twelve requirements}, each labelled functional or
        non-functional, each numbered so you can refer to it later. At least
        four must be non-functional, with a number attached: not ``must be
        fast'' but ``a category page must render in under one second with
        5{,}000 products''. You will test some of these in
        \chref{catalogue} and \chref{deploy}, so make them checkable.
  \item \textbf{Mark the discriminating ones.} Go through your twelve and strike
        out every requirement that all four options satisfy. Most students
        strike out eight or nine. What is left is your real requirements
        document.
  \item \textbf{Answer the five questions} from this chapter, in writing,
        including the honest answer to ``who maintains this in two years?''
  \item \textbf{A three-year total cost of ownership}, two columns, following
        the worked example. Mark every figure you guessed, and state the one
        number the answer is most sensitive to.
  \item \textbf{A decision matrix}, weights written before scores, and a
        two-sentence conclusion.
  \item \textbf{The strongest objection to your decision}, and your answer to
        it. If you cannot state a real objection, you have not understood the
        option you rejected.
\end{tightnum}

\smallskip
\textbf{What to notice.} Your decision does not have to be the one this course
makes. A student whose venture is three products sold through Instagram should
conclude ``marketplace only'' and say so --- and will still build the shop in
Parts II--V, because the course has to teach the general case. Being right
about your own venture and building something else deliberately is a perfectly
coherent position, and a more useful one than pretending.
\end{lab}

\begin{checkpoint}
\begin{tightnum}
  \item Classify each as functional or non-functional, and say which of them
        discriminates between the four options: (a)~customers can reset a
        forgotten password; (b)~prices are agreed per customer and may differ
        for the same product; (c)~the shop stays up during a sale that brings
        twenty times normal traffic.
  \item A shop compares Rs.~22{,}000/month hosted against ``free'' open source
        and concludes open source saves Rs.~792{,}000 over three years. Name
        the three largest costs they have omitted.
  \item Give a specific business for which this course's choice --- self-hosted
        open source --- would be the wrong answer, and say which option you
        would choose instead and why.
\end{tightnum}
\end{checkpoint}

\begin{chaptersummary}
\begin{tight}
  \item A shop can be obtained four ways: marketplace only, hosted platform,
        self-hosted open source, or bespoke build. They form a spectrum trading
        control against effort.
  \item Requirements that every option satisfies cannot help you choose. Strike
        them out; what remains is the document that matters.
  \item Five questions decide most cases: anything unusual about pricing or
        delivery, who maintains it in two years, volume and its unevenness, the
        cost of leaving, and three-year budget including labour.
  \item Compare total cost of ownership over three years, not setup cost. In
        the worked example maintenance labour was the largest self-hosted line,
        and entering it as zero reverses the conclusion.
  \item Vendor lock-in and the feature ceiling appear in no price list and are
        often the real cost of the hosted option.
  \item A decision matrix does not make the decision; it makes the weights
        visible so somebody can argue with them. Write the weights first.
  \item This course chooses self-hosted open source and teaches every clause of
        the outline through it, rather than writing a cart from scratch as the
        2010 book behind the outline assumed. You should be able to defend
        that, including knowing when it is the wrong answer.
\end{tight}
\end{chaptersummary}

\begin{reviewq}
  \item \textbf{(MCQ)} Which is a non-functional requirement? (a)~The shop
        emails an invoice after each order; (b)~the checkout completes in under
        two seconds at 200 concurrent users; (c)~customers can apply a voucher
        code; (d)~administrators can issue refunds.
  \item \textbf{(MCQ)} In the worked example, the largest single cost in the
        self-hosted column was: (a)~hosting; (b)~setup and theming; (c)~custom
        modules; (d)~maintenance labour.
  \item \textbf{(Short)} Explain why a requirement satisfied by all four options
        is useless for this decision, and give an example of one.
  \item \textbf{(Short)} What is vendor lock-in, and what concrete test would
        you run during a free trial to measure it?
  \item \textbf{(Short)} Give the strongest argument \emph{against} the choice
        this course makes, and then answer it.
  \item \textbf{(Applied)} A tailoring business takes measurements in the shop,
        makes garments to order over two weeks, and prices each job
        individually after seeing the cloth. Work through the five questions
        and recommend one of the four options, justifying it. Identify the one
        requirement that makes this case different from a normal shop.
  \item \textbf{(Applied)} Build a three-year total cost of ownership for a
        shop doing Rs.~3{,}000{,}000 a month --- nearly four times the worked
        example --- on both options, keeping all other assumptions. State which
        option wins, and identify the single assumption that, if changed, would
        reverse it.
\end{reviewq}


\part{Standing the Shop Up}
%=====================================================================
% PART II OPENER
%=====================================================================
\thispagestyle{plain}
\structuralfigures{P2.}

\begin{center}
\begin{tcolorbox}[enhanced, width=0.92\textwidth, colback=secondary!5,
  colframe=secondary, boxrule=0pt, leftrule=5pt, arc=3pt, top=8pt, bottom=8pt]
\large\itshape\color{secondary!80!black}
From here on there is a shop. Part~II brings it up, fills it with a catalogue that survives contact with real products, and makes it look like somebody's rather than nobody's --- and it does all of that without installing a single thing on your own machine.
\end{tcolorbox}
\end{center}

\vspace{6pt}

\dsfigh{%
\begin{tikzpicture}[
  cbox/.style={rounded corners=3pt, draw=secondary, line width=1pt, fill=secondary!7,
               text=secondary!80!black, font=\small\bfseries, align=center,
               minimum width=3.5cm, minimum height=1.0cm},
  hbox/.style={rounded corners=3pt, draw=secondary, line width=1.8pt, fill=secondary!16,
               text=secondary!80!black, font=\small\bfseries, align=center,
               minimum width=3.5cm, minimum height=1.0cm},
  arr/.style={-{Stealth[length=2.4mm]}, secondary!60, line width=1pt},
  note/.style={font=\scriptsize\itshape, text=secondary!80!black, align=center,
               fill=white, inner sep=2pt}
]
  \node[cbox] (c0) at (0.00,0) {Ch.\ 5\\The Stack};
  \node[hbox] (c1) at (4.30,0) {Ch.\ 6\\The Catalogue};
  \node[cbox] (c2) at (8.60,0) {Ch.\ 7\\Variations};
  \node[cbox] (c3) at (12.90,0) {Ch.\ 8\\Theme and UX};
  \draw[arr] (c0) -- (c1);
  \draw[arr] (c1) -- (c2);
  \draw[arr] (c2) -- (c3);
  \node[note] at (2.15,-1.0) {it runs};
  \node[note] at (6.45,-1.0) {it has things to sell};
  \node[note] at (10.75,-1.0) {in sizes and colours};
  \node[font=\scriptsize\itshape, text=accent, align=center] at (4.30,1.15)
       {screens on top, tables underneath --- both, or neither is understood};
\end{tikzpicture}}%
{Reading path through Part~II. Chapter 6 is the one to slow down on: it is the only chapter that looks at the same thing from the admin panel and from the database at once.}{part2map}

\newpage

%=====================================================================
\chapter{The Stack, and Your Shop in One Command}\label{ch:stack}
%=====================================================================
\locator{2}{Part II --- Standing the Shop Up}

\begin{outcomes}
By the end of this chapter you will be able to:
\begin{tight}
  \item \textbf{Name} the four layers a web shop runs on and \textbf{state}
        what each one is responsible for.
  \item \textbf{Explain} what a container is, and \textbf{distinguish} it from
        a virtual machine and from an installer.
  \item \textbf{Read} a Compose file and \textbf{say} what each service, port
        and volume in it does.
  \item \textbf{Bring up} a working storefront, \textbf{verify} it from outside,
        and \textbf{reset} it to a known state.
  \item \textbf{Measure} what the stack costs to start and to store, and
        \textbf{distinguish} the cost paid once from the cost paid daily.
\end{tight}
\end{outcomes}

\begin{prereq}
\chref{planning}: the decision to use a self-hosted open-source platform, and
why. Appendix~A must be finished before this chapter --- you need a working
container runtime, and if it is going to fight you, it will fight you now
rather than in Week~9.
\end{prereq}

\begin{keyterms}
stack $\bullet$ LAMP $\bullet$ web server $\bullet$ application server
$\bullet$ database server $\bullet$ container $\bullet$ image $\bullet$ volume
$\bullet$ port mapping $\bullet$ Docker Compose $\bullet$ service $\bullet$
environment variable $\bullet$ reproducibility $\bullet$ idempotent
$\bullet$ pinning $\bullet$ checksum
\end{keyterms}

\section{Four Layers, and What Each One Does}

When you type a shop's address into a browser, four distinct pieces of software
answer in turn. Knowing which is which is the difference between fixing a
problem and restarting things until it goes away.

\dsfig{%
\begin{tikzpicture}[
  lyr/.style={rounded corners=3pt, draw=#1, line width=1.1pt, fill=#1!8,
              text=#1!60!black, font=\bfseries\small, align=center,
              text width=3.1cm, minimum height=1.15cm},
  job/.style={font=\scriptsize, text=neutral, align=left, text width=7.4cm,
              anchor=west},
  arr/.style={-{Stealth[length=2.2mm]}, primary!50, line width=0.9pt}
]
  \node[lyr=neutral]   (os) at (0,0)    {Operating system\\{\scriptsize
       \mdseries Debian Linux}};
  \node[lyr=secondary] (ws) at (0,1.55) {Web server\\{\scriptsize\mdseries
       Apache 2.4}};
  \node[lyr=greenhl]   (ap) at (0,3.10) {Application\\{\scriptsize\mdseries
       PHP 8.2 + OpenCart}};
  \node[lyr=purplehl]  (db) at (0,4.65) {Database\\{\scriptsize\mdseries
       MariaDB 11}};

  \foreach \a/\b in {os/ws, ws/ap, ap/db}{\draw[arr] (\a) -- (\b);}

  \node[job] at (2.0,0)    {Files, processes, the network socket. You will
       almost never touch it.};
  \node[job] at (2.0,1.55) {Accepts the HTTP request, decides which file
       answers it, serves images and stylesheets itself.};
  \node[job] at (2.0,3.10) {Runs the shop's code: works out what the page
       should say, asks the database, builds the HTML.};
  \node[job] at (2.0,4.65) {Remembers everything --- products, prices, baskets,
       orders, customers. The only layer whose loss is unrecoverable.};

  \node[font=\scriptsize\itshape, text=accent, align=center] at (4.8,5.85)
       {A slow page is usually this layer. A lost order is always this one.};
\end{tikzpicture}}%
{The four layers. The traditional name for this arrangement is LAMP --- Linux,
Apache, MySQL, PHP --- and although every letter has alternatives now, the
shape has not changed in twenty-five years.}{lampstack}

\begin{conceptbox}[Which layer is this problem in?]
A diagnostic habit worth building on day one, because it turns ``the site is
broken'' into a question with an answer.

\begin{tight}
  \item \textbf{Nothing answers at all}, the browser cannot connect: the web
        server is down, or the port is wrong. Nothing above it has run.
  \item \textbf{A page appears but says ``500''}: the web server is fine; the
        application failed. The reason is in the PHP error log, never on the
        screen --- and \chref{security} explains why you want it that way.
  \item \textbf{A page appears, looks right, but data is missing or stale}: the
        application ran and the database answered something unexpected.
  \item \textbf{Everything works and is slow}: usually the database, and
        \chref{catalogue} measures exactly that.
\end{tight}
\end{conceptbox}

\section{Containers, and What They Are Not}

The traditional way to get those four layers is to install them: a web server, a
database, a PHP runtime, each with its own installer, its own configuration
file in its own place, and its own way of being slightly different on your
machine from on everybody else's.

A \term{container} replaces that with a package that already contains the
software \emph{and} everything it depends on, down to the system libraries,
arranged exactly as its author arranged it.

\begin{definitionbox}[Image, Container, Volume]
An \term{image} is a read-only, layered filesystem containing an application
and its dependencies. It is built once and is identical everywhere.

\smallskip
A \term{container} is a running instance of an image. It has its own filesystem
view and its own network identity, but it shares the host's kernel --- which is
what makes it start in a second rather than a minute.

\smallskip
A \term{volume} is storage that lives \emph{outside} the container and survives
it. This is the one that matters: everything inside a container is thrown away
when the container is removed, so anything you want to keep --- your database,
your shop --- must be in a volume.
\end{definitionbox}

\begin{alertbox}[Three things a container is not]
\textbf{It is not a virtual machine.} A VM carries a whole operating system
including its own kernel, takes a minute to boot and gigabytes of memory. A
container shares the host kernel and starts in about a second. The isolation is
correspondingly weaker, which matters in production and not in this course.

\smallskip
\textbf{It is not a backup.} ``It is in the container'' means it will be gone
when the container is replaced, which happens every time you rebuild. Your shop
survives because it is in a \emph{volume}, and \chref{deploy} is about what
happens when somebody has not understood that distinction.

\smallskip
\textbf{It is not automatically secure.} A container with a weak password is a
weak password that starts quickly. \chref{security} attacks this very stack and
finds things.
\end{alertbox}

\section{The Three Services in This Course's Stack}

The whole stack is described by one file, \texttt{code/docker-compose.yml},
and it declares three \term{services}.

\begin{lstlisting}[language=yaml]
services:
  db:                              # MariaDB 11: remembers everything
    image: mariadb:11
    volumes:
      - dbdata:/var/lib/mysql      # the shop's memory, kept outside
    healthcheck: ...               # "accepts a login", not "started"

  shop:                            # Apache + PHP 8.2 + OpenCart 4.1.0.4
    build: ./shop
    depends_on:
      db: { condition: service_healthy }
    ports:
      - "8090:80"                  # your 8090 -> the container's 80
    volumes:
      - shopdata:/var/www/html

  adminer:                         # a database console in the browser
    image: adminer:latest
    ports:
      - "8091:8080"
\end{lstlisting}

Three decisions in that file are worth understanding, because each one is a
mistake people make.

\textbf{The healthcheck, not just \texttt{depends\_on}.} Ordering container
\emph{starts} is not enough. MariaDB accepts TCP connections several seconds
before it will accept a login, so a shop that waits only for the database to
have started will try to install itself against a database that refuses it ---
intermittently, and more often on a fast machine. The condition
\texttt{service\_healthy} waits for the database to say it is actually ready.

\textbf{The port mapping is a translation, not a setting.} \texttt{8090:80}
means ``traffic arriving at port 8090 on my machine goes to port 80 inside the
container''. The shop always believes it is on port 80. If 8090 is busy on your
machine --- and on a university laptop it often is --- change the left-hand
number only.

\textbf{Named volumes, so that a reset is deliberate.} \texttt{docker compose
down} stops everything and keeps your shop. \texttt{docker compose down -v}
removes the volumes too, and the next start reinstalls from nothing. That
second command is how you recover from a shop you have broken beyond
understanding, and it is also how you destroy a semester's work in week 14. It
is worth knowing which is which before you need either.

\begin{conceptbox}[Why the image is built here rather than downloaded]
Most teaching material says ``pull the OpenCart image''. There is no longer one
to pull: there has never been an official OpenCart image, and the widely-used
third-party one was withdrawn. Every such instruction on the web is now dead.

\smallskip
So \texttt{code/shop/Dockerfile} builds it, from the official PHP image plus a
\emph{pinned} OpenCart release --- pinned by version and by SHA-256 checksum,
so that the shop you build this year is the shop a classmate builds next year.
Two lines in that file do the real work:

\begin{tight}
  \item \texttt{ARG OC\_VERSION=4.1.0.4} --- the exact release.
  \item \texttt{echo "\$\{OC\_SHA256\}  /tmp/oc.zip | sha256sum -c -} --- the
        download really is the file we tested against. A version number alone
        is not reproducible, because a release asset can be replaced.
\end{tight}

\smallskip
There is a second reason, and it is the pedagogical one: \dsref{lampstack} is a
diagram of something, and pulling one opaque image would have left you unable
to point at any of it. Here the layers are in the file, in order.
\end{conceptbox}

\begin{alertbox}[OpenCart ships its own Dockerfile, and it does not work]
The release archive contains a \texttt{Dockerfile} and a
\texttt{docker-compose.yml} at the top level. This course uses neither, for two
reasons you can verify yourself in five minutes.

\smallskip
The Dockerfile contains a shell syntax error --- \texttt{docker-php-ext-install
zip \&\& \&\& docker-php-ext-enable zip} --- so it has never built. The compose
file pins \texttt{mysql:5.7}, which reached end of life in October~2023 and
receives no further security fixes, and starts Redis, Memcached and PostgreSQL
that a stock install does not use.

\smallskip
The lesson is not that OpenCart is bad software; it is a mature project that
runs a large share of the world's small shops. The lesson is that an artefact
being official is evidence and not authority. You are allowed to check, and in
this course you are expected to.
\end{alertbox}

\section{What Happens the First Time}

One command brings the whole thing up:

\begin{lstlisting}[language=bash]
docker compose up -d --build
\end{lstlisting}

On a first run this does rather a lot: downloads about 1.5~GB of images, builds
the shop image, starts three containers, waits for the database, and then runs
OpenCart's command-line installer, which creates 159 tables and populates them
with a demonstration catalogue. On every later run it does almost none of it.

That difference is worth measuring rather than asserting.

\begin{measured}[What the Stack Costs to Start]
Produced by \texttt{code/m05\_coldstart.py --scratch}, which removes every
image and volume first, on the machine in Appendix~A --- 4 cores, 7.7~GiB
available to the container runtime.

\rowcolors{2}{white}{lightbg}
\begin{longtable}{@{}L{5.8cm}L{2.6cm}L{6.0cm}@{}}
\toprule
\rowcolor{primary!15}
\textbf{Start} & \textbf{To first 200} & \textbf{What it had to do} \\
\midrule
\endhead
From scratch (no images at all) & \textbf{338.4 s} & Download 1.5 GB, compile
three PHP extensions, fetch and verify OpenCart, install 159 tables. Of this,
214.6 s was the build alone \\
Cold (images built, no shop) & \textbf{34.1 s} & Start three containers, wait
for the database, install OpenCart \\
Warm (the shop already exists) & \textbf{16.4 s} & Start three containers.
Nothing else \\
\bottomrule
\end{longtable}

\smallskip
\textbf{Footprint.} Images: shop 882~MB, MariaDB 462~MB, Adminer 164~MB ---
\textbf{1{,}508~MB} in all. Volumes after installation: database 181~MB, shop
files 40~MB. About 1.7~GB of disk for a complete shop.

\smallskip
\textbf{The ratio is the point.} From scratch to warm is \textbf{20.6$\times$}.
You pay the 338 seconds once, in Week~3, and the 16 seconds perhaps sixty times
over the semester. Quoting a single ``startup time'' for this stack would
describe an experience nobody has.

\smallskip
\textbf{On variance.} An earlier run of the same measurement on the same
machine gave 17.6 s cold and 7.7 s warm --- roughly half --- because less was
running at the time. Your numbers will differ from these by more than you
expect. The \emph{ratio} between cold and warm should not.
\end{measured}

\section{Looking At What You Were Given}

The installer leaves you a working shop with a demonstration catalogue. It is
worth knowing its shape, because every chapter in Part~II changes part of it.

\begin{examplebox}[The shop, counted]
Asked of the database directly, through Adminer or the command line:

\begin{lstlisting}[language=sql]
SELECT (SELECT COUNT(*) FROM oc_product)    AS products,
       (SELECT COUNT(*) FROM oc_category)   AS categories,
       (SELECT COUNT(*) FROM oc_order)      AS orders,
       (SELECT COUNT(*) FROM oc_customer)   AS customers;
\end{lstlisting}

\begin{tight}
  \item \textbf{159 tables} in total.
  \item \textbf{19 products} and \textbf{38 categories} of demonstration data.
  \item \textbf{0 orders} and \textbf{0 customers} --- nobody has shopped here
        yet.
  \item \textbf{18 tables whose name begins \texttt{oc\_product}}. One product
        is not one row, and \chref{catalogue} is about why.
\end{tight}

\smallskip
That last number is the one to sit with. A beginner expects a products table.
There are eighteen, because a product has descriptions in several languages,
images, options, attributes, discounts, related items and category memberships
--- and each of those is a different relationship.
\end{examplebox}

\begin{pitfall}
\begin{tight}
  \item \textbf{Running \texttt{down -v} when you meant \texttt{down}.} The
        \texttt{-v} removes the volumes, which is your shop. There is no undo.
  \item \textbf{Editing files inside a running container.} They are gone at the
        next rebuild. Changes that must survive belong in the image or in a
        volume.
  \item \textbf{Assuming \texttt{up} means ready.} It means \emph{started}.
        On a first run the shop answers about twenty seconds later; every
        measurement in this book calls \texttt{waitfor.py} first for exactly
        this reason.
  \item \textbf{Changing the right-hand side of a port mapping.} In
        \texttt{8090:80}, the 80 is the port inside the container and the shop
        expects it. Change the left number.
  \item \textbf{Treating a container as a backup.} It is the opposite of one.
  \item \textbf{Giving up when it does not start.} Appendix~A lists the five
        failures that account for nearly all of them, and four are a busy port
        or a runtime that is not actually running.
\end{tight}
\end{pitfall}

\begin{lab}[Your shop, up, verified, broken and restored]
From here to \chref{project} this is the same shop. Treat it accordingly.

\begin{lstlisting}[language=bash]
# --- 1. Bring the whole stack up -----------------------------------
# First run: expect about five minutes and 1.5 GB of download.
cd code
cp .env.example .env
docker compose up -d --build

# --- 2. Watch it install -------------------------------------------
# Not optional. You should see "SUCCESS! OpenCart successfully
# installed" go past; if you never see it, nothing below will work.
docker compose logs -f shop

# --- 3. Verify from outside, not from the logs ---------------------
# Three URLs must answer 200. waitfor.py checks all three and is the
# answer to "is my shop up?" for the rest of the semester.
python waitfor.py

# --- 4. Look at the shop, as a customer and as the owner -----------
# Storefront   http://localhost:8090/
# Admin        http://localhost:8090/admin/   (admin / admin123)
# Database     http://localhost:8091/         (shop / shoppw / opencart)

# --- 5. Count what you were given ----------------------------------
# In Adminer, or here. 159 tables; 19 products; 0 orders.
docker compose exec db mariadb -ushop -pshoppw opencart \
  -e "SELECT COUNT(*) FROM information_schema.tables \
      WHERE table_schema='opencart';"

# --- 6. Prove the volume is doing what you think -------------------
# Change something in the admin panel first -- edit a product name.
# Then stop WITHOUT -v and start again. Your change is still there.
docker compose down
docker compose up -d
python waitfor.py

# --- 7. Now reset it deliberately ----------------------------------
# -v throws the volumes away. The shop reinstalls from nothing and
# your edit is gone. Do this once, now, while it costs you nothing.
docker compose down -v
docker compose up -d
python waitfor.py

# --- 8. Measure it -------------------------------------------------
# Produces the table quoted above, on YOUR machine.
python m05_coldstart.py --warm
\end{lstlisting}

\smallskip
\textbf{What to notice.} Step~6 and step~7 differ by two characters and by
everything else. Having done both deliberately, in a week when the shop
contains nothing you care about, is the cheapest possible way to learn the
difference.

\smallskip
Also notice that step~3 does not trust step~2. The log said the install
succeeded; step~3 asks the shop. Throughout this course, the thing that
reports success and the thing that is true are kept apart on purpose.
\end{lab}

\begin{checkpoint}
\begin{tightnum}
  \item A classmate's shop returns ``500 Internal Server Error'' on every page.
        Which of the four layers in \dsref{lampstack} are known to be working,
        which is known to have failed, and where would you look first?
  \item Explain why \texttt{depends\_on} alone is not enough to order the
        database before the shop, and describe the failure it produces.
  \item You run \texttt{docker compose down -v} by mistake in Week~12. State
        exactly what is lost, what is not, and what would have had to exist
        beforehand for this to be a minor inconvenience.
\end{tightnum}
\end{checkpoint}

\begin{chaptersummary}
\begin{tight}
  \item A shop runs on four layers --- operating system, web server,
        application, database. Knowing which layer a symptom belongs to is most
        of diagnosis; only the database's loss is unrecoverable.
  \item An image is a built filesystem, a container is a running instance of
        one, and a volume is the only part that survives the container. A
        container is not a virtual machine, not a backup, and not secure by
        itself.
  \item The whole stack is three services in one Compose file: database, shop,
        database console.
  \item \texttt{depends\_on} orders starts, not readiness. A healthcheck is what
        makes the shop wait until the database will actually accept a login.
  \item \texttt{down} keeps your shop; \texttt{down -v} destroys it. Learn the
        difference in Week~3.
  \item The image is built from a pinned release verified by checksum, because
        no maintained OpenCart image exists to pull and because a version
        number alone is not reproducible.
  \item Measured on this machine: 338.4 s from scratch, 34.1 s cold, 16.4 s
        warm --- a ratio of 20.6$\times$ between the cost paid once and the cost
        paid daily. The stack occupies about 1.7 GB.
  \item The installed shop has 159 tables and 19 products, and eighteen of
        those tables are about products. One product is not one row.
\end{tight}
\end{chaptersummary}

\begin{reviewq}
  \item \textbf{(MCQ)} In the port mapping \texttt{8090:80}, the 80 is:
        (a)~the port on your machine; (b)~the port inside the container;
        (c)~the database port; (d)~an arbitrary label.
  \item \textbf{(MCQ)} Which survives \texttt{docker compose down}?
        (a)~Nothing; (b)~the containers; (c)~data written to a named volume;
        (d)~files edited inside a running container.
  \item \textbf{(Short)} Distinguish a container from a virtual machine, and
        give one consequence of the difference that matters in this course.
  \item \textbf{(Short)} Why does this course build the shop image from a pinned
        release rather than pulling a ready-made one? Give both reasons.
  \item \textbf{(Short)} A page loads but shows yesterday's price. Which layer
        is at fault, and what are the two most likely explanations?
  \item \textbf{(Applied)} Your stack will not start: port 8090 is already in
        use. Describe exactly what you would change, in which file, and what you
        would then have to change in the shop's own configuration for links to
        keep working. (Look at \texttt{OC\_HTTP\_SERVER} before answering.)
  \item \textbf{(Applied)} Design a measurement that would show whether the
        \emph{database} or the \emph{application} layer is responsible for a
        slow category page, using only tools in this chapter. State what you
        would record, at what catalogue sizes, and what result would point at
        each layer. \chref{catalogue} runs your experiment; write it first.
\end{reviewq}

%=====================================================================
\chapter{Products, Categories and Their Tables}\label{ch:catalogue}
%=====================================================================
\locator{2}{Part II --- Standing the Shop Up}

\begin{outcomes}
By the end of this chapter you will be able to:
\begin{tight}
  \item \textbf{Explain} why one product occupies eighteen tables, and
        \textbf{say} what decides which table a given fact belongs in.
  \item \textbf{Create} products and categories through the admin panel, and
        \textbf{find} in the database exactly what each screen wrote.
  \item \textbf{Describe} how a category hierarchy is stored, and \textbf{say}
        what the path table buys and what it costs.
  \item \textbf{Read} the query a category page runs, and \textbf{use}
        \texttt{EXPLAIN} to say what the database does with it.
  \item \textbf{Measure} how the catalogue's size affects the page, and
        \textbf{decide}, from evidence, whether an index would help.
\end{tight}
\end{outcomes}

\begin{prereq}
\chref{stack}: a running shop. \emph{Database Systems} for joins, keys and
indexes; Appendix~B if those need refreshing. You will spend as much of this
chapter in the database console as in the admin panel.
\end{prereq}

\begin{keyterms}
catalogue $\bullet$ product $\bullet$ category $\bullet$ SKU $\bullet$ model
$\bullet$ normalisation $\bullet$ junction table $\bullet$ many-to-many
$\bullet$ localisation $\bullet$ materialised path $\bullet$ denormalisation
$\bullet$ index $\bullet$ full table scan $\bullet$ filesort $\bullet$
correlated subquery $\bullet$ pagination
\end{keyterms}

\section{One Product Is Not One Row}

Ask a student to design a products table and you get something like this:

\begin{lstlisting}[language=sql]
CREATE TABLE product (
  product_id  INT PRIMARY KEY,
  name        VARCHAR(255),
  description TEXT,
  price       DECIMAL(10,2),
  image       VARCHAR(255),
  category    VARCHAR(100)
);
\end{lstlisting}

It is a reasonable first attempt and every part of it is wrong in a way that
matters. Your shop has \textbf{eighteen} tables whose names begin
\texttt{oc\_product}, and the reasons for each are worth knowing, because the
same reasons will apply to whatever you build after this course.

\begin{lstlisting}[language=sql]
SELECT table_name FROM information_schema.tables
 WHERE table_schema = 'opencart'
   AND table_name LIKE 'oc_product%'
 ORDER BY table_name;
\end{lstlisting}

\begin{conceptbox}[Four questions that split a table]
Every one of those eighteen tables exists because of one of four questions.
Learn the questions and the table list stops being arbitrary.

\begin{tightnum}
  \item \textbf{``Does this fact depend on the language?''} A price does not;
        a name does. So \texttt{price} lives in \texttt{oc\_product} and
        \texttt{name} lives in \texttt{oc\_product\_description}, keyed on
        \emph{both} the product and the language.
  \item \textbf{``Can there be more than one?''} A product has one price and
        many photographs, so photographs are rows in
        \texttt{oc\_product\_image} rather than columns.
  \item \textbf{``Does this relate two things that both exist already?''} A
        product belongs to several categories and a category holds several
        products. Neither can hold the other, so a junction table ---
        \texttt{oc\_product\_to\_category} --- holds the relationship.
  \item \textbf{``Is this about the product, or about something that happened
        to it?''} \texttt{oc\_product\_viewed} and \texttt{oc\_product\_report}
        record events, not properties, and grow without bound.
\end{tightnum}
\end{conceptbox}

\subsection{The two tables you will use constantly}

\begin{lstlisting}[language=sql]
DESCRIBE oc_product;              -- the language-independent facts
DESCRIBE oc_product_description;  -- the words, per language
\end{lstlisting}

\texttt{oc\_product} holds what is true regardless of who is reading:
\texttt{price}, \texttt{quantity}, \texttt{weight}, \texttt{status},
\texttt{date\_available}, \texttt{sort\_order}. Its primary key is
\texttt{product\_id}.

\texttt{oc\_product\_description} holds \texttt{name},
\texttt{description}, \texttt{meta\_title} and the rest --- and its primary key
is the \emph{pair} \texttt{(product\_id, language\_id)}. One product, one row
per language.

\begin{examplebox}[Why the split is not over-engineering]
It looks like extra work for a shop that will only ever be in English. It is
not, and the reason is worth more than the example.

\smallskip
Suppose name and price shared a row, and the shop later added Urdu. You now
need a second row for the Urdu name --- which duplicates the price. Change the
price and you must remember to change it in both rows, and the day somebody
forgets, the shop charges two different prices depending on the language the
customer is reading. That is not a bug that gets noticed quickly; it is a bug
that gets noticed by an accountant, months later.

\smallskip
The split makes that mistake \emph{impossible to express}. There is exactly one
place the price lives, so there is exactly one price. This is what
normalisation is for, and it is the clearest small example of it you will meet.
\end{examplebox}

\section{The Screen, and What It Writes}

Create a product: \adminpath{Catalog, Products, Add New}. The form has tabs ---
\uiel{General}, \uiel{Data}, \uiel{Links}, \uiel{SEO} --- and the tabs are
not a presentational choice. They are very nearly the table list.

\rowcolors{2}{white}{lightbg}
\begin{longtable}{@{}L{2.6cm}L{5.4cm}L{6.8cm}@{}}
\toprule
\rowcolor{primary!15}
\textbf{Tab} & \textbf{What you type} & \textbf{Where it lands} \\
\midrule
\endhead
\uiel{General} & Product name, description, meta title and tags & One row in
\texttt{oc\_product\_description}, for the language you were editing \\
\uiel{Data} & Model, price, quantity, weight, dimensions, date available,
status & The single row in \texttt{oc\_product} \\
\uiel{Links} & Manufacturer, categories, stores, related products & Rows in
\texttt{oc\_product\_to\_category}, \texttt{oc\_product\_to\_store} and
\texttt{oc\_product\_related} --- one per tick \\
\uiel{Image} & Main image and additional images & \texttt{image} on the product
row; the extras become rows in \texttt{oc\_product\_image} \\
\uiel{SEO} & The URL keyword & \texttt{oc\_seo\_url} --- not a product table at
all, which \chref{seo} explains \\
\bottomrule
\end{longtable}

\begin{alertbox}[Do this once, deliberately, and you will never be confused again]
Create one product. Then, in the database console at
\texttt{http://localhost:8091/}, run:

\begin{lstlisting}[language=sql]
SELECT MAX(product_id) FROM oc_product;   -- yours, the newest
\end{lstlisting}

and look at that id in \texttt{oc\_product}, then in
\texttt{oc\_product\_description}, then in
\texttt{oc\_product\_to\_category}. Three tables, one product, and you typed it
all on one screen.

\smallskip
Every chapter after this one assumes you can do this. When something in the
shop is not what you expected, the question is always ``what did it actually
write?'', and the admin panel is the wrong place to look for the answer.
\end{alertbox}

\section{Categories, and the Table That Is Not Normalised}

Categories form a tree: \adminpath{Components, Monitors} sits under
\adminpath{Components}, which sits at the top. The obvious storage is a
\texttt{parent\_id} on each category, and that is indeed what
\texttt{oc\_category} has.

It is not enough, and the reason is the one interesting thing about category
storage.

\begin{alertbox}[The question a parent pointer cannot answer cheaply]
A customer opens \adminpath{Components}. The page must show every product in
\emph{Components or anything beneath it}, however deep.

\smallskip
With only \texttt{parent\_id}, answering that means: find the children of
Components; find their children; find theirs; stop when there are no more.
That is a loop with a query inside it, and the number of queries depends on how
deep the tree happens to be --- so the page gets slower as the shop's
categories get better organised, which is precisely the wrong incentive.
\end{alertbox}

So OpenCart keeps a second table that stores the answer directly.

\begin{lstlisting}[language=sql]
DESCRIBE oc_category_path;
-- category_id, path_id, level

SELECT * FROM oc_category_path WHERE category_id = 27 ORDER BY level;
-- category_id  path_id  level
--          27       20      0      <- Components, the ancestor
--          27       27      1      <- Monitors, itself
\end{lstlisting}

\begin{definitionbox}[Materialised Path]
A \term{materialised path} stores, for every node, the complete list of its
ancestors --- one row per ancestor, with the depth. The tree is still in
\texttt{parent\_id}; this table is a \emph{precomputed answer} to ``who is
above me?'', and therefore also to ``who is below me?''

\smallskip
It is a deliberate \term{denormalisation}: the same fact is now stored twice,
in \texttt{oc\_category.parent\_id} and in \texttt{oc\_category\_path}. That
buys a single-query answer to the descendants question, and costs the
obligation to keep the two in step --- so moving one category in the admin
panel rewrites every path row beneath it.

\smallskip
This is the trade-off in miniature: \textbf{denormalisation buys read speed
with write complexity and the risk of disagreement.} Make it deliberately, and
only where the reads vastly outnumber the writes --- which for a category tree
they do, by perhaps a million to one.
\end{definitionbox}

\section{What the Category Page Actually Asks}

Here is where this chapter stops being about design and starts being about
evidence. The shop's own source is in your container, and you can read it:

\begin{lstlisting}[language=bash]
docker compose exec shop \
  cat /var/www/html/catalog/model/catalog/product.php
\end{lstlisting}

The listing query is built up in \texttt{getProducts()}, and reduced to its
skeleton it looks like this:

\begin{lstlisting}[language=sql]
SELECT DISTINCT *, pd.name, p.image,
       (SELECT ... FROM oc_product_discount pd2
         WHERE pd2.product_id = p.product_id ... LIMIT 1) AS discount,
       (SELECT ... FROM oc_product_discount ps
         WHERE ps.product_id = p.product_id ... LIMIT 1) AS special,
       (SELECT pr.points FROM oc_product_reward pr
         WHERE pr.product_id = p.product_id ...)         AS reward,
       (SELECT COUNT(*) FROM oc_review r
         WHERE r.product_id = p.product_id ...)          AS reviews
  FROM oc_category_to_store c2s
  LEFT JOIN oc_category_path cp  ON (cp.category_id = c2s.category_id)
  LEFT JOIN oc_product_to_category p2c ON (p2c.category_id = cp.category_id)
 INNER JOIN oc_product_to_store p2s ON (p2s.product_id = p2c.product_id)
  LEFT JOIN oc_product p  ON (p.product_id = p2s.product_id
                              AND p.status = 1
                              AND p.date_available <= NOW())
  LEFT JOIN oc_product_description pd ON (p.product_id = pd.product_id)
 WHERE ...
 ORDER BY p.sort_order
 LIMIT 0, 20;
\end{lstlisting}

Four things in that query are worth naming, because each is a decision somebody
made and each has a consequence you can measure.

\begin{tightnum}
  \item \textbf{\texttt{SELECT DISTINCT *}} --- every column of every joined
        table, then deduplicated. The \texttt{DISTINCT} is there because the
        path join can produce a product twice; it forces a temporary table.
  \item \textbf{Four correlated subqueries.} Discount, special, reward and
        review count are each evaluated \emph{per row returned}. With
        \texttt{LIMIT 20} that is eighty extra queries' worth of work per page,
        and it is the SQL form of the N+1 problem.
  \item \textbf{\texttt{ORDER BY p.sort\_order}} is the default sort, and
        \texttt{oc\_product} has \textbf{no index on \texttt{sort\_order}} ---
        or on \texttt{price}, \texttt{status} or \texttt{date\_available}. It
        has a primary key and nothing else.
  \item \textbf{The whole join runs twice.} \texttt{getTotalProducts()} repeats
        it with \texttt{COUNT(DISTINCT p.product\_id)}, because the page needs
        to know how many pages there are. Every category page you have ever
        loaded ran that join twice.
\end{tightnum}

\begin{examplebox}[Reading EXPLAIN on the real query]
Put \texttt{EXPLAIN} in front of it and the database says what it intends:

\smallskip
\texttt{p:ALL [Using where; Using temporary; Using filesort]}

\smallskip
Three findings in one line. \textbf{\texttt{ALL}} is a full table scan of
\texttt{oc\_product} --- every row read, at any catalogue size.
\textbf{\texttt{Using temporary}} is the \texttt{DISTINCT} building a temporary
table. \textbf{\texttt{Using filesort}} is the \texttt{ORDER BY} sorting the
result afterwards, because no index offers it in that order.

\smallskip
Now ask for the same page sorted by name, which is \texttt{ORDER BY
LCASE(pd.name)}:

\smallskip
\texttt{pd:index [Using where; Using index; Using temporary; Using filesort]}

\smallskip
Still a filesort --- and \texttt{oc\_product\_description} \emph{does} have an
index on \texttt{name}. \chref{catalogue}'s key lesson is in that
contradiction, and the next section measures it.
\end{examplebox}

\begin{conceptbox}[A function around a column hides the column]
\texttt{ORDER BY LCASE(pd.name)} cannot use an index on \texttt{name}.

\smallskip
An index on \texttt{name} stores the \emph{names}, in order. The query does not
ask for names in order --- it asks for \texttt{LCASE(name)} in order, and the
database has no index on that. It cannot assume the two orderings agree,
because for some collations they do not. So it reads every row, computes
\texttt{LCASE} on each, and sorts the results.

\smallskip
\textbf{The rule, which is worth more than this example:} the moment you wrap a
column in a function in \texttt{WHERE} or \texttt{ORDER BY}, any index on that
column stops being usable. \texttt{WHERE YEAR(date\_added) = 2026} and
\texttt{WHERE UPPER(email) = ?} are the same mistake, and they are extremely
common.

\smallskip
The fixes are to store the value you actually search by, to use a
case-insensitive collation so \texttt{LCASE} is unnecessary, or --- in modern
MariaDB and MySQL --- to build an index on the expression itself.
\end{conceptbox}

\section{What It Costs, and What an Index Buys}

The argument so far is that this query ought to be slow and that an index ought
to help. Both halves are worth checking, and only one of them survives.

\begin{measured}[The Category Page Against Catalogue Size]
Produced by \texttt{code/m06\_catalogue.py} on the machine in Appendix~A.
Products are seeded into one category; the page is requested seven times at
each size and the median taken; then the same page is requested once more with
MariaDB's query log recording \emph{every} statement it ran and how long each
took.

The whole run was done \textbf{twice}, because one of its results was
surprising enough that a single observation would not have been worth
reporting. Both runs are shown.

\rowcolors{2}{white}{lightbg}
\begin{longtable}{@{}L{2.0cm}L{2.5cm}L{2.2cm}L{2.6cm}L{2.6cm}@{}}
\toprule
\rowcolor{primary!15}
\textbf{Products} & \textbf{Page} & \textbf{Queries} & \textbf{DB time} &
\textbf{DB as \% of page} \\
\midrule
\endhead
100 & 112 / 111\,ms & \textbf{144} & 60 / 35\,ms & 54\% / 32\% \\
1{,}000 & 203 / 101\,ms & \textbf{144} & 60 / 60\,ms & 30\% / 60\% \\
10{,}000 & 171 / 302\,ms & \textbf{144} & 114 / 128\,ms & 67\% / 43\% \\
50{,}000 & 541 / 476\,ms & \textbf{146 / 144} & 402 / 477\,ms & 74\% / 100\% \\
\bottomrule
\end{longtable}

\smallskip
\textbf{One page, one hundred and forty-four queries} --- and that is the most
stable number in the whole experiment. It is the same at a hundred products as
at fifty thousand, in both runs. Twenty products are displayed, each needing
its images, its discount, its special price and its review count, and around
forty more queries fetch settings, languages, currencies, layout modules and
the category tree. The $N+1$ problem here is per \emph{displayed} product,
which is why growing the catalogue does not grow it.

\smallskip
\textbf{The catalogue's size does cost, and it lands on the database.} Five
hundred times the products cost roughly 4$\times$ the page time and 8--13$\times$
the database time, and the database's share rises from about a third to
between 74\% and 100\%. At fifty thousand products this page is
database-bound; at a hundred it is not.

\smallskip
\textbf{Now the surprise. The indexes help in the middle and hurt at the top.}

\rowcolors{2}{white}{lightbg}
\begin{longtable}{@{}L{2.2cm}L{3.0cm}L{3.0cm}L{4.0cm}@{}}
\toprule
\rowcolor{primary!15}
\textbf{Products} & \textbf{DB before} & \textbf{DB after} &
\textbf{Change (run 1 / run 2)} \\
\midrule
\endhead
100 & 60 / 35\,ms & 43 / 42\,ms & $-28\%$ / $+19\%$ --- noise \\
1{,}000 & 60 / 60\,ms & 41 / 48\,ms & $\mathbf{-31\%}$ / $\mathbf{-20\%}$ \\
10{,}000 & 114 / 128\,ms & 107 / 106\,ms & $-6\%$ / $\mathbf{-18\%}$ \\
50{,}000 & 402 / 477\,ms & 616 / 574\,ms & $\mathbf{+53\%}$ / $\mathbf{+20\%}$ \\
\bottomrule
\end{longtable}

\smallskip
At a thousand and ten thousand products the indexes make the database faster,
in both runs. At fifty thousand they make it \textbf{slower}, in both runs. The
direction reproduces even though the magnitudes do not, which is what makes it
worth believing.

\smallskip
\textbf{Why an index can lose.} \texttt{EXPLAIN} shows what changed:
\texttt{p:ALL [\ldots\ Using filesort]} became \texttt{p:index [\ldots]}. Given
an index on \texttt{sort\_order}, the optimiser walks the index in order and
avoids the sort --- but it must then fetch each row by primary key, which is
random access. Without the index it reads the table straight through, which is
sequential, and sorts afterwards. The \texttt{DISTINCT} means it has to touch
nearly every row either way, so the question is only whether those reads are
sequential or scattered. At fifty thousand rows, \textbf{scattered reads plus
no sort lose to sequential reads plus a sort.}

\smallskip
\textbf{The chapter's conclusion, which is not the one the chapter expected.}
``Add an index'' is not a diagnosis, and removing a filesort is not
automatically a win. The index is right for a mid-sized catalogue and wrong for
a large one, and nothing in \texttt{EXPLAIN} told us that --- only the clock
did. Meanwhile the page still runs 144 queries at every size, which no index
touches at all. That is a structural problem, and \chref{deploy} meets it with
the only tool that works at that level: caching.
\end{measured}

\begin{alertbox}[Three earlier versions of this measurement were wrong]
All three are recorded because all three mistakes are ones you will make. The
third is the worst kind.

\smallskip
\textbf{Third, and the one that should frighten you: the products were not
there.} The seeder wrote \texttt{oc\_product\_to\_store} with
\texttt{store\_id = 1}. The default store is \textbf{0}. Every storefront query
filters on \texttt{p2s.store\_id = config\_store\_id}, so fifty thousand
products sat in the database, appeared in the admin panel, and were invisible
to every customer --- and to the measurement, which was timing a category page
showing the nineteen demonstration products.

\smallskip
Nothing failed. No error, no warning, no empty page: the page rendered
perfectly, returned 200, and meant something else. The numbers it produced were
real numbers about a different question, and they were only caught because a
\emph{later} chapter tried to open one of the seeded products and got a 404.

\smallskip
That is the characteristic shape of a measurement bug. It does not crash; it
quietly answers a question you did not ask. The defence is to check that your
setup did what you think before you trust what it tells you --- here, one
\texttt{curl} for a seeded product's page would have caught it in week one.

\smallskip
\textbf{First: the harness timed itself.} It ran the listing query
twenty-five times through \texttt{docker compose exec} and divided the wall
clock by twenty-five. It reported \textbf{707\,ms of query time inside a
277\,ms page} --- impossible, and the impossibility was the only reason it was
questioned. What it was really measuring was the cost of spawning a process and
opening a connection, which dwarfs the query. Timing a trivial statement first
and subtracting it fixed the small sizes.

\smallskip
\textbf{Second, and worse: it was timing the wrong query.} The harness used a
hand-written query that \emph{resembled} the one OpenCart runs. At 50{,}000
products it still reported more database time than the whole page took, because
the hand-copied version lacked the real one's restrictions and did far more
work. Two queries that look alike are not alike.

\smallskip
\textbf{The fix was to stop guessing.} MariaDB will record every statement it
executes, with its duration, if you set \texttt{long\_query\_time=0} and turn
the slow query log on. That is what produced the table above --- the real
queries, timed by the server, including the 146 of them.

\smallskip
Two rules, both expensive to learn any other way. \textbf{If a measurement
gives you a number that cannot be true, believe the impossibility.} And
\textbf{never hand-copy the query you mean to measure} --- ask the server what
it ran.
\end{alertbox}

\begin{pitfall}
\begin{tight}
  \item \textbf{Putting the name in \texttt{oc\_product}.} It is the one column
        that depends on the language, and the one students most want to move.
  \item \textbf{Editing \texttt{oc\_category\_path} by hand.} It is derived.
        Move categories in the admin panel and let it rewrite the paths, or the
        tree and the paths disagree and the storefront shows products in
        categories they are not in.
  \item \textbf{Wrapping a column in a function and expecting the index to
        work.} \texttt{LCASE(name)}, \texttt{YEAR(date\_added)},
        \texttt{UPPER(email)} --- all the same mistake.
  \item \textbf{Adding indexes because a page is slow.} Measure first. This
        chapter added three, removed a full scan and a filesort, and gained
        nothing.
  \item \textbf{Forgetting that the join runs twice.} Optimising the listing
        query and ignoring the \texttt{COUNT} halves your improvement at best.
  \item \textbf{Testing a catalogue with nineteen products.} Every interesting
        behaviour in this chapter is invisible below about a thousand.
\end{tight}
\end{pitfall}

\begin{lab}[Fifty thousand products, and where the time goes]
Your shop has nineteen products. By the end of this laboratory it will have had
fifty thousand, and you will know what that costs.

\begin{lstlisting}[language=bash]
# --- 1. Look at what one product is, before changing anything -------
# Create a product in the admin panel: Catalog > Products > Add New.
# Fill in General (name), Data (model, price) and Links (a category).
# Then find it three times, in three tables.
cd code
docker compose exec db mariadb -ushop -pshoppw opencart

# --- 2. Count the tables one product touches -----------------------
# 18 of them begin oc_product. Read the list and, for each, say which
# of the four questions in this chapter put it there.

# --- 3. Fill the catalogue ------------------------------------------
# Seeded from a fixed seed, so your classmate's shop matches yours.
python seed/catalogue.py 50000

# --- 4. Load the category page and feel it --------------------------
# http://localhost:8090/index.php?route=product/category&path=20
# Then sort by name, and by price. One of those is slower; predict
# which before you try it.

# --- 5. Ask the database what it is doing ---------------------------
# EXPLAIN the listing query. Find ALL, Using temporary, Using filesort.
# Then EXPLAIN it with ORDER BY LCASE(pd.name) and explain why the
# existing index on name does not save you.

# --- 6. Run the measurement -----------------------------------------
# Four catalogue sizes, with and without the missing indexes. Watch
# the query COUNT as the catalogue grows -- it is the finding.
python m06_catalogue.py

# --- 7. Disagree with the book --------------------------------------
# The book found 146 queries per page, and that the indexes halved
# the database time without reliably moving the page time. Does
# yours? Then answer the question the book leaves open: of those 146
# queries, how many are for the 20 products shown, and how many would
# still run on an EMPTY category? Find out by requesting a category
# with no products in it, with the query log on.

# --- 8. Put your shop back ------------------------------------------
python seed/catalogue.py --restore
\end{lstlisting}

\smallskip
\textbf{What to notice.} Step 7 is the real exercise. The measurement in this
book is a claim about one machine on one afternoon, and your job is not to
reproduce the number but to find out whether the \emph{shape} holds on yours
--- and, if it does, to say where the time is actually going. A student who can
answer that can optimise a shop; one who can only add indexes cannot.
\end{lab}

\begin{checkpoint}
\begin{tightnum}
  \item A shop adds Urdu. Which of these move to
        \texttt{oc\_product\_description} and which stay in
        \texttt{oc\_product}, and why: price, product name, weight, meta
        description, stock quantity?
  \item \texttt{oc\_category\_path} stores facts already derivable from
        \texttt{oc\_category.parent\_id}. Name what this buys, what it costs,
        and the one operation that makes the cost visible.
  \item A colleague reports that a category page is slow, runs
        \texttt{EXPLAIN}, sees \texttt{Using filesort}, adds an index, and the
        page is no faster --- although the database's own time halved. Give two
        distinct explanations for why the page did not improve, and say what you
        would measure to tell them apart.
\end{tightnum}
\end{checkpoint}

\begin{chaptersummary}
\begin{tight}
  \item One product occupies eighteen tables. Four questions decide which:
        does the fact depend on language, can there be more than one, does it
        relate two existing things, and is it a property or an event?
  \item \texttt{oc\_product} holds the language-independent facts;
        \texttt{oc\_product\_description} holds the words, keyed on
        \texttt{(product\_id, language\_id)}. The split makes a price that
        differs by language impossible to express.
  \item The admin form's tabs are very nearly the table list. Create one
        product and find it in three tables; every later chapter assumes you
        can.
  \item A parent pointer cannot answer ``everything beneath this category''
        in one query, so \texttt{oc\_category\_path} stores the answer --- a
        deliberate denormalisation buying read speed with write complexity.
  \item The category listing query does \texttt{SELECT DISTINCT *} with four
        correlated subqueries, sorts on an unindexed column, and the whole join
        runs twice because the page needs a count.
  \item \texttt{ORDER BY LCASE(name)} cannot use the index on \texttt{name}.
        A function around a column hides the column from the optimiser.
  \item Measured: \textbf{one category page runs 146 queries}, and that number
        barely moves between a hundred products and fifty thousand. The $N+1$
        is per displayed product, not per catalogue product.
  \item Five hundred times the catalogue cost about 2.5$\times$ the page time,
        because \texttt{LIMIT 20} bounds the work.
  \item The three missing indexes cut database time by 45--62\% at every size,
        and the page time did not reliably follow --- the gain is real and
        smaller than the run-to-run noise. Measure the database to see it.
  \item An index makes one query faster; it cannot make 145 others disappear.
        The page's real problem is structural, and \chref{deploy} meets it with
        caching.
\end{tight}
\end{chaptersummary}

\begin{reviewq}
  \item \textbf{(MCQ)} The primary key of \texttt{oc\_product\_description} is:
        (a)~\texttt{product\_id}; (b)~\texttt{(product\_id, language\_id)};
        (c)~\texttt{name}; (d)~an auto-increment id.
  \item \textbf{(MCQ)} \texttt{Using filesort} in \texttt{EXPLAIN} means:
        (a)~the query wrote a temporary file to disk; (b)~the result had to be
        sorted after fetching, because no index gave that order; (c)~the table
        has no primary key; (d)~the query failed.
  \item \textbf{(Short)} Give the four questions that decide which table a
        product fact belongs in, with one example each.
  \item \textbf{(Short)} Explain why \texttt{WHERE YEAR(date\_added) = 2026}
        cannot use an index on \texttt{date\_added}, and rewrite it so it can.
  \item \textbf{(Short)} What is a materialised path, and what has to happen
        when a category is moved in the admin panel?
  \item \textbf{(Applied)} A shop wants products to carry a ``minimum order
        quantity'' that differs per customer group. Say which existing table
        this belongs in or what new table is needed, give the columns and the
        primary key, and justify it against the four questions.
  \item \textbf{(Applied)} The measurement in this chapter found that removing
        a full table scan and a filesort did not speed the page up. Design an
        experiment that would locate where the time actually goes, using only
        what is in your container. State what you would vary, what you would
        hold fixed, what you would record, and what result would point at each
        of the four layers in \dsref{lampstack}.
\end{reviewq}

%=====================================================================
\chapter{Product Variations, Options and Uploads}\label{ch:variations}
%=====================================================================
\locator{2}{Part II --- Standing the Shop Up}

\begin{outcomes}
By the end of this chapter you will be able to:
\begin{tight}
  \item \textbf{Distinguish} an option from a variant, and \textbf{say} which
        of the two a given requirement needs.
  \item \textbf{Explain} why options do not store the combinations they imply,
        and \textbf{state} what that makes impossible.
  \item \textbf{Configure} option sets on a product and \textbf{predict} the
        number of rows each one creates.
  \item \textbf{Describe} what a shop must check before accepting an uploaded
        file, and \textbf{identify} which of those checks a client can defeat.
  \item \textbf{Test} your own shop's upload handling and \textbf{report} what
        it accepts.
\end{tight}
\end{outcomes}

\begin{prereq}
\chref{catalogue}: the eighteen product tables, and the habit of looking at
what a screen wrote. Your shop should still have its demonstration catalogue;
\texttt{seed/catalogue.py --restore} puts it back if not.
\end{prereq}

\begin{keyterms}
option $\bullet$ option value $\bullet$ variant $\bullet$ SKU $\bullet$
modifier $\bullet$ price prefix $\bullet$ combinatorial explosion $\bullet$
stock keeping $\bullet$ required option $\bullet$ file upload $\bullet$
allow-list $\bullet$ MIME type $\bullet$ content sniffing $\bullet$ stored
cross-site scripting
\end{keyterms}

\section{Two Ways to Say ``Comes in Three Colours''}

A shirt in three colours and four sizes. Twelve things a customer might buy.
How many rows is that?

The honest answer is that it depends on a decision nobody will prompt you to
make, and that the two answers have different consequences for the rest of the
shop's life.

\begin{definitionbox}[Option and Variant]
An \term{option} is a choice attached to one product, whose values
\emph{modify} it. The shirt is one product; ``Large'' adds Rs.~200 and
``Red'' adds nothing. The twelve combinations are \textbf{not stored
anywhere} --- they are computed when a customer picks.

\smallskip
A \term{variant} is a product in its own right, linked to a master. ``Red
Large Shirt'' is a row in \texttt{oc\_product}, with
\texttt{master\_id} pointing at the parent, its own model number, its own
stock and its own price. Twelve combinations are twelve rows.

\smallskip
OpenCart~4 supports both --- \texttt{oc\_product} carries
\texttt{master\_id}, \texttt{variant} and \texttt{override} columns alongside
the older option tables. The admin panel shows options first, so options are
what most shops end up with, including shops that needed variants.
\end{definitionbox}

\subsection{What options actually store}

Two tables do the work, and a third names the choices globally.

\rowcolors{2}{white}{lightbg}
\begin{longtable}{@{}L{5.4cm}L{9.6cm}@{}}
\toprule
\rowcolor{primary!15}
\textbf{Table} & \textbf{What a row means} \\
\midrule
\endhead
\texttt{oc\_option} & A choice that exists in the shop at all --- ``Size'',
``Colour'' --- with its \texttt{type}: \texttt{select}, \texttt{radio},
\texttt{checkbox}, \texttt{text}, \texttt{textarea}, \texttt{date},
\texttt{time}, \texttt{datetime}, \texttt{file} \\
\texttt{oc\_option\_value} & One possible value of it: ``Large''. Shared
across every product that uses the option \\
\texttt{oc\_product\_option} & \emph{This} product offers \emph{that} option,
and whether it is \texttt{required} \\
\texttt{oc\_product\_option\_value} & This product offers that value, and
what choosing it does: \texttt{price} with a \texttt{price\_prefix} of
\texttt{+} or \texttt{-}, plus \texttt{quantity}, \texttt{weight} and
\texttt{points} adjustments \\
\bottomrule
\end{longtable}

\begin{alertbox}[What options make impossible]
Look at where \texttt{quantity} lives: on
\texttt{oc\_product\_option\_value}, which is \emph{per option value}, not per
combination.

\smallskip
So the shop can know it has 40 Large shirts and 30 Red shirts. It
\textbf{cannot} know how many Red Large shirts it has, because there is no row
for that --- and no amount of configuration will add one. Three consequences
follow, and all three are discovered late:

\begin{tight}
  \item \textbf{No stock per combination.} You will oversell the popular
        combination and sit on the unpopular one.
  \item \textbf{No SKU or barcode per combination}, so the shop cannot speak
        to a warehouse system that works in SKUs --- which is all of them.
  \item \textbf{No price history or reporting per combination.} ``How many
        Red Large did we sell?'' has no answer.
\end{tight}

\smallskip
If any of those three matter, you need variants, and you need to know that
\emph{before} the catalogue is built. Migrating two thousand option-based
products to variants is not a setting; it is a project.
\end{alertbox}

\begin{conceptbox}[Which one does this requirement need?]
A test that settles most cases in one question: \textbf{does the shop need to
count, price or identify the combination on its own?}

\smallskip
\textbf{Options are right} when the choice does not change what you hold in
stock: an engraving message, a delivery date, a gift-wrap tick, a cable length
that you cut to order. Also when the combination genuinely shares stock ---
which is rarer than people assume.

\smallskip
\textbf{Variants are right} whenever the combination is a thing in a box on a
shelf. Clothing, shoes, anything with a barcode. If a warehouse person could
pick up exactly one combination and scan it, it is a variant.

\smallskip
\textbf{The common mistake} is to choose options because the admin panel
offers them first and they are quicker to set up, and to discover in month
nine --- during the first real stock-take --- that the shop has never known
what it holds.
\end{conceptbox}

\section{What Each Choice Costs}

Options are cheap to store and variants are not. That much is obvious. What is
worth measuring is \emph{how fast} the difference grows, because it decides
whether a catalogue is manageable.

\begin{measured}[Options Against Variants]
Produced by \texttt{code/m07\_variations.py} on the machine in Appendix~A.
Option sets are attached to one product and the product page is requested seven
times, median taken; the variant column is the number of
\texttt{oc\_product} rows the same catalogue would need.

\rowcolors{2}{white}{lightbg}
% Six columns inside a tcolorbox: 10.4 cm of content plus five 0.42 cm gaps.
% Seven columns overflowed by 12.3 pt, and "option rows" was the column the
% prose could carry instead.
\begin{longtable}{@{}L{1.9cm}L{1.7cm}L{1.7cm}L{1.6cm}L{1.4cm}L{2.1cm}@{}}
\toprule
\rowcolor{primary!15}
\textbf{Options} & \textbf{Combin.} & \textbf{Value rows} & \textbf{Page} &
\textbf{KB} & \textbf{As variants} \\
\midrule
\endhead
3 & 3 & 3 & 87\,ms & 30.0 & 4 rows \\
3$\times$4 & 12 & 7 & 104\,ms & 30.7 & 13 rows \\
3$\times$4$\times$5 & 60 & 12 & 101\,ms & 31.3 & 61 rows \\
6$\times$8$\times$5 & \textbf{240} & \textbf{19} & 103\,ms & 31.3 &
\textbf{241 rows} \\
\bottomrule
\end{longtable}

\smallskip
\textbf{Options grow by addition; variants grow by multiplication.} Three
option sets of 6, 8 and 5 values are \textbf{19 rows} --- $6+8+5$ --- and
express 240 combinations. The same catalogue as variants is \textbf{241
product rows}, each with its own description row, its own store row and its
own category rows. Add a sixth colour and options gain one row while variants
gain forty.

\smallskip
\textbf{And the product page does not care.} From 3 combinations to 240, the
page stayed at about 100\,ms and grew by \textbf{1.3 kilobytes} --- because the
page renders the \emph{choices}, not the combinations. Nineteen values in three
dropdowns, whatever the 240 products those values imply.

\smallskip
\textbf{So the cost of variants is not page speed; it is catalogue
management.} Two hundred and forty-one rows that must each be priced, stocked,
described, photographed and kept in step. That is the real bill, it is paid by
a person rather than a server, and it is the reason shops that should have
chosen variants did not.
\end{measured}

\section{Uploads, and What Your Shop Will Accept}

One of the nine option types is \texttt{file}. Switch it on and a customer can
send you a file --- a logo to print, a document to quote against, a photograph
to frame.

You have now accepted arbitrary bytes from a stranger, and everything that
follows is about not regretting it.

\begin{examplebox}[What OpenCart checks, in order]
From \texttt{catalog/controller/tool/upload.php} in your own container. Read
it; it is forty lines.

\begin{tightnum}
  \item \textbf{The filename is sanitised} --- everything outside
        \texttt{[a-zA-Z0-9.\-\textvisiblespace+]} is stripped, then
        \texttt{basename()} removes any path. This defeats
        \texttt{../../etc/passwd} and it is solid.
  \item \textbf{The length is checked}, 3 to 64 characters.
  \item \textbf{The extension is checked against an allow-list}
        (\texttt{config\_file\_ext\_allowed}). An allow-list is the right
        shape: anything not named is refused.
  \item \textbf{The MIME type is checked against a second allow-list}
        (\texttt{config\_file\_mime\_allowed}).
  \item \textbf{The stored name is randomised} --- the file is saved as the
        name plus a 32-character token, so it cannot be guessed, and the real
        name is kept in the database.
\end{tightnum}
\end{examplebox}

\begin{alertbox}[Check 4 is worth almost nothing, and it is not a bug in OpenCart]
The MIME type it compares is
\texttt{\$\_FILES['file']['type']} --- and that value is \textbf{sent by the
browser}. It is a claim made by the client about its own file. Anyone using
anything other than a browser sets it to whatever passes.

\smallskip
The default allow-list makes this worse by including
\texttt{application/octet-stream}, which means ``unidentified binary'' ---
so the honest answer for almost any file is already on the list.

\smallskip
\textbf{This is not a criticism of OpenCart so much as a lesson about where a
check must run.} A check performed on data the attacker controls is not a
check. The extension allow-list is real because the server derives it from the
name it has already sanitised. The MIME check is advisory, and the only
server-side equivalent is to inspect the file's actual contents ---
\texttt{finfo\_file()} in PHP, which reads the magic bytes.
\end{alertbox}

\begin{alertbox}[Two defaults on your shop, which you can verify in five minutes]
Both of these are true of the stock install in \chref{stack}, and both are
measured in \chref{security}.

\smallskip
\textbf{The allow-list includes \texttt{svg}.} An SVG is not a picture; it is
an XML document, and it may contain \texttt{<script>}. A browser asked to
display one from your origin will run it. The default list also carries
\texttt{eps} and \texttt{ps}, which are PostScript --- a programming language
--- and \texttt{msi}, \texttt{cab} and \texttt{rar}.

\smallskip
\textbf{Uploads are served over HTTP.} They land in
\texttt{system/storage/upload/}, which is \emph{inside} the web root. Requesting
a file there returns it: a text file comes back as \texttt{text/plain} and an
SVG as \texttt{image/svg+xml}. The directory cannot be listed --- there is an
empty \texttt{index.html} --- and names carry a 32-character token, so this is
not trivially exploitable. But the protection is secrecy of the name, not
inaccessibility of the directory, and Apache is configured with
\texttt{AllowOverride None}, so the \texttt{.htaccess} OpenCart ships cannot
change it.

\smallskip
\textbf{What to do about it} is \chref{security}'s subject. The short version:
remove \texttt{svg} from the list unless you need it, serve uploads through a
script that sets \texttt{Content-Disposition: attachment} rather than from the
filesystem, and keep the storage directory outside the web root.
\end{alertbox}

\begin{conceptbox}[Two size limits, and the smaller one wins]
\texttt{config\_file\_max\_size} is a setting in the admin panel. PHP has its
own, \texttt{upload\_max\_filesize}, in \texttt{code/shop/php.ini}, and this
course sets it to 8~MB.

\smallskip
If the admin setting is larger, it has no effect: PHP refuses the upload
before any OpenCart code runs, and the error the customer sees is unhelpful.
Whenever a limit appears not to work, look for the second limit --- there is
nearly always a second limit, in the web server, the runtime or the proxy, and
the smallest one is the one in force.
\end{conceptbox}

\begin{pitfall}
\begin{tight}
  \item \textbf{Choosing options when the requirement needs variants.} The
        symptom arrives at the first stock-take, and the fix is a migration.
  \item \textbf{Believing option \texttt{quantity} tracks combinations.} It
        tracks one value. Red Large and Red Small share Red's stock.
  \item \textbf{Trusting the uploaded MIME type.} The client supplies it.
  \item \textbf{Allowing \texttt{svg} without thinking about it.} It is script
        that renders as a picture.
  \item \textbf{Setting a size limit in one place only.} The smaller of the
        admin setting and \texttt{upload\_max\_filesize} wins.
  \item \textbf{Making every option \texttt{required}.} The customer cannot
        buy without choosing, and a required text option with no sensible
        default costs sales.
  \item \textbf{Forgetting that option values are global.} Editing ``Large''
        edits it for every product that uses it.
\end{tight}
\end{pitfall}

\begin{lab}[Options, variants, and what your shop will swallow]
\begin{lstlisting}[language=bash]
# --- 1. Look at a product that already has every option type ---------
# Product 42 in the demo catalogue carries all nine. Open it in the
# admin panel: Catalog > Products > edit > Option.
cd code
docker compose exec db mariadb -ushop -pshoppw opencart

# --- 2. Count what one option set costs ------------------------------
# Add a Size option with three values, through the admin panel. Then
# count the rows it created. Predict the number before you look.
#   SELECT COUNT(*) FROM oc_product_option       WHERE product_id = 42;
#   SELECT COUNT(*) FROM oc_product_option_value WHERE product_id = 42;

# --- 3. Find the row that does not exist -----------------------------
# Set stock on two option values. Then try to write a query that
# answers "how many Red Large do we hold?". You cannot. Write down,
# in one sentence, why -- that sentence is the whole of section 7.1.

# --- 4. Measure options against variants -----------------------------
python m07_variations.py

# --- 5. Read the upload controller -----------------------------------
# Forty lines. Find the five checks, and find the one the client
# controls.
docker compose exec shop \
  cat /var/www/html/catalog/controller/tool/upload.php

# --- 6. Find out what your shop allows -------------------------------
docker compose exec db mariadb -ushop -pshoppw opencart -e \
  "SELECT \`key\`,\`value\` FROM oc_setting \
    WHERE \`key\` LIKE 'config_file%';"

# --- 7. Prove the upload directory is served -------------------------
# Put a file there yourself and fetch it. This is YOUR container; see
# the scope statement in Appendix A before you try this anywhere else.
docker compose exec shop sh -c \
  'echo hello > /var/www/html/system/storage/upload/probe.txt'
curl -i http://localhost:8090/system/storage/upload/probe.txt
docker compose exec shop rm /var/www/html/system/storage/upload/probe.txt

# --- 8. Decide, and write it down ------------------------------------
# For the shop you described in Chapter 4: options or variants? One
# paragraph, with the stock question answered explicitly. Chapter 23
# marks you on whether the catalogue matches this decision.
\end{lstlisting}

\smallskip
\textbf{What to notice.} Step 3 is the one that teaches. Everybody believes
option stock tracks combinations until they try to write the query, and the
minute spent failing to write it is worth more than this chapter's prose.
\end{lab}

\begin{checkpoint}
\begin{tightnum}
  \item A shop sells mugs with a customer-supplied photograph printed on them,
        in two sizes. Which choices are options and which are variants, and
        why? Say what the shop must hold in stock.
  \item A product has option sets of 4, 3 and 2 values. How many rows in
        \texttt{oc\_product\_option\_value}, and how many
        \texttt{oc\_product} rows would the same catalogue need as variants?
  \item An attacker uploads a file named \texttt{logo.svg} containing
        \texttt{<script>}, sending \texttt{Content-Type:} \texttt{image/png}.
        Which of OpenCart's five checks does it pass, which does it fail, and
        what would have to be true for this to harm a visitor?
\end{tightnum}
\end{checkpoint}

\begin{chaptersummary}
\begin{tight}
  \item An option modifies one product; a variant is a product of its own with
        \texttt{master\_id} pointing at a master. OpenCart~4 supports both and
        shows options first.
  \item Options store choices, not combinations. Stock, price and weight sit on
        the option \emph{value}, so there is no row for ``Red Large'' and no
        query can produce one.
  \item That makes three things impossible: stock per combination, a SKU per
        combination, and reporting per combination. If any matters, use
        variants, and decide before the catalogue is built.
  \item The test: could a warehouse person pick up exactly this combination and
        scan it? Then it is a variant.
  \item Measured: option sets of 6, 8 and 5 values are 19 rows and express 240
        combinations; the same catalogue as variants is 241 product rows.
        Options grow by addition, variants by multiplication --- and the page
        time is unaffected either way, because the page renders choices.
  \item The real cost of variants is catalogue management by a person, not
        server time.
  \item A \texttt{file} option accepts arbitrary bytes from a stranger.
        OpenCart sanitises the name, checks the extension against an allow-list
        and randomises the stored name --- all sound. Its MIME check is not,
        because the client supplies the MIME type.
  \item On the stock install \texttt{svg} is allowed and the upload directory
        is served over HTTP. The protection is an unguessable name, not an
        unreachable directory. \chref{security} fixes both.
\end{tight}
\end{chaptersummary}

\begin{reviewq}
  \item \textbf{(MCQ)} Stock quantity for an option lives on:
        (a)~\texttt{oc\_product}; (b)~\texttt{oc\_product\_option};
        (c)~\texttt{oc\_product\_option\_value}; (d)~a combination table.
  \item \textbf{(MCQ)} Three option sets of 5, 4 and 3 values imply how many
        combinations, and cost how many \texttt{oc\_product\_option\_value}
        rows? (a)~12 and 12; (b)~60 and 12; (c)~60 and 60; (d)~12 and 60.
  \item \textbf{(Short)} Give the one-question test for options against
        variants, and apply it to a shop selling curtains cut to a
        customer-specified width.
  \item \textbf{(Short)} Why is checking \texttt{\$\_FILES['file']['type']}
        not a security control, and what is the server-side equivalent?
  \item \textbf{(Short)} An administrator raises the upload size limit in the
        admin panel and large files are still refused. Explain, and say where
        else to look.
  \item \textbf{(Applied)} A clothing shop has 300 styles, each in 5 sizes and
        4 colours. Compute the row counts for options and for variants across
        \texttt{oc\_product} and \texttt{oc\_product\_description}. Then make a
        recommendation, naming the single requirement that decides it.
  \item \textbf{(Applied)} Design the upload handling you would write for a
        shop that must accept customer artwork. State every check, say which
        side it runs on, say what each one defends against, and say what you
        would do with the file after accepting it so that serving it later is
        safe.
\end{reviewq}

%=====================================================================
\chapter{Themes, Layout and the Customer's Experience}\label{ch:ux}
%=====================================================================
\locator{2}{Part II --- Standing the Shop Up}

\begin{outcomes}
By the end of this chapter you will be able to:
\begin{tight}
  \item \textbf{Describe} how a theme is built from templates, and
        \textbf{override} one without editing the original.
  \item \textbf{Explain} what a layout is, and \textbf{place} a module on a
        chosen page without touching a template at all.
  \item \textbf{Measure} a page's weight and \textbf{say} which part of it is
        the weight.
  \item \textbf{Distinguish} the optimisations that are worth doing from the
        ones that are merely popular, using evidence.
  \item \textbf{Judge} a storefront change by its effect on the customer rather
        than by taste.
\end{tight}
\end{outcomes}

\begin{prereq}
\chref{catalogue} and \chref{variations}: a catalogue worth looking at.
Appendix~B's Twig section --- this chapter edits templates. HTML and CSS from
\emph{Web Technologies}.
\end{prereq}

\begin{keyterms}
theme $\bullet$ template $\bullet$ Twig $\bullet$ template override $\bullet$
layout $\bullet$ module $\bullet$ position $\bullet$ page weight $\bullet$
request count $\bullet$ render-blocking $\bullet$ minification $\bullet$
responsive design $\bullet$ above the fold $\bullet$ conversion
\end{keyterms}

\section{What a Theme Actually Is}

``Change the theme'' sounds like changing a colour. It is not. A theme in
OpenCart is a directory of Twig templates, and each one is responsible for one
piece of the page.

\begin{lstlisting}[language=bash]
docker compose exec shop \
  ls /var/www/html/catalog/view/template/product/
# category.twig  product.twig  search.twig  special.twig ...
\end{lstlisting}

The default theme lives in
\texttt{catalog/view/template/}. The important rule about it is one you should
adopt before you change a single line:

\begin{alertbox}[Never edit the default theme]
Edit \texttt{catalog/view/template/product/category.twig} directly and three
things happen. Your change is lost at the next upgrade. Nobody can tell what
you changed, because there is no original to compare against. And a second
developer, or you in March, will assume the file is stock and read it as
documentation of how OpenCart works.

\smallskip
Make a \term{child theme} instead: a directory of your own that contains
\emph{only} the templates you changed. OpenCart looks there first and falls
back to the default for everything else, so your directory \emph{is} the list
of your changes. A child theme with four files in it is a complete and honest
statement that you changed four things.
\end{alertbox}

\section{Layouts, or How to Change a Page Without Editing It}

A surprising amount of what people reach for a template to do can be done
without touching one.

\begin{definitionbox}[Layout, Module, Position]
A \term{module} is a self-contained block --- a category list, a banner, a
carousel, ``featured products''.

\smallskip
A \term{layout} says which modules appear on which \emph{kind} of page, and
where. \adminpath{Design, Layouts} lists them: Home, Category, Product,
Checkout, and so on.

\smallskip
A \term{position} is a slot in the page: \uiel{Column Left},
\uiel{Column Right}, \uiel{Content Top}, \uiel{Content Bottom}. A module is
placed in a position of a layout, with a sort order.

\smallskip
So ``put a banner above the products on every category page'' is three clicks
in \adminpath{Design, Layouts, Category} and no code at all.
\end{definitionbox}

\begin{conceptbox}[Ask this before editing a template]
\textbf{``Is this a layout change or a template change?''}

\smallskip
If you want something \emph{moved, added or removed} and it already exists as a
module, it is a layout change --- do it in the admin panel, where it is
visible to whoever comes next and survives upgrades.

\smallskip
Edit a template only when the \emph{markup itself} must differ: a field shown
that the theme does not show, a different arrangement within a block, a tag the
theme gets wrong. Then override exactly that file in your child theme and
nothing else.

\smallskip
The commonest waste of effort in this chapter's laboratory is a student editing
\texttt{category.twig} to do something \adminpath{Design, Layouts} would have
done.
\end{conceptbox}

\section{What a Page Weighs}

Here the chapter stops being about arrangement. A customer's experience of your
shop is mostly its speed, and speed is mostly \emph{weight} --- how many bytes
and how many separate requests the browser must complete before the page is
usable.

Note that this is a different question from \chref{catalogue}'s. That chapter
measured how long the \emph{server} took to produce the HTML. This one measures
everything that happens afterwards, and it is usually larger.

\begin{measured}[What the Storefront Weighs]
Produced by \texttt{code/m08\_pageweight.py} on the machine in Appendix~A. The
harness fetches a page, finds every stylesheet, script and image the HTML
names, requests each one once, and totals them.

\rowcolors{2}{white}{lightbg}
\begin{longtable}{@{}L{2.0cm}L{1.9cm}L{1.6cm}L{1.6cm}L{1.6cm}L{1.6cm}L{2.0cm}@{}}
\toprule
\rowcolor{primary!15}
\textbf{Page} & \textbf{Requests} & \textbf{HTML} & \textbf{CSS} &
\textbf{JS} & \textbf{Images} & \textbf{Total} \\
\midrule
\endhead
Home & 25 & 32\,KB & 348\,KB & 223\,KB & 293\,KB & \textbf{895\,KB} \\
Category & 20 & 46\,KB & 348\,KB & 223\,KB & 122\,KB & \textbf{738\,KB} \\
Product & 18 & 54\,KB & 354\,KB & 243\,KB & 61\,KB & \textbf{713\,KB} \\
\bottomrule
\end{longtable}

\smallskip
\textbf{The framework is the page.} CSS and JavaScript come to about
\textbf{570\,KB on every page}, and they do not vary with the content ---
the same bytes on the home page, a category of 50{,}000 products and a single
product. On the product page they are \textbf{84\% of the weight}. The HTML
that the whole of \chref{catalogue} was spent optimising is 6--7\%.

\smallskip
\textbf{Images are not the problem here, although everyone optimises them
first.} They are 33\% of the home page, 16\% of a category page and 9\% of a
product page. The largest single image in the demonstration catalogue is
13.9\,KB. You could halve every image on the category page and save 61\,KB ---
about 8\%.

\smallskip
\textbf{The single largest file is 265\,KB of unminified CSS.}
\texttt{bootstrap.css} ships at 271{,}637 bytes, and \textbf{no minified
version ships with it} --- there is no \texttt{bootstrap.min.css} anywhere in
the tree. A minified build of the same framework is around a tenth of that.
That one file is 36\% of the category page.

\smallskip
\textbf{This measurement undercounts, and you should know how.} It follows only
assets the \emph{HTML} names. Fonts are requested by the CSS, so they are
invisible to it: the FontAwesome directory alone holds 248\,KB of
\texttt{.woff2} files, of which a browser fetches whichever subsets it needs.
The totals above are therefore a floor, not a figure. A tool that drove a real
browser would see more; this one is honest about seeing less.
\end{measured}

\begin{conceptbox}[The order to optimise in, which is not the popular order]
From the numbers above, for this shop:

\begin{tightnum}
  \item \textbf{Minify the CSS.} One file, 265\,KB to roughly 27\,KB. The
        largest single saving available, and it changes nothing a customer can
        see.
  \item \textbf{Drop what you do not use.} Bootstrap and FontAwesome are
        general-purpose libraries; a shop uses a fraction of either. Removing
        the unused parts is the second-largest saving and the fiddliest.
  \item \textbf{Cut the request count.} Twenty requests is twenty round trips.
        On a fast link this hardly matters; on a phone on a slow connection in
        Bahawalpur it is most of the wait.
  \item \textbf{Then images} --- resize to the size actually displayed, and
        serve modern formats.
\end{tightnum}

\smallskip
Most advice reverses this list, because images are the thing people can see and
a 2\,MB photograph is a vivid story. On \emph{this} shop images are 16\% and
the framework is 77\%. On a shop with photographic product shots the balance
moves. That is the point: the order depends on the measurement, and the
measurement takes two minutes.
\end{conceptbox}

\begin{alertbox}[Render-blocking, and why position matters as much as size]
A stylesheet in the \texttt{<head>} is \term{render-blocking}: the browser will
not paint \emph{anything} until it has fetched and parsed it. So 348\,KB of CSS
is not merely 348\,KB of transfer --- it is 348\,KB standing between the
customer and the first pixel.

\smallskip
Scripts block too, unless marked \texttt{defer} or \texttt{async}. A
\texttt{<script src>} in the head stops the parser dead while it downloads.

\smallskip
This is why page weight and page time are not the same measurement, and why
``it is only 200\,KB'' can still be the reason a shop feels slow.
\end{alertbox}

\section{Judging a Change by the Customer}

The hardest discipline in this chapter is not technical. It is resisting the
urge to change the storefront because you would prefer it differently.

\begin{examplebox}[Three changes, and how you would know if they helped]
\textbf{``Make the Add to Cart button green.''} You cannot know without an
experiment: show half your visitors one colour and half the other, and compare
the proportion who buy. With thirty visitors a day you will wait months for an
answer, which is itself the answer --- a shop this size should spend its
effort elsewhere.

\smallskip
\textbf{``Show the delivery charge on the product page.''} This you can reason
about. \chref{what}'s laboratory found that shops which hide the delivery
charge until checkout lose customers there; it is the most-cited reason for
abandonment. No experiment needed.

\smallskip
\textbf{``Minify the CSS.''} Measurable directly, in bytes, with no customers
involved at all. 265\,KB to 27\,KB is a fact, and it cannot make anything
worse.

\smallskip
\textbf{The pattern.} Prefer changes you can verify without an experiment you
cannot afford to run. Most small shops cannot generate enough traffic to
measure a colour, and should spend their effort on the things that are
measurable in bytes or known from everybody else's evidence.
\end{examplebox}

\begin{pitfall}
\begin{tight}
  \item \textbf{Editing the default theme.} Lost at upgrade, invisible to the
        next person. Override in a child theme.
  \item \textbf{Editing a template to do a layout's job.} If the block exists
        as a module, move it in \adminpath{Design, Layouts}.
  \item \textbf{Optimising images first because they are visible.} On this shop
        that addresses 16\% of the weight and ignores 77\%.
  \item \textbf{Assuming a framework file is minified.} \texttt{bootstrap.css}
        here is 265\,KB and there is no minified copy in the tree.
  \item \textbf{Measuring page weight by counting every \texttt{<img>} tag.}
        The same thumbnail on twenty cards is one request; counting it twenty
        times overstates the weight several-fold.
  \item \textbf{Forgetting fonts.} They are requested by CSS, so they are
        invisible to any harness that reads only the HTML --- including this
        chapter's.
  \item \textbf{Redesigning on taste with no way to tell whether it helped.}
\end{tight}
\end{pitfall}

\begin{lab}[Weigh your shop, then make it lighter]
\begin{lstlisting}[language=bash]
# --- 1. Weigh three pages -------------------------------------------
cd code
python m08_pageweight.py --compare

# --- 2. Find the biggest single file --------------------------------
# Predict it before you look. Most people guess an image.
docker compose exec shop sh -c \
  'ls -laS /var/www/html/catalog/view/stylesheet/ | head -5'

# --- 3. Make a child theme ------------------------------------------
# Admin: Design > Theme Editor. Override ONE template -- the category
# product card -- and add the model number under each product name.
# Your override directory is now the list of your changes.

# --- 4. Do a layout change with no code -----------------------------
# Admin: Design > Layouts > Category. Add a module to Content Top.
# Reload a category page. No template was edited.

# --- 5. Weigh it again ----------------------------------------------
# Did your two changes move the number? By how much, and which one?
python m08_pageweight.py

# --- 6. Make a real saving ------------------------------------------
# Minify bootstrap.css and serve the minified copy instead. Any
# minifier will do. Record the before and after byte counts.

# --- 7. Weigh it a third time, and account for the difference -------
# The saving should be most of 240 KB. If it is not, find out what
# else changed -- do not just record the number.
python m08_pageweight.py

# --- 8. Say what you would do next, and why -------------------------
# One paragraph, ordered by measured saving rather than by what is
# easiest to see. Chapter 20 returns to this with caching.
\end{lstlisting}

\smallskip
\textbf{What to notice.} Step 2 against step 1: almost everyone predicts an
image and the answer is a stylesheet. That gap between where people look for
weight and where weight is, is the whole chapter.
\end{lab}

\begin{checkpoint}
\begin{tightnum}
  \item A product page is 713\,KB, of which 597\,KB is CSS and JavaScript. A
        colleague proposes compressing all the product photographs. What is the
        best case for that change, and what would you propose instead?
  \item Distinguish a layout change from a template change, and give one
        example of each for the request ``show a trust badge under the Add to
        Cart button on every product page''.
  \item This chapter's harness reports 738\,KB for a category page, and a
        browser's developer tools report more. Give two reasons the harness is
        lower, and say which of them it could fix.
\end{tightnum}
\end{checkpoint}

\begin{chaptersummary}
\begin{tight}
  \item A theme is a directory of Twig templates. Never edit the default:
        override only the files you change in a child theme, so that the
        directory is the list of your changes and survives upgrades.
  \item A layout places modules in positions on a kind of page. Moving, adding
        or removing an existing block is a layout change made in the admin
        panel, not a template edit.
  \item Measured: the storefront weighs 713--895\,KB per page across 18--25
        requests. CSS and JavaScript are about 570\,KB of that \emph{on every
        page} and do not vary with content --- 84\% of the product page.
  \item Images are 9--33\% depending on the page, although they are what people
        optimise first.
  \item The largest single file is \texttt{bootstrap.css} at 265\,KB,
        unminified, with no minified copy anywhere in the tree.
  \item The harness undercounts: it follows only assets named in the HTML, so
        fonts requested by CSS (248\,KB in one directory) are invisible to it.
        Know what your instrument cannot see.
  \item Weight and time are not the same: a stylesheet in the head is
        render-blocking and stands between the customer and the first pixel.
  \item Prefer changes you can verify without an experiment you cannot afford
        to run. A small shop cannot measure a button colour and can measure
        240\,KB.
\end{tight}
\end{chaptersummary}

\begin{reviewq}
  \item \textbf{(MCQ)} On the measured product page, CSS and JavaScript are
        what share of the weight? (a)~9\%; (b)~about a third; (c)~84\%;
        (d)~it depends on the catalogue size.
  \item \textbf{(MCQ)} ``Show the manufacturer's name under the product title''
        where the theme does not show it at all is: (a)~a layout change;
        (b)~a template override; (c)~a module setting; (d)~a database change.
  \item \textbf{(Short)} Why is a child theme better than editing the default,
        giving two reasons that are not about upgrades.
  \item \textbf{(Short)} Explain render-blocking, and why page weight and page
        time are different measurements.
  \item \textbf{(Short)} This chapter's harness counts each unique asset once
        rather than once per mention. Explain what would go wrong otherwise,
        with a number.
  \item \textbf{(Applied)} You are given a shop whose category page is 2.1\,MB:
        HTML 40\,KB, CSS 90\,KB, JS 150\,KB, images 1.82\,MB across 40
        requests. Give your optimisation order with a predicted saving for each
        step, and say how this differs from the order you would use on the shop
        in this chapter and why.
  \item \textbf{(Applied)} Design a measurement that would tell you whether
        showing the delivery charge on the product page increases or decreases
        completed orders. State what you would record, how long you would need
        to run it, and --- honestly --- whether a shop taking thirty orders a
        day could reach a conclusion at all.
\end{reviewq}


\part{From Basket to Order}
%=====================================================================
% PART III OPENER
%=====================================================================
\thispagestyle{plain}
\structuralfigures{P3.}

\begin{center}
\begin{tcolorbox}[enhanced, width=0.92\textwidth, colback=greenhl!5,
  colframe=greenhl, boxrule=0pt, leftrule=5pt, arc=3pt, top=8pt, bottom=8pt]
\large\itshape\color{greenhl!60!black}
A catalogue is a brochure until somebody can buy from it. Part~III is the path from a product page to a confirmed order: where the basket lives, what the checkout is really doing, and the two things that change the total after the customer thought they knew it --- shipping with tax, and discounts.
\end{tcolorbox}
\end{center}

\vspace{6pt}

\dsfigh{%
\begin{tikzpicture}[
  cbox/.style={rounded corners=3pt, draw=greenhl, line width=1pt, fill=greenhl!7,
               text=greenhl!55!black, font=\small\bfseries, align=center,
               minimum width=3.5cm, minimum height=1.0cm},
  hbox/.style={rounded corners=3pt, draw=greenhl, line width=1.8pt, fill=greenhl!16,
               text=greenhl!55!black, font=\small\bfseries, align=center,
               minimum width=3.5cm, minimum height=1.0cm},
  arr/.style={-{Stealth[length=2.4mm]}, greenhl!60, line width=1pt},
  note/.style={font=\scriptsize\itshape, text=greenhl!55!black, align=center,
               fill=white, inner sep=2pt}
]
  \node[cbox] (c0) at (0.00,0) {Ch.\ 9\\The Basket};
  \node[hbox] (c1) at (4.30,0) {Ch.\ 10\\Checkout};
  \node[cbox] (c2) at (8.60,0) {Ch.\ 11\\Shipping and Tax};
  \node[cbox] (c3) at (12.90,0) {Ch.\ 12\\Discounts};
  \draw[arr] (c0) -- (c1);
  \draw[arr] (c1) -- (c2);
  \draw[arr] (c2) -- (c3);
  \node[note] at (2.15,-1.0) {an intention};
  \node[note] at (6.45,-1.0) {a commitment};
  \node[note] at (10.75,-1.0) {what it really costs};
  \node[font=\scriptsize\itshape, text=accent, align=center] at (4.30,1.15)
       {the order state machine: everything after it is a transition};
\end{tikzpicture}}%
{Reading path through Part~III. Chapter 10 is the hinge of the whole book --- an order is the first object in the course that cannot simply be deleted and redone.}{part3map}

\newpage

%=====================================================================
\chapter{The Shopping Basket, and Where State Lives}\label{ch:basket}
%=====================================================================
\locator{3}{Part III --- From Basket to Order}

\begin{outcomes}
By the end of this chapter you will be able to:
\begin{tight}
  \item \textbf{Explain} why HTTP has no memory, and \textbf{name} the three
        places a shop can keep what it remembers about a visitor.
  \item \textbf{Locate} a basket in a running shop, and \textbf{state} which
        store holds it, with evidence rather than with belief.
  \item \textbf{Measure} what a session costs in rows, bytes and queries, and
        \textbf{distinguish} the cost of a visitor from the cost of a customer.
  \item \textbf{Predict} what survives a restart, a redeployment and a volume
        deletion, and \textbf{verify} each prediction.
  \item \textbf{Judge} a session configuration --- engine, expiry, cookie
        flags --- against what the shop actually does.
\end{tight}
\end{outcomes}

\begin{prereq}
\chref{stack}: containers, volumes, and the difference between restarting a
container and deleting one. \chref{catalogue}: the tables behind the screens,
and how to read a query the server actually ran rather than one you assumed it
ran. From Web Technologies: HTTP requests and responses, and what a cookie is.
\end{prereq}

\begin{keyterms}
state $\bullet$ stateless protocol $\bullet$ cookie $\bullet$ session
$\bullet$ session identifier $\bullet$ session engine $\bullet$ server-side
state $\bullet$ client-side state $\bullet$ persistent cart $\bullet$
abandoned basket $\bullet$ session expiry $\bullet$ garbage collection
$\bullet$ volume $\bullet$ cache miss
\end{keyterms}

\section{The Problem: A Protocol With No Memory}

A customer puts a laptop in the basket, browses to another category, comes back,
and the laptop is still there. Nothing about HTTP makes that happen.

Every HTTP request is independent. The server receives a request, answers it,
and keeps nothing. The next request from the same person arrives looking exactly
like a request from a stranger. If the shop is to know that these two requests
belong to the same visitor, \emph{something carried between them} must say so,
and that something can only travel in the request itself --- because the request
is all there is.

\begin{definitionbox}[State, Cookie, Session]
\term{State} is anything the shop remembers between one request and the next:
what is in the basket, who is logged in, which currency is being shown.

\smallskip
A \term{cookie} is a small named value the server asks the browser to store and
send back on every subsequent request to that site. It is the only standard way
to make a second request identifiably related to a first. It lives on the
customer's machine, which means the customer can read it, change it and delete
it.

\smallskip
A \term{session} is state the \emph{server} keeps, filed under a key. The
browser carries only the key --- the \term{session identifier} --- in a cookie;
the data stays on the server, where the customer cannot touch it.

\smallskip
The distinction is the whole of this chapter. A cookie is a place. A session is
an arrangement that \emph{uses} a cookie to point at a place somewhere else.
\end{definitionbox}

\section{Three Places State Can Live}

Given that the browser must carry something, a shop has three choices about
\emph{how much} it carries, and they trade off against each other in ways worth
seeing before measuring anything.

\dsfig{%
\begin{tikzpicture}[
  store/.style={rounded corners=3pt, draw=#1, line width=1.1pt, fill=#1!8,
                text=#1!60!black, font=\bfseries\small, align=center,
                text width=3.4cm, minimum height=1.6cm},
  note/.style={font=\scriptsize, text=neutral, align=left, text width=4.9cm,
               anchor=west},
  flow/.style={-{Stealth[length=2.2mm]}, primary!55, line width=0.9pt},
  tagline/.style={font=\scriptsize\bfseries, text=#1, align=center}
]
  \node[store=accent]    (ck) at (0,4.3)  {In the cookie\\{\scriptsize
       \mdseries the browser carries it all}};
  \node[store=secondary] (ss) at (0,2.15) {In the session\\{\scriptsize
       \mdseries server keeps it, cookie points}};
  \node[store=purplehl]  (tb) at (0,0)    {In a table\\{\scriptsize
       \mdseries a row per basket line}};

  \node[note] at (2.1,4.3)  {Survives nothing you control. The customer can
       edit it, and so can anyone who reads the traffic. Capped at about
       4\,KB. No server storage at all.};
  \node[note] at (2.1,2.15) {Server decides what is true. Lost when the
       session expires or its store is cleared. Grows with visitors, not with
       customers.};
  \node[note] at (2.1,0)    {Queryable, joinable, reportable, and still there
       tomorrow. Costs a write on every visit, and the rows outlive the
       visitors who made them.};

  \draw[flow] (ck) -- (ss);
  \draw[flow] (ss) -- (tb);
  \node[tagline=accent, anchor=west] at (-1.9,5.55)
       {Less server cost, less trust, less memory \ \faIcon{arrow-down}};
  \node[tagline=purplehl, anchor=west] at (-1.9,-1.25)
       {\faIcon{arrow-up} \ More server cost, more trust, more memory};
\end{tikzpicture}}%
{The three stores, and what each one buys. Nothing forces a shop to pick one:
most use all three at once, for different things. The question worth asking of
any shop is not which one it uses, but which one holds the basket.}{threestores}

\begin{conceptbox}[A rule of thumb, before any measurement]
Put in the \textbf{cookie} only what you would not mind the customer editing ---
a chosen language, a dismissed banner. Put in the \textbf{session} what must be
true for this visit and need not outlive it. Put in a \textbf{table} anything
you will later want to count, report on, or still have tomorrow.

\smallskip
Applying that rule to a basket gives an answer immediately: a shop wants to know
how many baskets were abandoned last week, so the basket belongs in a table.
That is an argument. The rest of this chapter is evidence.
\end{conceptbox}

\section{Where This Shop Actually Puts It}

Almost every tutorial on this subject says the cart is kept in the session.
Before believing that of the shop in front of you, look.

The method is the one \chref{catalogue} arrived at the hard way: do not reason
about what the code probably does, make the shop do it and read what changed.
Add three products, and after each one measure all three stores at once.

\begin{measured}[Add three products; watch all three stores]
Produced by \texttt{code/m09\_basket.py --where} against the stack in
Appendix~A. ``Cookie'' is what the browser sends up on every request;
``session'' is the length of the stored session record; ``cart'' is rows in
\texttt{oc\_cart} and the bytes in them.

\rowcolors{2}{white}{lightbg}
\begin{longtable}{@{}L{4.0cm}L{2.3cm}L{2.4cm}L{2.3cm}L{2.2cm}@{}}
\toprule
\rowcolor{primary!15}
\textbf{After adding} & \textbf{Cookie} & \textbf{Session} &
\textbf{Cart rows} & \textbf{Cart bytes} \\
\midrule
\endhead
(nothing) & 49\,B & 18\,B & 0 & 0 \\
2 $\times$ MacBook & 49\,B & 18\,B & 1 & 53 \\
1 $\times$ iPhone & 49\,B & 18\,B & 2 & 106 \\
3 $\times$ HTC Touch HD & 49\,B & 18\,B & 3 & 159 \\
\bottomrule
\end{longtable}

\smallskip
\textbf{The cookie does not move, and the session does not move.} Both are
constant across an empty basket and a three-line one. The entire session record
is eighteen bytes:

\smallskip
\hspace*{1em}\texttt{\{"currency":"USD"\}}

\smallskip
\textbf{The basket is a table.} It is rows in \texttt{oc\_cart}, one per line,
about 53 bytes each, keyed by the session identifier. The cookie carries a
26-character identifier and a currency, 49 bytes in all, and that is the whole
of what the browser knows.

\smallskip
So the common claim is false here. It is worth knowing \emph{why} it is so
widely believed: in a simpler shop it is usually true, and the arrangement
OpenCart uses --- session for identity, table for the basket --- is the one a
shop adopts when it wants to report on baskets that were never checked out.
\end{measured}

\begin{alertbox}[``The cart is in the session'' --- check before you repeat it]
This matters beyond trivia. If you believe the basket is in the session, you
will also believe:

\begin{tight}
  \item that lengthening the session expiry lengthens the life of a basket
        (here it does both, but for two different reasons, in two places);
  \item that switching the session engine moves the basket (it does not --- see
        the next measurement);
  \item that clearing sessions clears baskets (it does not; it orphans them);
  \item that you cannot report on abandoned baskets without extra work (you
        can: they are rows, and \chref{dashboard} reads them).
\end{tight}

Each of those is a decision someone will make on a Friday afternoon. The
evidence took one command.
\end{alertbox}

\section{What a Session Costs}

The basket is in a table, but the shop still keeps sessions, and sessions are
not free. Two questions: how many are created, and what does each one cost per
request?

\begin{measured}[A session is written before the visitor asks for anything]
From \texttt{code/m09\_basket.py --sessions}. Twenty requests made as twenty
first-time visitors, then twenty made by one visitor who keeps the cookie.

\smallskip
\begin{tabular}{@{}L{7.2cm}L{6.0cm}@{}}
\textbf{20 first-time visitors} & \textbf{20} new session rows \\
\textbf{1 visitor, 20 page views} & \textbf{1} new session row \\
\end{tabular}

\smallskip
The session table on this machine held \textbf{645 rows, about 35.6\,KB},
accumulated entirely by the measurements in this handout --- no human has ever
shopped here.

\smallskip
\textbf{Every request writes a session, including the first.} A visitor who
loads one page and leaves has cost a row. A crawler that ignores cookies costs
one row per URL it touches, and the rows are removed on a clock rather than on
departure: \texttt{config\_session\_expire} is \textbf{86{,}400 seconds},
a full day.

\smallskip
This is the first number in the course that scales with \emph{traffic} rather
than with \emph{business}. The catalogue in \chref{catalogue} grew because the
shop sold more things. This grows because more robots visited.
\end{measured}

\begin{measured}[Database sessions against file sessions]
From \texttt{code/m09\_basket.py --engine}. OpenCart ships three session
engines --- \texttt{db}, \texttt{file} and \texttt{redis} --- selected in
\texttt{system/config/catalog.php}. The storefront ships configured for
\texttt{db}. Both rows below are the same home page, one engine apart.

\rowcolors{2}{white}{lightbg}
\begin{longtable}{@{}L{2.4cm}L{2.0cm}L{2.0cm}L{2.6cm}L{2.1cm}L{1.7cm}@{}}
\toprule
\rowcolor{primary!15}
\textbf{Engine} & \textbf{Page} & \textbf{Queries} & \textbf{On session} &
\textbf{Cart rows} & \textbf{Files} \\
\midrule
\endhead
\texttt{db} & 89\,ms & 117 & 2 & 1 & 1 \\
\texttt{file} & 71\,ms & 115 & 0 & 1 & 14 \\
\bottomrule
\end{longtable}

\smallskip
\textbf{Exactly two queries per request are the session's} --- one
\texttt{SELECT} to read it and one \texttt{REPLACE} to write it back. Moving to
files removes both and takes about 20\% off the page on this machine, where the
database and the web server are two containers on one host.

\smallskip
\textbf{The basket did not move.} The \texttt{Cart rows} column is 1 under both
engines. The engine decides where eighteen bytes live; it has nothing to do
with the basket.

\smallskip
\textbf{Files are not free either, they are just billed elsewhere.} Each session
file holds those same eighteen bytes and occupies a 4\,KB filesystem block ---
an amplification of more than two hundred. And a file session is tied to one
machine's disk, so the moment the shop runs on two web servers, a customer
served by the second one is a stranger. That, not speed, is why the shipped
default is \texttt{db}.
\end{measured}

\begin{conceptbox}[Reading that 20\% honestly]
It was measured on one host where the database is a container a virtual network
hop away. On a deployment where the database is a separate machine, removing two
round trips per request would save more. On a deployment behind a load balancer,
the file engine would not work at all.

\smallskip
A measurement tells you what happened on the machine it ran on. Appendix~A names
that machine for exactly this reason. Quoting 20\% as a property of OpenCart
would be a misuse of your own result.
\end{conceptbox}

\section{Two Things the Source Does That the Screens Do Not Show}

Reading \texttt{system/library/cart/cart.php} and
\texttt{system/library/session.php} turns up two behaviours that no admin screen
mentions and both measurements confirm.

\begin{measured}[An empty basket costs more than a full one]
From \texttt{code/m09\_basket.py --cache}. The same home page, with the basket
holding 0, 1, 2 and 3 lines.

\smallskip
\rowcolors{2}{white}{lightbg}
\begin{longtable}{@{}L{3.2cm}L{3.0cm}L{3.4cm}L{3.4cm}@{}}
\toprule
\rowcolor{primary!15}
\textbf{Basket lines} & \textbf{Queries} & \textbf{\texttt{SELECT} on cart} &
\textbf{\texttt{DELETE} on cart} \\
\midrule
\endhead
0 & 117 & \textbf{6} & 1 \\
1 & 119 & 1 & 1 \\
2 & 126 & 1 & 1 \\
3 & 129 & 1 & 1 \\
\bottomrule
\end{longtable}

\smallskip
\textbf{An empty basket is read six times; a full one is read once.} The cache
in \texttt{Cart::getProducts()} begins \texttt{if (!\$this->data)}, which is
true both when the basket has not been loaded yet \emph{and} when it is empty.
PHP cannot tell those apart that way, so an empty basket misses its own cache
on every call --- and a shop's commonest visitor has an empty basket.

\smallskip
\textbf{Every storefront page deletes before it reads.} The
\texttt{DELETE} column is 1 on every row, including the empty basket. The
\texttt{Cart} constructor removes expired guest baskets each time it is built,
so a page view is a write even when the visitor only looked.
\end{measured}

\begin{measured}[One request in five rebuilds the session table]
From \texttt{code/m09\_basket.py --gc}, over two runs.

\smallskip
\begin{tabular}{@{}L{4.6cm}L{4.2cm}L{4.4cm}@{}}
\textbf{25 requests} & \texttt{OPTIMIZE} ran 7 (28\%) & 555\,ms in total \\
\textbf{60 requests} & \texttt{OPTIMIZE} ran 15 (25\%) & 889\,ms in total \\
\end{tabular}

\smallskip
\texttt{Session::\_\_construct} registers \texttt{gc()} as a shutdown function,
so garbage collection is \emph{considered} on every request;
\texttt{session\_probability} and \texttt{session\_divisor} make it \emph{run}
on about one in five. When it runs it deletes expired rows and then issues
\texttt{OPTIMIZE TABLE}, which on InnoDB is a full table rebuild --- around
60\,ms here, on a 645-row table.

\smallskip
Both runs are printed because neither is the answer. The configured rate is
one in five; 7 of 25 and 15 of 60 are what one in five looks like in small
samples, and reporting only the first would have claimed 28\%.
\end{measured}

\section{What Survives}

\chref{stack} introduced volumes with the claim that they are what makes a
container's data outlive the container. A basket is the smallest honest test of
that claim, because you can put one somewhere and then try to destroy it three
different ways.

\begin{measured}[Three ways to stop a shop, and what is left afterwards]
From \texttt{code/m09\_basket.py --restart --destroy}. A three-line basket is
created, then each operation is applied in turn and the basket looked for again
--- in the tables, and in the page the same browser is shown.

\rowcolors{2}{white}{lightbg}
\begin{longtable}{@{}L{3.8cm}L{1.7cm}L{2.3cm}L{2.5cm}L{2.7cm}@{}}
\toprule
\rowcolor{primary!15}
\textbf{Operation} & \textbf{Took} & \textbf{Cart rows} &
\textbf{Session rows} & \textbf{Basket shown} \\
\midrule
\endhead
--- (before) & --- & 3 & 1 & yes \\
\texttt{restart shop} & 3\,s & 3 & 1 & yes \\
\texttt{down}, \texttt{up -d} & 10\,s & 3 & 1 & yes \\
\texttt{down -v}, \texttt{up -d} & 17\,s & \textbf{0} & \textbf{0} &
\textbf{no} \\
\bottomrule
\end{longtable}

\smallskip
\textbf{Destroying the containers destroyed nothing.} \texttt{down} removes the
containers; the named volumes stay, so the database files stay, so the basket
stays. The browser's cookie still pointed at a session that still existed.

\smallskip
\textbf{\texttt{-v} is the whole difference.} It deletes the volumes. The shop
reinstalled itself from scratch in 17 seconds --- which is a good property of
the stack and a bad surprise for anyone who typed it to free disk space.

\smallskip
\textbf{Notice what the third row does \emph{not} say.} It does not say the
shop broke. It came back, fast, looking perfectly healthy, with every customer's
basket gone. \chref{deploy} returns to this with backups, where the lesson is
the same one \chref{catalogue} taught about the store identifier: a shop that
answers 200 is not thereby a shop that is right.
\end{measured}

\section{The Settings That Govern Any of This}

Three of the behaviours above are configurable, and two of the three are in the
admin panel rather than in a file.

\begin{examplebox}[Where each decision is actually made]
\begin{tight}
  \item \textbf{How long an abandoned basket lives.}
        \adminpath{System, Settings, Edit, Option} --- the session expiry,
        \texttt{config\_session\_expire}, 86{,}400 seconds by default. It
        governs both the session record and the \texttt{DELETE} in the
        \texttt{Cart} constructor, so shortening it discards baskets sooner.
  \item \textbf{Whether a basket needs an account at all.}
        \adminpath{System, Settings, Edit, Option} --- \uiel{Guest Checkout}.
        With it on, the basket is keyed by session alone and vanishes with the
        cookie. With it off, every basket belongs to a customer and follows
        them to another machine.
  \item \textbf{Where the session itself is kept.} Not in the admin panel at
        all, but in a file: the \texttt{session\_engine} line of
        \texttt{system/config/catalog.php}. A setting that changes both a
        performance characteristic and a deployment constraint, and that can
        be reached only by editing a file, is worth noticing as a design
        choice --- it is not meant to be changed by the shopkeeper.
\end{tight}
\end{examplebox}

\begin{pitfall}
\begin{tight}
  \item \textbf{Assuming where the basket is instead of looking.} Every wrong
        decision in this chapter follows from that one.
  \item \textbf{Putting a price in the cookie.} Anything the browser carries,
        the customer can edit. \texttt{oc\_cart} has a \texttt{price} column
        and the shop recomputes rather than trusting it; your own extensions
        must do the same.
  \item \textbf{Lengthening session expiry to ``keep baskets''} without
        noticing that it also keeps every robot's session row for the same
        period.
  \item \textbf{Clearing \texttt{oc\_session} to tidy up.} It does not delete
        baskets; it strips them of the identifier that reaches them, leaving
        rows nothing will ever read and nothing will ever remove.
  \item \textbf{Typing \texttt{docker compose down -v} to free space.} It is
        the one command in this chapter that cannot be undone.
  \item \textbf{Reading a timing from one machine as a property of the
        software.} The 20\% above is a fact about this host.
\end{tight}
\end{pitfall}

\begin{lab}[Find the basket, then try to destroy it]
\begin{lstlisting}[language=bash]
# --- 1. Predict, in writing, before running anything -----------------
# Where is the basket: cookie, session, or a table? Commit to one.

# --- 2. Look ---------------------------------------------------------
cd code
python m09_basket.py --where

# --- 3. Read the session record yourself -----------------------------
docker compose exec db mariadb -ushop -pshoppw opencart \
  -e "SELECT session_id, LENGTH(data), data FROM oc_session LIMIT 3;"

# --- 4. Find your own basket in the table ----------------------------
docker compose exec db mariadb -ushop -pshoppw opencart \
  -e "SELECT * FROM oc_cart;"

# --- 5. Move the session to files, and check the basket --------------
python m09_basket.py --engine
# Did the basket move? Write down why not, in one sentence.

# --- 6. Count what an empty basket costs -----------------------------
python m09_basket.py --cache
# Explain the first row to someone who has not read the source.

# --- 7. Try to destroy it, in increasing order of violence -----------
python m09_basket.py --restart
# Predict each row before you read it.

# --- 8. Now the one that works ---------------------------------------
# Only when you are ready to reinstall the shop:
python m09_basket.py --restart --destroy

# --- 9. Change the expiry and show that it did something -------------
# Admin: System > Settings > Edit > Option. Set session expiry to 60
# seconds. Add a product, wait, reload. Record what you see and which
# of the two mechanisms in this chapter produced it.
\end{lstlisting}

\smallskip
\textbf{What to notice.} Step 1 against step 2 is the point of the lab. Keep
your written prediction: the gap between what you expected and what the shop
does is the only part of this exercise you will still remember in a year.
\end{lab}

\begin{checkpoint}
\begin{tightnum}
  \item A colleague says ``move the sessions to Redis and the basket will be
        faster''. Using this chapter's measurements, say what would and would
        not change, and give the number that settles it.
  \item The session table holds 645 rows on a shop that has never had a
        customer. Explain where they came from, why they are not removed when
        the visitor leaves, and what you would measure before deciding to
        shorten the expiry.
  \item \texttt{docker compose down} left a three-line basket intact and
        \texttt{docker compose down -v} did not. Explain the difference in
        terms of what each command removes, and name one other thing in your
        shop that the second command destroys.
\end{tightnum}
\end{checkpoint}

\begin{chaptersummary}
\begin{tight}
  \item HTTP keeps nothing between requests. A cookie is the only standard way
        to make the second request identifiably related to the first.
  \item A cookie is a place on the customer's machine; a session is server-side
        state that \emph{uses} a cookie to point at itself. Confusing the two
        confuses who can change what.
  \item Measured: adding three products left the cookie at 49\,bytes and the
        session at 18\,bytes, and added three rows to \texttt{oc\_cart}. The
        basket in this shop is a table, not a session.
  \item Measured: 20 first-time visitors created 20 session rows; one visitor
        making 20 page views created one. Sessions scale with traffic, not with
        trade.
  \item Measured: the \texttt{db} engine costs exactly two queries per request;
        the \texttt{file} engine costs none and took about 20\% off the page on
        this host --- while moving no part of the basket, and while making a
        second web server impossible.
  \item Measured: an empty basket is read six times per page and a full one
        once, because the cache test cannot distinguish ``empty'' from ``not
        loaded''. Every storefront page also deletes from \texttt{oc\_cart}
        before reading it.
  \item Measured: about one request in five rebuilds the session table with
        \texttt{OPTIMIZE TABLE}, at roughly 60\,ms.
  \item Measured: \texttt{restart} and \texttt{down}/\texttt{up} left the
        basket untouched; \texttt{down -v} destroyed it and the shop
        reinstalled itself in 17 seconds looking perfectly healthy.
  \item Where state lives decides what survives, what can be reported on, and
        who is able to tamper with it. It is a design decision, not a detail.
\end{tight}
\end{chaptersummary}

\begin{reviewq}
  \item \textbf{(MCQ)} In the measured shop, adding a product to the basket
        changes: (a)~the cookie; (b)~the session record; (c)~rows in a table;
        (d)~all three.
  \item \textbf{(MCQ)} Switching \texttt{session\_engine} from \texttt{db} to
        \texttt{file} moved: (a)~the basket; (b)~the session record only;
        (c)~both; (d)~neither.
  \item \textbf{(Short)} Explain why a session identifier must be unguessable,
        and what an attacker who obtained one could do in this shop.
  \item \textbf{(Short)} Give two things that belong in a cookie and two that
        must never go in one, with your reason in each case.
  \item \textbf{(Short)} An empty basket costs six cart queries and a full one
        costs one. Explain the cause, and propose a one-line change to the
        condition that would fix it.
  \item \textbf{(Applied)} A shop reports 400{,}000 session rows and 12 orders
        a day. Work out what is producing the rows, state what you would
        measure to confirm it, and give two changes you would make --- one to
        configuration and one to the shop's robots policy.
  \item \textbf{(Applied)} You are asked to make baskets survive for thirty
        days so that customers can come back to them. Say exactly what you
        would change, what each change costs, what it would break for a guest
        with no account, and how you would measure whether it increased
        completed orders at all.
\end{reviewq}

\input{parts/10_checkout}
\input{parts/11_shipping}
\input{parts/12_discounts}

\part{Money and Identity}
%=====================================================================
% PART IV OPENER
%=====================================================================
\thispagestyle{plain}
\structuralfigures{P4.}

\begin{center}
\begin{tcolorbox}[enhanced, width=0.92\textwidth, colback=purplehl!5,
  colframe=purplehl, boxrule=0pt, leftrule=5pt, arc=3pt, top=8pt, bottom=8pt]
\large\itshape\color{purplehl!70!black}
This is the part where mistakes cost money rather than marks. Taking payment, knowing who the customer is, and keeping both facts away from everybody else --- three chapters of building, and one of trying to break what you built.
\end{tcolorbox}
\end{center}

\vspace{6pt}

\dsfigh{%
\begin{tikzpicture}[
  cbox/.style={rounded corners=3pt, draw=purplehl, line width=1pt, fill=purplehl!7,
               text=purplehl!70!black, font=\small\bfseries, align=center,
               minimum width=3.5cm, minimum height=1.0cm},
  hbox/.style={rounded corners=3pt, draw=purplehl, line width=1.8pt, fill=purplehl!16,
               text=purplehl!70!black, font=\small\bfseries, align=center,
               minimum width=3.5cm, minimum height=1.0cm},
  arr/.style={-{Stealth[length=2.4mm]}, purplehl!60, line width=1pt},
  note/.style={font=\scriptsize\itshape, text=purplehl!70!black, align=center,
               fill=white, inner sep=2pt}
]
  \node[cbox] (c0) at (0.00,0) {Ch.\ 13\\Payment Systems};
  \node[cbox] (c1) at (4.30,0) {Ch.\ 14\\The Gateway};
  \node[cbox] (c2) at (8.60,0) {Ch.\ 15\\Accounts};
  \node[hbox] (c3) at (12.90,0) {Ch.\ 16\\Security};
  \draw[arr] (c0) -- (c1);
  \draw[arr] (c1) -- (c2);
  \draw[arr] (c2) -- (c3);
  \node[note] at (2.15,-1.0) {who carries the risk};
  \node[note] at (6.45,-1.0) {and what if it fails};
  \node[note] at (10.75,-1.0) {who the customer is};
  \node[font=\scriptsize\itshape, text=accent, align=center] at (12.90,1.15)
       {the only chapter that attacks the shop you built};
\end{tikzpicture}}%
{Reading path through Part~IV. Chapter 16's laboratory runs against your own container and nothing else; the chapter says so in its first paragraph, and means it.}{part4map}

\newpage

\input{parts/13_payments}
\input{parts/14_gateway}
\input{parts/15_accounts}
\input{parts/16_security}

\part{Running and Extending the Shop}
%=====================================================================
% PART V OPENER
%=====================================================================
\thispagestyle{plain}
\structuralfigures{P5.}

\begin{center}
\begin{tcolorbox}[enhanced, width=0.92\textwidth, colback=tealhl!5,
  colframe=tealhl, boxrule=0pt, leftrule=5pt, arc=3pt, top=8pt, bottom=8pt]
\large\itshape\color{tealhl!70!black}
Launching a shop is a day's work; running one is the job. Part~V is the half nobody demonstrates --- the dashboard and which of its numbers deserve action, the back office where orders are actually fulfilled and refunded, the first thing the platform cannot do, and everything that happens after you deploy.
\end{tcolorbox}
\end{center}

\vspace{6pt}

\dsfigh{%
\begin{tikzpicture}[
  cbox/.style={rounded corners=3pt, draw=tealhl, line width=1pt, fill=tealhl!7,
               text=tealhl!70!black, font=\small\bfseries, align=center,
               minimum width=3.5cm, minimum height=1.0cm},
  hbox/.style={rounded corners=3pt, draw=tealhl, line width=1.8pt, fill=tealhl!16,
               text=tealhl!70!black, font=\small\bfseries, align=center,
               minimum width=3.5cm, minimum height=1.0cm},
  arr/.style={-{Stealth[length=2.4mm]}, tealhl!60, line width=1pt},
  note/.style={font=\scriptsize\itshape, text=tealhl!70!black, align=center,
               fill=white, inner sep=2pt}
]
  \node[cbox] (c0) at (0.00,0) {Ch.\ 17\\The Dashboard};
  \node[cbox] (c1) at (4.30,0) {Ch.\ 18\\Orders, Refunds};
  \node[hbox] (c2) at (8.60,0) {Ch.\ 19\\Your First Module};
  \node[cbox] (c3) at (12.90,0) {Ch.\ 20\\Deployment};
  \draw[arr] (c0) -- (c1);
  \draw[arr] (c1) -- (c2);
  \draw[arr] (c2) -- (c3);
  \node[note] at (2.15,-1.0) {what to watch};
  \node[note] at (6.45,-1.0) {the daily work};
  \node[note] at (10.75,-1.0) {when ``no'' becomes ``yes''};
  \node[font=\scriptsize\itshape, text=accent, align=center] at (8.60,1.15)
       {the first chapter in which configuring is not enough};
\end{tikzpicture}}%
{Reading path through Part~V. Chapter 19 is where the course earns the words \emph{Application Development} in its title; everything before it could have been done without writing code.}{part5map}

\newpage

\input{parts/17_dashboard}
\input{parts/18_orders}
\input{parts/19_modules}
\input{parts/20_deploy}

\part{Reach, Law and Practice}
%=====================================================================
% PART VI OPENER
%=====================================================================
\thispagestyle{plain}
\structuralfigures{P6.}

\begin{center}
\begin{tcolorbox}[enhanced, width=0.92\textwidth, colback=goldhl!5,
  colframe=goldhl, boxrule=0pt, leftrule=5pt, arc=3pt, top=8pt, bottom=8pt]
\large\itshape\color{goldhl!70!black}
A shop that cannot be found has no customers, and a shop that ignores the law has them only briefly. Part~VI closes the course with being found, being lawful, and defending what you built.
\end{tcolorbox}
\end{center}

\vspace{6pt}

\dsfigh{%
\begin{tikzpicture}[
  cbox/.style={rounded corners=3pt, draw=goldhl, line width=1pt, fill=goldhl!7,
               text=goldhl!70!black, font=\small\bfseries, align=center,
               minimum width=3.8cm, minimum height=1.0cm},
  hbox/.style={rounded corners=3pt, draw=goldhl, line width=1.8pt, fill=goldhl!16,
               text=goldhl!70!black, font=\small\bfseries, align=center,
               minimum width=3.8cm, minimum height=1.0cm},
  arr/.style={-{Stealth[length=2.4mm]}, goldhl!60, line width=1pt},
  note/.style={font=\scriptsize\itshape, text=goldhl!70!black, align=center,
               fill=white, inner sep=2pt}
]
  \node[cbox] (c0) at (0.00,0)  {Ch.\ 21\\SEO};
  \node[cbox] (c1) at (5.40,0)  {Ch.\ 22\\Law and Ethics};
  \node[hbox] (c2) at (10.80,0) {Ch.\ 23\\The Project};
  \draw[arr] (c0) -- (c1);
  \draw[arr] (c1) -- (c2);
  \node[note] at (2.70,-1.0) {being found};
  \node[note] at (8.10,-1.0) {being answerable};
  \node[font=\scriptsize\itshape, text=accent, align=center] at (10.80,1.15)
       {everything, defended};
\end{tikzpicture}}%
{Reading path through Part~VI. Chapter 22 is the only chapter whose content is set by statute rather than by the platform, and the statute it uses is Pakistan's.}{part6map}

\newpage

\input{parts/21_seo}
\input{parts/22_legal}
\input{parts/23_project}

%---------------------------------------------------------------------
\appendix
% \part* under titlesec already emits its own contents line, so do not add
% one explicitly here or the entry appears twice.
\part*{Appendices}
%=====================================================================
\chapter{The Stack: Setting It Up and Keeping It}\label{app:stack}
%=====================================================================
\normalfigures

Everything in this book from \chref{stack} onward assumes a shop you can reach
in a browser. This appendix gets you one, in a single sitting, and lists what
to do when it refuses --- which on a university laptop it sometimes will.

\section{What You Install}

\textbf{One thing: a container runtime.} Not a web server, not PHP, not a
database. Those three run inside containers and never touch your own machine,
which is why the shop is identical on every laptop in the class and why
uninstalling it is one command.

\begin{tight}
  \item \textbf{Windows or macOS:} Docker Desktop, from
        \texttt{docker.com}. On Windows it needs WSL~2, which its installer
        offers to set up.
  \item \textbf{Linux:} the \texttt{docker.io} and \texttt{docker-compose-v2}
        packages from your distribution, then add yourself to the
        \texttt{docker} group and log out and in again.
\end{tight}

Check it with:

\begin{lstlisting}[language=bash]
docker --version
docker compose version
docker run --rm hello-world
\end{lstlisting}

The third is the one that matters. The first two only prove the program is
installed; the third proves the \emph{engine} is running, which is a different
claim and the one that usually fails.

\section{Starting the Shop}

\begin{lstlisting}[language=bash]
cd code
cp .env.example .env
docker compose up -d --build
python waitfor.py
\end{lstlisting}

The first run takes a few minutes and downloads about 1.5~GB; later runs take
seconds. \chref{stack} measures both and explains why the difference matters.

\rowcolors{2}{white}{lightbg}
% The middle column must hold "http://localhost:8090/admin/" unbroken at
% 10.95 pt typewriter, which needs 5.8 cm; at 5.4 it overflows by 7.3 pt.
\begin{longtable}{@{}L{2.9cm}L{5.9cm}L{6.0cm}@{}}
\toprule
\rowcolor{primary!15}
\textbf{What} & \textbf{Where} & \textbf{Sign in with} \\
\midrule
\endhead
The storefront & \texttt{http://localhost:8090/} & nobody; you are a customer \\
The admin panel & \texttt{http://localhost:8090/admin/} & \texttt{admin} /
\texttt{admin123} \\
The database console & \texttt{http://localhost:8091/} & server \texttt{db},
user \texttt{shop}, password \texttt{shoppw}, database \texttt{opencart} \\
\bottomrule
\end{longtable}

\begin{alertbox}[Those credentials are deliberately bad]
\texttt{admin} / \texttt{admin123} is exactly the kind of credential
\chref{security} goes looking for and finds. It is the default so that you can
measure the shop before and after hardening it; Chapter~16 asks you to change
it, and until then every student in the class can sign in to every other
student's shop if they can reach it.

\smallskip
Which they cannot, because nothing here is published to the internet. If you
ever expose this stack beyond your own machine, change the credentials in
\texttt{.env} \emph{first}.
\end{alertbox}

\section{The Four Commands Worth Memorising}

\rowcolors{2}{white}{lightbg}
\begin{longtable}{@{}L{5.0cm}L{10.0cm}@{}}
\toprule
\rowcolor{primary!15}
\textbf{Command} & \textbf{What it does} \\
\midrule
\endhead
\texttt{docker compose up -d} & Start everything. Safe to run when it is
already running \\
\texttt{docker compose down} & Stop everything. \textbf{Your shop is kept} \\
\texttt{docker compose logs -f shop} & Watch the shop's log. This is where PHP
errors go, and they go nowhere else \\
\texttt{docker compose down -v} & Stop everything \textbf{and delete your
shop}. The reset button, and the one mistake with no undo \\
\bottomrule
\end{longtable}

\section{When It Will Not Start}

Five failures account for nearly all of them, and four are not really about
Docker.

\begin{enumerate}[leftmargin=2.2em, itemsep=5pt, topsep=4pt]
  \item \textbf{``Cannot connect to the Docker daemon''} --- the engine is not
        running. On Windows and macOS, start Docker Desktop and wait for its
        icon to stop animating; the command line is ready several seconds
        before the whale stops moving.
  \item \textbf{``Ports are not available''} or ``address already in use'' ---
        something else on your machine already holds port 8090 or 8091. This is
        common; the development machine for this book had both 8080 and 8081
        taken, which is why the defaults are 8090 and 8091 in the first place.
        Edit \texttt{.env}, change \texttt{SHOP\_PORT}, and bring it up again.
        Change only the port in \texttt{.env} --- \texttt{OC\_HTTP\_SERVER}
        reads the same variable, so links keep working.
  \item \textbf{The shop answers but every page is an error} --- the install
        did not finish. Read \texttt{docker compose logs shop} and look for
        ``SUCCESS! OpenCart successfully installed''. If it is absent, the
        database was not ready; \texttt{docker compose down -v} and start
        again.
  \item \textbf{WSL 2 is not installed, or virtualisation is disabled} ---
        Windows only. Docker Desktop's installer handles the first; the second
        is a BIOS setting (often called ``Intel VT-x'', ``AMD-V'' or
        ``SVM Mode'') that some laptops ship disabled.
  \item \textbf{It worked yesterday and today it does not} --- almost always
        disk space. \texttt{docker system df} will tell you;
        \texttt{docker system prune} reclaims what is not in use. It does not
        touch named volumes, so your shop survives --- but read what it
        proposes before agreeing.
\end{enumerate}

\begin{conceptbox}[Read the log before changing anything]
The single most useful habit in this appendix. A shop that is broken is
broken for a reason, and the reason is almost always written down in
\texttt{docker compose logs shop}.

\smallskip
Restarting, rebuilding and resetting are the three things students try first,
and all three destroy the evidence. Look at the log, then act.
\end{conceptbox}

\section{Backing Up, and Why You Will Wish You Had}

Your shop lives in two named volumes. \texttt{down -v} deletes them, and so
does an afternoon of clearing disk space without reading the prompt.

\begin{lstlisting}[language=bash]
# Dump the database into the repository, where git can see it.
docker compose exec -T db mariadb-dump -ushop -pshoppw opencart \
  > backups/shop-$(date +%Y%m%d).sql

# Put it back.
docker compose exec -T db mariadb -ushop -pshoppw opencart \
  < backups/shop-20261001.sql
\end{lstlisting}

Take one after every chapter's laboratory. It costs two seconds, and
\chref{deploy} measures what the alternative costs.

\begin{alertbox}[Where to keep your own work]
Anything you write that must survive belongs in \texttt{code/}, in the
repository --- not inside a running container, where it is deleted at the next
rebuild, and not only in the admin panel, where it is deleted by
\texttt{down -v}.

\smallskip
This matters most in \chref{modules}, where you write a module. Write it in
\texttt{code/modules/} and mount it; do not edit it in place inside the
container. Students who learn this in Week~14 learn it the expensive way.
\end{alertbox}

\section{The Machine Every Measurement Was Taken On}

Every figure in a \faIcon{tachometer-alt}~\textbf{Measured} box in this book was
produced on this machine. Yours will differ; the shapes and ratios should not.

\rowcolors{2}{white}{lightbg}
\begin{longtable}{@{}L{4.6cm}L{10.4cm}@{}}
\toprule
\rowcolor{primary!15}
\textbf{Item} & \textbf{Value} \\
\midrule
\endhead
Host & Windows 11 Pro, 4 CPU cores, 7.7 GiB available to the container
runtime \\
Container runtime & Docker 28.1.1, Linux containers \\
Web and application & Apache 2.4.68, PHP 8.2.34 with OPcache enabled \\
Database & MariaDB 11.8.9 \\
Application & OpenCart 4.1.0.4, installed from the pinned release in
\texttt{code/shop/Dockerfile} \\
Everything on & one machine, so no figure here includes real network latency \\
\bottomrule
\end{longtable}

\begin{alertbox}[The scope of Chapter 16's security work]
\chref{security} asks you to attack a shop: SQL injection, cross-site
scripting, missing security headers, default credentials, an exposed
administrative path. Every one of those exercises is scoped, in the chapter and
here, to \textbf{the container running on your own machine, reached at
\texttt{localhost}, which you installed and which contains no real person's
data}.

\smallskip
That is the whole of it. Running the same checks against any system you do not
own and have not been authorised in writing to test is a criminal offence in
Pakistan under the Prevention of Electronic Crimes Act 2016, and a sackable one
nearly everywhere else. The skill is valuable precisely because it is
dangerous, and the boundary is not a formality.

\smallskip
\chref{legal} returns to this with the statute in front of it.
\end{alertbox}

%=====================================================================
\chapter{Enough PHP, SQL and Twig to Read the Shop}\label{app:languages}
%=====================================================================
\normalfigures

The sibling courses give this appendix to mathematics. This one spends it on
language, because that is what this course actually assumes and does not teach.

You are not asked to be a strong PHP programmer. Until \chref{modules} you will
read far more code than you write, and this appendix is calibrated to that:
enough to follow what the platform is doing, and enough to change it when
\chref{modules} asks you to.

\section{PHP, in Two Pages}

\subsection{The shape of a file}

PHP is a language for producing a response to an HTTP request. A file runs top
to bottom, every time the page is requested, and whatever it prints becomes the
page.

\begin{lstlisting}[language=php]
<?php
declare(strict_types=1);

$name  = 'Keyboard';       // string, single quotes: no interpolation
$price = 2499.00;          // float
$qty   = 3;                // int
$total = $price * $qty;

// Double quotes interpolate; single quotes do not. This is the single
// most common source of confusion in a first week of PHP.
echo "Total for $name: $total";   // Total for Keyboard: 7497
echo 'Total for $name';           // Total for $name
\end{lstlisting}

\begin{alertbox}[Three things that surprise everyone]
\textbf{Variables start with \texttt{\$}, always} --- including inside strings
and in function parameters.

\smallskip
\textbf{\texttt{==} is not \texttt{===}.} The two-character form converts types
before comparing, so \texttt{'0' == false} is true and \texttt{'abc' == 0} was
true in older versions. \textbf{Use \texttt{===} unless you have a reason not
to}, and in \chref{security} you will see what the other one costs.

\smallskip
\textbf{\texttt{declare(strict\_types=1)} is not the default.} Without it PHP
quietly converts \texttt{"5 apples"} into \texttt{5} when a function wants an
integer. OpenCart's own code does not use it; your modules should.
\end{alertbox}

\subsection{Arrays, which are really two things}

\begin{lstlisting}[language=php]
<?php
// A list.
$sizes = ['S', 'M', 'L'];
echo $sizes[0];                      // S
echo count($sizes);                  // 3

// A map -- the same type, with keys. Almost every piece of data in
// OpenCart arrives as one of these.
$product = [
    'product_id' => 42,
    'name'       => 'Keyboard',
    'price'      => 2499.00,
];
echo $product['name'];               // Keyboard

// Walking one. Note the "as $key => $value" form.
foreach ($product as $field => $value) {
    echo "$field = $value\n";
}

// The ?? operator: use this value, or that one if it is absent.
// Without it, a missing key is a warning and a null.
$image = $product['image'] ?? 'placeholder.png';
\end{lstlisting}

\subsection{Functions, classes and the arrow}

\begin{lstlisting}[language=php]
<?php
function formatPrice(float $amount, string $symbol = 'Rs. '): string
{
    return $symbol . number_format($amount, 2);
}

class ShippingQuote
{
    private float $rate;

    public function __construct(float $rate)
    {
        $this->rate = $rate;         // $this-> reaches a property
    }

    public function forWeight(float $kg): float
    {
        return $this->rate * max(1.0, $kg);
    }
}

$q = new ShippingQuote(120.0);
echo formatPrice($q->forWeight(2.5));    // Rs. 300.00
\end{lstlisting}

\begin{conceptbox}[The three arrows, which are different]
PHP uses three symbols that beginners merge into one.

\begin{tight}
  \item \texttt{\$obj->method()} --- call a method on an object, or read a
        property. An \emph{object}.
  \item \texttt{\$array['key']} --- read a key from an array. An \emph{array}.
        Square brackets, quoted key.
  \item \texttt{Class::CONSTANT} --- reach something on the class itself rather
        than on an instance.
\end{tight}

\smallskip
OpenCart hands you arrays far more often than objects, so \texttt{['key']} is
the one you will type most. \texttt{\$this->request->get['product\_id']} is all
three ideas in one expression: a property of a property, then an array key.
\end{conceptbox}

\subsection{Where request data arrives, and why it is dangerous}

\begin{lstlisting}[language=php]
<?php
// Submitted by a form with method="get"  -> $_GET
// Submitted by a form with method="post" -> $_POST
$qty = $_POST['quantity'] ?? 1;

// EVERY value in those arrays is a STRING, and every one of them was
// typed by somebody who may not be a customer. Chapter 16 is about
// what happens when this is forgotten; Chapter 15 is about sessions,
// which is how the shop remembers who you are between requests.
$qty = (int)$qty;                    // now it is certainly an integer
if ($qty < 1 || $qty > 999) {
    $qty = 1;                        // and certainly in range
}
\end{lstlisting}

\section{SQL, as This Course Uses It}

You met SQL in \emph{Database Systems}. This section is a reminder of the five
statements this course needs and the two ideas it leans on.

\begin{lstlisting}[language=sql]
-- Read. The only one you will type often.
SELECT p.product_id, pd.name, p.price
  FROM oc_product p
  JOIN oc_product_description pd ON pd.product_id = p.product_id
 WHERE pd.language_id = 1
   AND p.status = 1
 ORDER BY p.price DESC
 LIMIT 20;

-- Count. Chapter 6 shows why a category page runs one of these as
-- well as the query above, every single time.
SELECT COUNT(*) FROM oc_product WHERE status = 1;

-- Change one row. The WHERE is not optional; without it you have
-- just repriced the entire catalogue.
UPDATE oc_product SET price = 2999.00 WHERE product_id = 42;

-- Remove. Same warning, more so.
DELETE FROM oc_product_to_category WHERE product_id = 42;

-- Make a query fast, or fail to -- see Chapter 6.
CREATE INDEX idx_sort_order ON oc_product (sort_order);
\end{lstlisting}

\begin{definitionbox}[The two ideas everything else rests on]
\term{A join} combines rows from two tables on a matching column. OpenCart
splits a product across eighteen tables (\chref{catalogue}), so almost every
useful question needs at least one join.

\smallskip
\term{An index} is a sorted structure the database can search instead of
reading every row. It makes reads faster and writes slower, and it is useless
if the query hides the column inside a function --- \texttt{ORDER BY
LCASE(name)} cannot use an index on \texttt{name}, which is the central finding
of \chref{catalogue}.
\end{definitionbox}

\begin{examplebox}[EXPLAIN, which answers "why is this slow?"]
Put \texttt{EXPLAIN} in front of any \texttt{SELECT} and the database tells you
what it intends to do rather than doing it. Three words in the output carry
nearly all the information:

\begin{tight}
  \item \textbf{\texttt{type: ALL}} --- a full table scan. Every row read. At
        fifty thousand products this is the one to fix.
  \item \textbf{\texttt{Using filesort}} --- the result had to be sorted after
        it was fetched, because no index gave it in the right order.
  \item \textbf{\texttt{Using temporary}} --- a temporary table was built to
        hold intermediate results, usually because of \texttt{DISTINCT} or
        \texttt{GROUP BY}.
\end{tight}

\smallskip
\chref{catalogue} runs \texttt{EXPLAIN} against the real category query and
finds all three at once.
\end{examplebox}

\section{Twig, Which Is Just the Page}

OpenCart's templates are Twig. A template is HTML with three kinds of hole in
it, and that is very nearly the whole language.

\begin{lstlisting}[language=twig]
{# A comment. It does not reach the browser. #}

{# 1. Print something. #}
<h1>{{ heading_title }}</h1>
<span class="price">{{ product.price }}</span>

{# 2. Do something -- a loop or a condition. #}
{% if products %}
  <ul>
    {% for product in products %}
      <li>
        <a href="{{ product.href }}">{{ product.name }}</a>
        {% if product.special %}
          <s>{{ product.price }}</s> {{ product.special }}
        {% else %}
          {{ product.price }}
        {% endif %}
      </li>
    {% endfor %}
  </ul>
{% else %}
  <p>{{ text_empty }}</p>
{% endif %}

{# 3. Filter -- transform a value on the way out. #}
<p>{{ product.description|striptags|slice(0, 120) }}</p>
<p>{{ product.name|escape }}</p>
\end{lstlisting}

\begin{alertbox}[\texttt{|escape} and what it prevents]
Twig escapes output by default in most configurations, and that default is the
single thing standing between a product description written by somebody else
and a cross-site scripting hole in your shop.

\smallskip
Twig also offers \texttt{|raw}, which turns the escaping off. There are
legitimate uses, and every one of them is a decision to trust the source of
that value completely. \chref{security} looks for places where somebody has
used it without meaning to.
\end{alertbox}

\section{YAML, in Ten Lines}

The Compose file is YAML, and YAML has exactly two rules worth knowing.

\begin{lstlisting}[language=yaml]
# 1. Indentation is structure. Two spaces per level, NEVER tabs --
#    a tab is a syntax error and the message will not say so clearly.
services:
  db:
    image: mariadb:11
    environment:
      MARIADB_DATABASE: opencart

# 2. A leading dash is a list item.
    ports:
      - "8090:80"
      - "3306:3306"
\end{lstlisting}

That is enough to read \texttt{code/docker-compose.yml} and to change the one
thing \chref{stack} asks you to change, which is a port number.

\section{Where To Look Next}

\rowcolors{2}{white}{lightbg}
\begin{longtable}{@{}L{3.2cm}L{11.8cm}@{}}
\toprule
\rowcolor{primary!15}
\textbf{For} & \textbf{Go to} \\
\midrule
\endhead
PHP & \texttt{php.net/manual} --- the official manual, and unusually good. Each
function page has user-contributed notes; treat those as hints, not authority \\
SQL & The MariaDB knowledge base, and your \emph{Database Systems} notes for
the theory this appendix skips \\
Twig & \texttt{twig.symfony.com/doc} --- the ``for template designers'' page is
all you need \\
OpenCart & \texttt{docs.opencart.com}, and the source in your own container,
which is more current than the documentation \\
\bottomrule
\end{longtable}

\begin{conceptbox}[Read the source; it is right there]
The shop's entire source is inside your container, and you can read it:

\smallskip
\texttt{docker compose exec shop \textbackslash}\\
\texttt{\phantom{xx}cat /var/www/html/catalog/model/catalog/product.php}

\smallskip
This is not a last resort. \chref{catalogue} finds the category page's real
query this way, and it is the only way to answer ``what does this page actually
do?'' with certainty rather than with a plausible guess. The documentation
describes a version; the container holds yours.
\end{conceptbox}

%=====================================================================
\chapter{Answers to the Checkpoints}\label{app:answers}
%=====================================================================
\normalfigures

Three answers per chapter, in order. Attempt the checkpoint from memory before
reading these, and argue with them where you disagree --- several are judgement
calls rather than facts, and they say so.

% The slide pipeline parses this appendix: \ansch{<number> --- <title>} must be
% followed immediately by an ans environment of exactly three items.
\newcommand{\ansch}[1]{\subsection*{\color{primary}Chapter #1}}
\newenvironment{ans}
  {\begin{enumerate}[leftmargin=*, itemsep=4pt, topsep=3pt,
                     label=\textbf{\arabic*.}]}
  {\end{enumerate}}

%---------------------------------------------------------------------
\ansch{1 --- What E-Commerce Is, and What It Changed}
\begin{ans}
  \item \textbf{Electronic business, not electronic commerce.} Viewing results
        and downloading transcripts are business processes conducted over a
        network, but no value changes hands through the system --- the
        transaction happens at a bank counter. The moment fee payment moves
        online it becomes e-commerce, and the university acquires every problem
        to the right of the dashed line in \dsref{fourfunctions}: a payment
        that may have succeeded while the network dropped, a receipt that must
        exist exactly once, a refund path.
  \item Many are acceptable. One good answer: \textbf{refusing to sell the last
        item when it is visibly damaged.} In software this requires a stock
        record that distinguishes ``held'' from ``sellable'', somebody whose
        job is to update it, a rule for what the storefront does when the two
        disagree, and a decision about what the customer sees --- out of stock,
        or a quiet substitution? The shopkeeper does all of that in half a
        second without recording that a decision was made, which is exactly why
        it is forgotten when the shop goes online.
  \item \textbf{Two uncounted costs:} the acquisition cost of finding customers
        who currently arrive free (\chref{models} puts this at Rs.~1{,}200 an
        order for a typical small shop, which can exceed the commission), and
        the labour of running the four functions themselves --- payment
        integration, delivery arrangements, customer service, returns. The
        commission was buying those, not just the listing. \textbf{What they
        gain that is not money:} the customer relationship. They learn who
        bought, can contact them again, and cannot be delisted. For a business
        with repeat custom that is frequently worth more than the 15\%.
\end{ans}

%---------------------------------------------------------------------
\ansch{2 --- Business Models and Revenue Models}
\begin{ans}
  \item \textbf{Business model:} a marketplace connecting restaurants with
        diners, with fulfilment performed in-house --- three and a half of the
        four functions. \textbf{Revenue models: three at once} ---
        \emph{commission} from restaurants, a \emph{service or delivery fee}
        from customers, and \emph{advertising} in the form of promoted
        placement. \textbf{Most profitable per unit: the advertising.} It has
        almost no marginal cost, where commission is earned against riders and
        support that must be paid for, and the delivery fee rarely covers the
        rider. Promoted placement is close to pure margin, which is why these
        platforms push it so hard.
  \item Viable when $\text{LTV} > \text{CAC}$, that is when
        $400 \times n > 1{,}500$, so when the average customer orders
        $n > 3.75$ --- \textbf{at least four times} --- before they stop.
        \emph{What to measure:} the repeat rate and the time over which it
        happens. Four orders over six months is a business; four orders over
        six years is not, because the cash was spent today and returns over a
        period during which the venture has to survive.
  \item B2C platforms store \textbf{one price per product} --- a column on the
        product row. B2B pricing is a \emph{relation} between customer and
        product, often with quantity breaks and expiry dates. That is not a
        setting you have failed to find; it is a different cardinality, and
        making it work means either a different platform or an extension that
        intercepts every price lookup in the system (\chref{modules} shows how
        much work that is). The assumption that breaks is ``a product has
        \emph{a} price''.
\end{ans}

%---------------------------------------------------------------------
\ansch{3 --- Auctions, Portals and Communities}
\begin{ans}
  \item \textbf{Bid Rs.~15{,}000} --- your true valuation. You win, and you pay
        the second-highest bid, \textbf{Rs.~13{,}500}, gaining Rs.~1{,}500.
        Bidding Rs.~18{,}000 does not improve this outcome at all: you still
        pay Rs.~13{,}500, because the price is set by the others. What it does
        is expose you to winning when somebody bids Rs.~16{,}500, at which
        point you pay Rs.~16{,}500 for a thing worth Rs.~15{,}000 to you. The
        overbid can only ever buy you an auction you did not want.
  \item It \textbf{never closes}. The auction stays open indefinitely; the
        winning bidder is never notified, the seller is never paid, and the
        listing accumulates further bids whenever somebody wanders past. The
        seller sees an item that appears to be permanently ``ending soon'' and
        a buyer who has given up. Worse, the failure is \emph{silent} and
        inversely proportional to interest: popular auctions close correctly
        and unpopular ones do not, so the bug is invisible in testing, where
        you load every page you create.
  \item Any two of: \textbf{be a seller yourself} until real ones arrive
        (cost: you take on stock and inventory risk, and you compete with the
        sellers you are recruiting); \textbf{narrow the market} to one city or
        one category (cost: you cap your growth and may pick the wrong niche);
        \textbf{give sellers a tool worth having without buyers} (cost: you are
        now building and supporting a second product); \textbf{list supply
        gathered elsewhere} (cost: legally risky and reputationally
        expensive if discovered).
\end{ans}

%---------------------------------------------------------------------
\ansch{4 --- Planning an E-Commerce Framework}
\begin{ans}
  \item (a)~\textbf{Functional}, and it \emph{does not} discriminate --- all
        four options reset passwords. (b)~\textbf{Functional}, and it
        \textbf{discriminates sharply}: per-customer pricing eliminates most
        hosted platforms and the marketplace route outright, and is the single
        most common reason to choose self-hosted or bespoke. (c)~\textbf{
        Non-functional}, and it discriminates \emph{in the opposite direction}
        to what students expect --- surviving a twentyfold spike is something a
        large hosted platform does effortlessly and your single virtual server
        does not.
  \item \textbf{Hosting} (Rs.~324{,}000 in the worked example),
        \textbf{setup and theming labour} (Rs.~175{,}000 plus
        Rs.~100{,}000 of custom modules), and above all \textbf{maintenance
        labour} (Rs.~360{,}000). Their stated saving of Rs.~792{,}000 is
        therefore closer to Rs.~339{,}000 --- still a saving, but less than
        half what they claimed, and it vanishes entirely if maintenance takes
        eight hours a month instead of four.
  \item Many are acceptable. A clear one: \textbf{a home baker selling six
        kinds of cake to one city, taking orders through Instagram.} Self-hosting
        would mean paying for a server and maintaining it to serve perhaps
        thirty orders a week, with nobody to patch it --- which is not a
        cheaper shop but an unpatched one. \textbf{Choose marketplace-only or a
        cheap hosted platform}: the business's scarce resource is the baker's
        time, the volume is low, nothing about the pricing is unusual, and
        every hour spent on a server is an hour not spent baking.
\end{ans}

%---------------------------------------------------------------------
\ansch{5 --- The Stack, and Your Shop in One Command}
\begin{ans}
  \item \textbf{Working:} the operating system and the web server --- a 500 is
        generated and served by Apache, so the request reached it and a
        response came back. \textbf{Failed:} the application layer; PHP raised
        an error and died before producing a page. The database is
        \emph{unknown} --- it may be the cause (the application could not
        connect) or entirely healthy. \textbf{Look first} at the PHP error log,
        with \texttt{docker compose logs shop}; the message is never on the
        screen, because \texttt{display\_errors} is off, and \chref{security}
        is about why you want it that way.
  \item \texttt{depends\_on} orders container \emph{starts}, and a container
        that has started is not a database that will accept a login --- MariaDB
        opens its TCP port several seconds before it finishes initialising. The
        shop therefore starts, runs the installer, and the installer is refused
        by a database that is listening but not ready. \textbf{The failure is
        intermittent and worse on a fast machine}, because the faster the shop
        container starts, the more certainly it arrives early. The
        \texttt{service\_healthy} condition waits for the database's own
        healthcheck to pass, which tests a login rather than a port.
  \item \textbf{Lost:} both named volumes --- the entire database (products,
        categories, orders, customers, settings) and the shop's installed
        files, including any module you had written in place. \textbf{Not
        lost:} anything in the repository --- the Compose file, the Dockerfile,
        \texttt{code/}, and any module kept there rather than edited inside the
        container. \textbf{What would have made it minor:} a database dump
        taken on a schedule and kept outside the volume, plus the discipline of
        writing modules in \texttt{code/modules/} and mounting them, so that
        the only thing destroyed is state you can regenerate. \chref{deploy}
        makes both of these routine.
\end{ans}

%---------------------------------------------------------------------
\ansch{6 --- Products, Categories and Their Tables}
\begin{ans}
  \item \textbf{Move to \texttt{oc\_product\_description}:} the product name
        and the meta description. Both are \emph{words}, and words differ by
        language. \textbf{Stay in \texttt{oc\_product}:} price, weight and
        stock quantity --- a kilogram is a kilogram in Urdu, and a shop that
        let the price depend on the reader's language would eventually charge
        two different prices for one product. The test is the first of the four
        questions: \emph{does this fact depend on who is reading?}
  \item \textbf{Buys:} a single-query answer to ``every category beneath this
        one'', so a category page costs one query instead of a loop whose
        length depends on how deep the tree happens to be.
        \textbf{Costs:} the same fact is now stored twice, so the two can
        disagree, and every write that changes the shape of the tree must
        rewrite the derived rows. \textbf{The operation that makes it visible:}
        \emph{moving a category}. Changing one \texttt{parent\_id} invalidates
        the path rows of that category and of everything beneath it, so the
        admin panel must rewrite them all --- which is why moving a category
        with many descendants is noticeably slower than renaming one.
  \item Two of several good answers. \textbf{(a)~The database was never the
        bottleneck}: if the database is 20\% of the page, halving it improves
        the page by 10\%, which can be smaller than the run-to-run variance.
        \textbf{(b)~The expensive thing is the \emph{number} of queries, not
        the cost of any one} --- the measurement in this chapter found 146
        queries for one page, and an index helps exactly one of them.
        \smallskip
        \textbf{To tell them apart:} turn on the query log with
        \texttt{long\_query\_time=0}, request the page once, and read
        \texttt{COUNT(*)} and \texttt{SUM(query\_time)} from
        \texttt{mysql.slow\_log}. A large total with few queries is (a); a
        small total spread over a hundred queries is (b). That single
        measurement separates them, and it is the one in
        \texttt{m06\_catalogue.py}.
\end{ans}

%---------------------------------------------------------------------
\ansch{7 --- Product Variations, Options and Uploads}
\begin{ans}
  \item \textbf{The photograph is an option} --- specifically a \texttt{file}
        option. It is supplied by the customer, it is different every time, and
        the shop holds no stock of it. \textbf{The size is a variant}, because
        a small mug and a large mug are two different objects on two different
        shelves, each with its own cost, its own barcode and its own count.
        \textbf{What the shop holds in stock:} blank mugs in two sizes. Nothing
        else --- which is exactly why the photograph cannot be a variant: there
        is no such thing as stock of ``the mug with Ahmed's photograph on it''
        until Ahmed orders one.
  \item \textbf{Option value rows: $4 + 3 + 2 = \mathbf{9}$}, plus three rows
        in \texttt{oc\_product\_option}. \textbf{As variants: $4 \times 3 \times
        2 = 24$ combinations, so \textbf{25} \texttt{oc\_product} rows} ---
        twenty-four plus the master --- and 25 description rows, 25 store rows
        and 25 category rows beside them. Nine against a hundred, for the same
        catalogue.
  \item \textbf{Passes:} the filename sanitiser (nothing illegal in
        ``logo.svg''), the length check, the extension allow-list
        (\texttt{svg} is on the stock list), and the MIME allow-list --- it
        claimed \texttt{image/png}, which is allowed, and nothing verifies that
        claim. \textbf{Fails:} nothing. All five checks pass.
        \smallskip
        \textbf{For it to harm a visitor}, the file must be \emph{served to
        them from your origin and rendered}. Two things would have to be true:
        the attacker must learn the stored URL, which carries a 32-character
        random token and so is not guessable --- but it \emph{is} shown to
        whoever uploaded it and may be forwarded to an administrator who opens
        it --- and the server must serve it with
        \texttt{Content-Type: image/svg+xml} rather than as a download, which
        \chref{stack}'s stock install does. The realistic victim is therefore
        the administrator reviewing the order, and the realistic fix is to drop
        \texttt{svg} from the list and serve uploads with
        \texttt{Content-Disposition: attachment} (\chref{security}).
\end{ans}

%---------------------------------------------------------------------
\ansch{8 --- Themes, Layout and the Customer's Experience}
\begin{ans}
  \item \textbf{Best case for compressing the photographs:} the images are
        $713 - 597 - 54 = $ about \textbf{61\,KB}. Halving them saves 30\,KB,
        which is \textbf{4\%} of the page --- and that is the \emph{best} case,
        because you cannot compress them to nothing. \textbf{Propose instead:}
        minify \texttt{bootstrap.css}, which is 265\,KB of the 597 and has no
        minified copy in the tree. That single change saves roughly 240\,KB,
        about \textbf{34\%} of the page --- eight times the whole image budget
        --- and changes nothing a customer can see. Then remove the unused
        parts of Bootstrap and FontAwesome, and only then look at the images.
  \item \textbf{Layout change} if a ``trust badge'' module already exists:
        \adminpath{Design, Layouts, Product}, add the module to
        \uiel{Content Bottom}, set its sort order. No code, visible to the next
        person, survives upgrades. \textbf{Template override} if it must sit
        \emph{immediately under the button}, inside the markup the theme
        generates: there is no position there, so you override
        \texttt{product/product.twig} in your child theme and add it. The
        general rule: a position that exists is a layout change; markup that
        must differ is a template change.
  \item \textbf{Two reasons it is lower.} (a)~\textbf{It follows only assets
        named in the HTML.} Fonts are requested by the CSS once the browser
        parses it, so none of the FontAwesome \texttt{.woff2} files --- 248\,KB
        in that directory --- are counted. (b)~\textbf{It does not execute
        JavaScript}, so anything a script loads afterwards (a tracking pixel,
        a lazy-loaded image, a carousel's later slides) is invisible to it.
        \smallskip
        \textbf{Which it could fix:} the first, by parsing the CSS for
        \texttt{url()} references and following them --- a dozen lines. The
        second it cannot fix, because fixing it means running a browser, which
        is a different kind of tool. Knowing which limitation is laziness and
        which is fundamental is the useful judgement here.
\end{ans}

%---------------------------------------------------------------------
\ansch{9 --- The Shopping Basket, and Where State Lives}
\begin{ans}
  \item \textbf{What would change:} the two queries per request that the
        \texttt{db} engine spends reading and writing the session record. On
        this host that was worth about 20\% of the page --- 89\,ms to 71\,ms
        --- when those two queries were removed altogether by the
        \texttt{file} engine, and Redis would remove them from the database in
        the same way. \textbf{What would not change: the basket.} It is rows in
        \texttt{oc\_cart}, and the \texttt{Cart rows} column was 1 under both
        engines. \textbf{The number that settles it:} the session record is
        \textbf{18 bytes} and holds \texttt{\{"currency":"USD"\}}. Whatever you
        move it to, you are moving eighteen bytes. The right place to look for
        basket performance is the six cart queries an empty basket costs, not
        the store the session sits in.
  \item \textbf{Where they came from:} a session row is written on every
        request, before the visitor has asked for anything, so every page view
        by every crawler, uptime check and measurement harness created one. The
        645 rows here are the by-product of writing this handout.
        \textbf{Why they are not removed on departure:} nothing tells a server
        that a visitor has left --- HTTP has no ``goodbye''. Rows can therefore
        only expire on a clock, and \texttt{config\_session\_expire} is 86{,}400
        seconds. \textbf{What to measure first:} the proportion of sessions
        that ever acquire a basket or a login. If almost none do, the expiry is
        not the problem and robots are; shortening the expiry would then also
        shorten the life of the few baskets that matter, which is the opposite
        of what was wanted. Measure before tuning the number that is easiest
        to reach.
  \item \textbf{What each removes:} \texttt{down} removes the containers ---
        the running processes and their writable layers. The named volumes
        \texttt{dbdata} and \texttt{shopdata} are separate objects and are left
        alone, so the database files survive and the basket with them;
        \texttt{up -d} attaches the new containers to the same volumes and the
        shop resumes mid-visit. \texttt{down -v} additionally deletes those
        volumes, which is where every byte of durable state lives, so the
        entrypoint finds an empty database and reinstalls from scratch.
        \textbf{One other thing it destroys:} every product, category, order,
        customer and setting you have made since \chref{stack} --- in this
        measurement the shop came back in 17 seconds looking perfectly
        healthy, with all of it gone. Uploaded images go with the
        \texttt{shopdata} volume at the same time.
\end{ans}

\input{parts/D_glossary}
\input{parts/E_resources}
\input{parts/F_boundary}
%=====================================================================
\chapter{The Sixteen-Week Schedule}\label{app:schedule}
%=====================================================================
\normalfigures

This appendix is the \textbf{source of record} for the teaching schedule. The
one-page \texttt{ECAD\_Fall2026\_Lecture\_Schedule.pdf} distributed to students
is generated from it, and the slide pipeline parses the table below. Change
this first; regenerate the one-page version afterwards.

Twenty-three chapters are taught across sixteen weeks, so seven weeks carry two
chapters. The pairings are not arbitrary. In every case the two chapters are
either one topic the printed book splits for length --- the catalogue and its
variants, the two halves of the back office --- or a pair in which the second
is unusable without the first: a basket with nothing to check out, a gateway
with no payment method chosen.

One pairing is deliberate in the other direction. \textbf{Chapter~16, Security,
has a week to itself}, although it would pair naturally with Chapter~15. Its
laboratory attacks the student's own shop and then hardens it, and that does
not fit in half a week alongside anything else.

\section{The Schedule}

% MACHINE-READ. The slide builder takes the FIRST longtable in this file and
% requires: five columns; a bare integer week; chapters as comma-separated
% digits or N--M ranges; literal --- for an empty assessment cell; and no
% chapter appearing in two weeks.
\rowcolors{2}{white}{lightbg}
% column 3 holds the bold head "Chapters", which needs 1.8 cm at 10.95 pt
\begin{longtable}{@{}L{1.1cm}L{3.2cm}L{1.8cm}L{4.4cm}L{3.9cm}@{}}
\toprule
\rowcolor{primary!15}
\textbf{Wk} & \textbf{Lecture topic} & \textbf{Chapters} & \textbf{Lab} &
\textbf{Assessment} \\
\midrule
\endhead
1 & What e-commerce changed; business and revenue models & 1, 2 & Read four
real shops; the arithmetic of one order & --- \\
2 & Auctions, portals and communities; planning a framework & 3, 4 & Three
shapes compared; your requirements and your decision & --- \\
3 & The stack, and your shop in one command & 5 & Compose up; verify, break and
reset the shop & --- \\
4 & Products, categories and the data model beneath them & 6 & 50{,}000
products, and the index that rescues the category page & --- \\
5 & Product variations, options and user uploads & 7 & Option combinations, and
what an upload field accepts & \textbf{Quiz 1 (Ch 1--5)} \\
6 & Themes, layout and the user experience & 8 & Page weight before and after,
measured in bytes and requests & \textbf{Assignment 1 set} \\
7 & The shopping basket, and where state lives & 9 & Session, cookie and
database baskets across a restart & --- \\
8 & The checkout and the order state machine & 10 & Midterm, then queries
counted per checkout step & \textbf{Midterm (Ch 1--9)} \\
9 & Shipping, tax and zones; discounts and vouchers & 11, 12 & Geo zones and
GST; a voucher applied in the wrong order & \textbf{Assignment 1 due} \\
10 & Web payment systems; taking payment for orders & 13, 14 & A sandbox
gateway, and what happens when it times out & \textbf{Project proposal} \\
11 & User account management & 15 & Password cost against logins per second &
--- \\
12 & Securing the shop & 16 & The OWASP checks against your own container,
before and after hardening & \textbf{Quiz 2 (Ch 10--15)} \\
13 & Administration: the dashboard, orders, customers, refunds & 17, 18 & A
refund through the order state machine & \textbf{Assignment 2 set} \\
14 & Extending the platform: your first module & 19 & A shipping method and a
dashboard widget, written from nothing & --- \\
15 & Deployment, maintenance and performance; SEO & 20, 21 & Load test with
cache on and off; crawlability measured & \textbf{Assignment 2 due} \\
16 & Social, legal and ethical issues; project defence & 22, 23 & Project
presentations & \textbf{Project due; presentations} \\
\bottomrule
\end{longtable}

\section{How the Assessment Adds Up}

HEC publishes the assessment instruments for this course --- sessional
examination, home assignments, quizzes, project, presentations, final
examination --- without weights. The weights below are the department's, and
are stated here so that a student can see the whole shape of the semester on
one page.

\rowcolors{2}{white}{lightbg}
\begin{longtable}{@{}L{4.6cm}L{2.0cm}L{8.4cm}@{}}
\toprule
\rowcolor{primary!15}
\textbf{Instrument} & \textbf{Weight} & \textbf{What it examines} \\
\midrule
\endhead
Two quizzes & 10\% & Weeks 5 and 12. Short, closed-book, on the chapters named
in the schedule \\
Two assignments & 15\% & Set in weeks 6 and 13. Both are done on your own shop
and submitted as a change to it plus a short write-up \\
Midterm examination & 20\% & Week 8, Chapters 1--9 \\
Semester project & 25\% & The shop itself, built across the whole semester
(\chref{project}), submitted in week 16 \\
Presentation and defence & 5\% & Week 16. Ten minutes, plus questions you are
expected to be able to answer about your own decisions \\
Final examination & 25\% & All chapters, with emphasis after the midterm \\
\bottomrule
\end{longtable}

\begin{alertbox}[The project is not a week-16 activity]
Twenty-five per cent of the marks are for a shop that is built in the
laboratory of every chapter from \chref{stack} onward. A student who treats the
laboratories as optional homework and intends to build the project in the last
fortnight will be building, in two weeks and without help, something the rest
of the class built in fourteen weeks with a lecturer in the room.

\smallskip
\chref{project} states what is required and how it is marked. Read it in
Week~3, not in Week~15.
\end{alertbox}


\end{document}
